% type the following command to generate pdf file:
%   tectonic "Report for JPMorgan Chase MLCOE internship - Q1 final report - extension.tex"

\documentclass[12pt,a4paper]{article}

% --- Packages ---
\usepackage[margin=2.5cm]{geometry}
\usepackage{hyperref}
\usepackage{listings}
\usepackage{xcolor}
\usepackage{caption}        % provides \ContinuedFloat for multi-page figures
\usepackage[breakable,skins]{tcolorbox}
\newtcolorbox{llmoutput}[1][]{
  colback=gray!5, colframe=gray!50,
  fonttitle=\bfseries, title={#1},
  breakable, enhanced
}
\usepackage{booktabs}
\usepackage{array}
\usepackage{enumitem}
\usepackage{titlesec}
\usepackage{fancyhdr}
\usepackage{graphicx}
\usepackage{float}
\usepackage{amsmath,amssymb}
\usepackage{lmodern}
\usepackage{parskip}
\usepackage{unicode-math}
\defaultfontfeatures{Scale=MatchLowercase}
\defaultfontfeatures[\rmfamily]{Ligatures=TeX,Scale=1}
\usepackage{xeCJK}
\setCJKmainfont{Hiragino Sans GB}
\usepackage{soul}

% --- Colour Definitions ---
\definecolor{codeblue}{RGB}{30,100,200}
\definecolor{codegray}{RGB}{100,100,100}
\definecolor{codegreen}{RGB}{0,140,0}
\definecolor{codered}{RGB}{180,0,0}
\definecolor{backcolour}{RGB}{245,245,245}
\definecolor{jpblue}{RGB}{0,45,114}

% --- Code Listing Styles ---
\lstdefinestyle{pythonstyle}{
    language=Python,
    backgroundcolor=\color{backcolour},
    commentstyle=\color{codegreen},
    keywordstyle=\color{codeblue}\bfseries,
    numberstyle=\tiny\color{codegray},
    stringstyle=\color{codered},
    basicstyle=\ttfamily\footnotesize,
    breakatwhitespace=false,
    breaklines=true,
    captionpos=b,
    keepspaces=true,
    numbers=left,
    numbersep=5pt,
    showspaces=false,
    showstringspaces=false,
    showtabs=false,
    tabsize=4,
    frame=single,
    rulecolor=\color{codegray},
    literate={\_}{\_}1,
}
\lstdefinestyle{json}{
    backgroundcolor=\color{backcolour},
    basicstyle=\ttfamily\footnotesize,
    stringstyle=\color{codegreen},
    breaklines=true,
    breakatwhitespace=false,
    frame=single,
    rulecolor=\color{codegray},
    showstringspaces=false,
    tabsize=2,
    literate={\_}{\_}1,
    captionpos=b,
}
\lstdefinestyle{console}{
    backgroundcolor=\color{backcolour},
    basicstyle=\ttfamily\footnotesize,
    breaklines=true,
    breakatwhitespace=false,
    frame=single,
    rulecolor=\color{codegray},
    showstringspaces=false,
    captionpos=b,
}
\lstdefinestyle{python}{style=pythonstyle}
\lstset{style=pythonstyle}

% --- Image scaling (kept from pandoc) ---
\makeatletter
\newsavebox\pandoc@box
\newcommand*\pandocbounded[1]{%
  \sbox\pandoc@box{#1}%
  \Gscale@div\@tempa{\textheight}{\dimexpr\ht\pandoc@box+\dp\pandoc@box\relax}%
  \Gscale@div\@tempb{\linewidth}{\wd\pandoc@box}%
  \ifdim\@tempb\p@<\@tempa\p@\let\@tempa\@tempb\fi
  \ifdim\@tempa\p@<\p@\scalebox{\@tempa}{\usebox\pandoc@box}%
  \else\usebox{\pandoc@box}%
  \fi%
}
\def\fps@figure{htbp}
\makeatother

% --- Pandoc compatibility ---
\setlength{\emergencystretch}{3em}
\providecommand{\tightlist}{%
  \setlength{\itemsep}{0pt}\setlength{\parskip}{0pt}}

% --- Header / Footer ---
\pagestyle{fancy}
\fancyhf{}
\rhead{\color{jpblue}\textbf{JPMorganChase --- Question 1 Final Report}}
% \lhead{\color{jpblue}Confidential}
\rfoot{Page \thepage}
\lfoot{\color{codegray}\today}

% --- Section Formatting ---
\titleformat{\section}{\large\bfseries\color{jpblue}}{\thesection}{0.75em}{}[\titlerule]
\titleformat{\subsection}{\normalsize\bfseries\color{jpblue}}{\thesubsection}{0.75em}{}
\titleformat{\subsubsection}{\normalsize\bfseries\color{jpblue}}{\thesubsubsection}{0.75em}{}

% --- Hyperref ---
\hypersetup{
  colorlinks=true,
  linkcolor=jpblue,
  urlcolor=jpblue,
  citecolor=jpblue,
  pdftitle={Final Report for JPMorgan Chase MLCOE Internship -- Q1},
  pdfauthor={Jaebum (Albert) Chung},
}

\begin{document}

% --- Title Page ---
\begin{titlepage}
  \centering
  \vspace*{2cm}
  {\Huge\bfseries\color{jpblue} Financial Statement Forecasting\par}
  \vspace{0.3cm}
  {\Huge\bfseries\color{jpblue} from Annual Reports\par}
  \vspace{0.5cm}
  {\large\color{codegray} A Structural No-Plug Approach with Machine Learning\par}
  \vspace{2cm}
  \rule{0.6\textwidth}{0.4pt}\\[0.5cm]
  {\large JPMorganChase --- MLCOE TSRL Internship Project (Question 1)\par}
  \vspace{0.5cm}
  {\normalsize Jaebum (Albert) Chung\par}
  \vspace{0.3cm}
  {\normalsize April 6, 2026\par}
  \vspace{3cm}
  \begin{abstract}
    \noindent This report implements a structural financial statement forecasting model
    based on the no-plug framework of \hyperlink{ref:pareja09}{Pareja(09)}, where the balance sheet is derived
    endogenously from the cash budget and income statement without circular references.
    Policy variables are trained from historical data using gradient-based optimisation,
    and Bayesian extensions provide confidence intervals via Monte Carlo simulation.
    The framework is extended with LLM-based extraction pipelines for automated
    ingestion of annual reports and strategic financial analysis.
  \end{abstract}
\end{titlepage}

\tableofcontents
\newpage

\section{Introduction}

\subsection{Literature Review}

\emph{The Limitations of ``Plug'' Variables in Financial Forecasting}

Financial statement forecasting is a foundational practice in corporate
finance, yet traditional methodologies often compromise the integrity of
the double-entry accounting system. As reviewed in \hyperlink{ref:pareja07}{Pareja et al.\ (2007)}
and \hyperlink{ref:pareja09}{Pareja (2009)}, a significant portion of existing literature relies
on the use of ``plugs'': balancing variables forced to reconcile the
difference between Assets and Liabilities plus Equity. Typically,
analysts assign a specific account (often Cash, Debt or Equity
\hyperlink{ref:pareja09}{[Pareja(09)]}) to absorb any discrepancies in the forecast.

While this approach ensures a balanced sheet mathematically, it is
methodologically flawed for robust forecasting. By using a plug, the
analyst effectively overrides the natural evolution of the balance
sheet, obscuring the causal relationships between operating decisions
and financial health. More critically, the use of a plug erodes the
primary validation mechanism of the double-entry principle. If the
balance sheet is forced to balance, it becomes impossible to detect
logical errors or inconsistencies in the underlying code. As noted in
the critiques within the Pareja framework, a model that balances by
force rather than by logic fails to serve as a reliable tool for stress
testing or error detection.

Financial statements are governed by strictly defined accounting
identities and deterministic valuation rules, such as the First-In,
First-Out (FIFO) inventory method \hyperlink{ref:iasb21}{[IASB(21)]}. Consequently, the
direct application of machine learning algorithms, such as Long
Short-Term Memory (LSTM) networks, to forecast individual line items is
often methodologically flawed. Because these `black
box' models treat variables as probabilistic time
series rather than components of an integrated system, they fail to
enforce the structural interconnectedness required by double-entry
bookkeeping. This limitation frequently results in forecasts that
violate the fundamental accounting equation (Assets = Liabilities +
Equity), thereby rendering the model incapable of accurately assessing a
company's overall financial health.

\subsection{The Double-Entry Principle}
\label{sec:double-entry}

Double-entry bookkeeping is the foundation of modern accounting. Its
central rule is that every financial transaction must be recorded as two
equal and opposite entries across ledger accounts, so that the fundamental
accounting equation is preserved at all times:

\[
\text{Assets} = \text{Liabilities} + \text{Equity}
\]

To see why, consider two scenarios and how they flow through the
three core financial statements: the Income Statement, the Statement
of Cash Flows, and the Balance Sheet.

\paragraph{Example 1: A credit sale and subsequent collection
(Income Statement $\to$ Balance Sheet $\to$ Cash Flow Statement).}
A company sells \$1{,}000 of goods on credit. At the point of sale,
the Income Statement records \$1{,}000 in Revenue, which, after
taxes and expenses, flows into Net Income and ultimately into
Retained Earnings (a component of Equity). On the Balance Sheet,
Accounts Receivable (an asset) increases by \$1{,}000. The equation
remains balanced: assets rose by \$1{,}000 (AR) and equity rose by
the same amount (via Retained Earnings absorbing Net Income).

When the customer later pays the \$1{,}000 in cash, the Cash Flow
Statement records a \$1{,}000 operating cash inflow. On the Balance
Sheet, Cash increases by \$1{,}000 while Accounts Receivable
decreases by \$1{,}000. One asset replaced another, and Total Assets
are unchanged. The equation still holds because this is a pure asset
swap. Notice how a single economic event (a sale) creates an accrual
on the Income Statement first, then resolves through the Cash Flow
Statement later, with the Balance Sheet tracking both stages.

\paragraph{Example 2: Debt-financed equipment purchase
(Cash Flow Statement $\to$ Balance Sheet $\to$ Income Statement).}
A company borrows \$2{,}000 in long-term debt and immediately uses
the proceeds to purchase equipment. The Cash Flow Statement records a
\$2{,}000 inflow under financing activities (the borrowing) and a
\$2{,}000 outflow under investing activities (the capital
expenditure), so the net cash impact is zero. On the Balance Sheet,
Non-Current Assets (equipment) increase by \$2{,}000 and Long-Term
Debt (a liability) increases by \$2{,}000. Both sides of the equation
grew by the same amount, so it remains balanced.

In subsequent periods, the transaction continues to ripple through the
statements. The Income Statement records Depreciation expense (reducing
Net Income as the asset is consumed) and Interest expense on the debt
(reducing Net Income further). Both expenses flow back to the Balance
Sheet: Depreciation reduces Non-Current Assets, and the lower Net
Income reduces Retained Earnings. Meanwhile, the Cash Flow Statement
captures the interest and principal payments as outflows, each
reducing Cash on the Balance Sheet while simultaneously reducing the
outstanding Debt liability. At every step, the equation holds.

\paragraph{Connection to our model.}
These examples reveal why double-entry bookkeeping makes
\emph{driver-based} financial forecasting without plugs or circularity possible. Because every value flow is
recorded on both sides of the ledger (every expense has a matching
cash outflow, every credit sale has a matching receivable, every
depreciation charge has a matching reduction in Non-Current
Assets), the financial state at time $t+1$ is fully determined by the
state at time $t$ together with the flows between them. Nothing escapes
the ledger.

To make these mechanics tangible, we built a simple interactive visualisation
as a Claude Artifact (Figure~\ref{fig:double-entry}). It lets the reader
select any of six representative transactions and trace their debit and
credit legs simultaneously across the Income Statement, Cash Budget, and
Balance Sheet, together with the formal journal entry and a
verification that $A = L + E$ is preserved.

\begin{figure}[H]
\centering
% TODO: replace with a screenshot of the live artifact
\includegraphics[width=\textwidth]{report_latex_media/media/double_entry_artifact.png}
\caption{Interactive visualisation of double-entry bookkeeping across
the three financial statements. Six transactions --- credit sale, cash
collection, depreciation, net-income close to retained earnings, CapEx
purchase, and debt repayment --- can be selected to trace their debit
(\textcolor{red}{Dr}, red) and credit (\textcolor{blue}{Cr}, blue) legs
simultaneously through the Income Statement, Cash Budget, and Balance
Sheet. Each selection displays the formal journal entry and verifies
that the accounting identity $A = L + E$ is preserved. The live
interactive version is available at
\url{https://claude.ai/public/artifacts/2706f0bf-ba12-4fa9-964a-6ba54ccc2838}.}
\label{fig:double-entry}
\end{figure}

\subsection{Proposed Methodology: A Structural ``No-Plug'' Approach}

To address the limitations of plug-based financial forecasts, this report implements a
structural forecasting model based on the framework proposed by
Pareja~(2009) with the revenue as the primary driver. The core distinction of this method is the rejection of
circular references and balancing plugs. Instead, the model derives
next year's Balance Sheet and Income Statement endogenously from this
year's state through a \emph{Cash Budget} and the projected revenue.

\paragraph{The three financial statements.}
A company's financial position at any point in time is captured by three
interconnected statements:

\begin{itemize}
\item \textbf{Balance Sheet}: A snapshot of what the company owns
  (Assets: non-current assets, receivables, inventory, cash, market
  securities) and owes (Liabilities: payables, short-term debt,
  long-term debt) plus the residual claim of owners (Equity).
\item \textbf{Income Statement}: The flow of revenues and expenses over
  a period, producing Net Income. Key elements include Sales, Cost of
  Goods Sold (COGS), Operating Expenses (OpEx), Depreciation, Interest
  Expense, and Income Tax.
\item \textbf{Cash Budget} (Statement of Cash Flows): The flow of cash
  in and out of the company, decomposed into operating, investing,
  financing, and owner transaction activities. The ending cash balance
  is the \emph{result} of these flows, not an input.
\end{itemize}

\paragraph{Forecast mechanism: driver + policies + Cash Budget.}
The model uses \textbf{Sales Revenue} as the single primary forecast
driver. Given a sales forecast, a set of \emph{policy ratios} determine
every other element:

\begin{itemize}
\item Sales drives Revenue on the Income Statement, which determines
  COGS (via a cost ratio), Receivables, Inventory, Payables, and
  Advance Payments on the Balance Sheet.
\item OpEx, Depreciation, and Interest Expense complete the Income
  Statement, producing Net Income.
\item The Cash Budget routes every revenue, expense, and accrual change
  into the appropriate net liquidity balance (operating, capital
  expense, financing, owner transactions, discretionary). The sum of
  these flows determines the new Cash and Investments in Market
  Securities balances.
\item Debt rollover, equity financing, dividends, and buybacks are
  computed from the Cash Budget's financing and owner transaction
  modules, updating Liabilities and Equity.
\end{itemize}

Because every cash inflow and outflow is explicitly routed through the
Cash Budget, the next year's Balance Sheet is fully determined: no
element is left as a residual plug. The accounting identity
$\text{Assets} = \text{Liabilities} + \text{Equity}$ must emerge from
the correctness of the logic, which is what makes the model auditable
and testable.

\paragraph{Policy variables and training.}
In Pareja~(2009)'s framework, all policy variables are assumed known.
In our case, we do not know the policy variables but know a company's
past financial statements. We train the policy ratios (e.g., \%AR,
\%AP, \%TL, \%Cash, cost ratio, payout ratio, depreciation rate) from
historical data using gradient-based optimization. Simple policies use
static ratios; advanced policies use logit-linear time trends or
Bayesian variational inference to capture evolving firm behavior and
quantify uncertainty.

\paragraph{Refinements to Pareja~(2009).}
We have further refined the framework to better reflect the treasury
strategies of modern multinational corporations like Apple or General
Motors. Where the original model might treat Marketable Securities
solely as a reservoir for excess cash, our implementation models these
securities as a distinct portion of the company's liquidity strategy
\hyperlink{ref:almeida14}{[Almeida(14)]}. This adjustment allows for a more
realistic depiction of how large firms manage capital allocation between
operational cash and interest-bearing assets.

Moreover, Pareja~(2009) assumes long-term debt to always have 10-year
maturity. In reality, various debt products with varying maturity
schedules exist. We, therefore, use an average maturity value as a
variable to be also trained.

Section~\ref{sec:part1} details the time evolution equations for every
financial element, showing precisely how Sales, the policy ratios, and
the Cash Budget connect to produce next year's Balance Sheet and Income
Statement without any plug.

\subsection{Testing and Validation Plan}

Because the model eschews plugs, the fundamental accounting identity
(\emph{Assets = Liabilities + Equity}) becomes a \emph{falsifiable} test
of correctness rather than a guaranteed outcome. A nonzero delta at any
forecast step constitutes a genuine failure, exposing structural errors
that plug-based models would silently absorb. This property is the
foundation of our validation strategy, which combines the accounting
identity invariant with a comprehensive programmatic test suite. The
full testing methodology is detailed in Section~\ref{sec:testing}.

The first three subsections of Part~1 (\emph{The Balance Sheet
Elements}, \emph{Time Evolution of Financial Elements}, and \emph{Time
Series Representation of the Balance Sheet}) answer questions~1 and~2;
the submitted codes answer questions~3 and~4; and the remaining
subsections (\emph{Training the Model}, \emph{Applying Machine
Learning}, and \emph{Forecast Results}) answer questions~5 through~8.
For Part~2, the answers to questions a)~through~i) are directly mapped
to its subsection headings.

\section{Part 1}
\label{sec:part1}

\subsection{The Balance Sheet Elements}

The elements of a balance sheet from \hyperlink{ref:pareja09}{Pareja(09)} are as follows:

Assets

Net fixed assets (NFA)

Advance payments to suppliers for purchases (AdvPP)

Account receivables (AR)

Inventory (Inventory)

Short-term investments (Invest)

Cash equivalent (Cash)

Liabilities

Account payables (AP)

Advance payments from customers for sales (AdvPS)

Short-term debt (STDebt)

Long-term debt (LTDebt)

Equity

Equity invested (EI)

Cumulated retained earnings (CRE)

Cumulated buyback of equity (CBB)

Net Income (NI)

On the balance sheets of real companies, there are likely to be more
elements. For example, there could be current and non-current `other
liabilities' other than the 4 mentioned above, and there could be
various other assets including deferred tax, non-current account
receivables, Goodwill, and other intangible assets. For the purpose of
this simple modeling exercise, we first identify the elements listed in
the above balance sheet, and lump all unaccounted-for elements as shown
below, using Apple's balance sheet as example:

\begin{itemize}
\item
  Assets

  \begin{itemize}
  \item
    All unaccounted-for non-interest-bearing elements are lumped with
    ``non-current assets'' along with net PPE.

    \begin{itemize}
    \item
      Deferred tax, non-current account receivables, goodwill, other
      assets
    \item
      Depreciation rate based on net fixed asset will be reduced
      accordingly.
    \end{itemize}
  \item
    We use other current assets as the proxy for advance payments to
    suppliers for purchases.
  \item
    Short-term and long-term investments in market securities are lumped
    into ``investments in market securities.''

    \begin{itemize}
    \item
      Earn a combined average interest.
    \end{itemize}
  \item
    We also combine cash equivalents with investments in market
    securities to for ``total liquidity'' because companies do not just
    hold onto a pile of cash, but rather invest a significant portion of
    cash in securities to earn interest, which is different from
    \hyperlink{ref:pareja09}{Pareja(09)}'s logic of only investing in excess cash after meeting a
    minimum cash balance.

    \begin{itemize}
    \item
      Total liquidity = cash equivalent + investments in market
      securities
    \end{itemize}
  \end{itemize}
\item
  Liabilities

  \begin{itemize}
  \item
    Advance payments from customers to be fulfilled next year are
    represented as deferred revenue in company balance sheets.
  \item
    All unaccounted-for elements are lumped with long-term debt, to be
    labeled ``non-current liabilities''

    \begin{itemize}
    \item
      Includes non-current long-term debt, employee benefits,
      non-current deferred revenue, non-current accrued expenses,
      long-term capital lease obligation

      \begin{itemize}
      \item
        These other non-current liabilities are like interest-free
        long-term debt
      \item
        The total long-term liability is assumed to have the average
        maturity of AvgM years, and it generates an average LT interest
        (AvgLTInt) owed by the company
      \end{itemize}
    \item
      Something important to note is that the ``long-term debt'' in an
      actual balance sheet EXCLUDES the principal owed next year.

      \begin{itemize}
      \item
        Non-current liabilities\textsubscript{t} = Total LT
        liability\textsubscript{t} * (1 -- 1/AvgM)
      \end{itemize}
    \end{itemize}
  \end{itemize}
\end{itemize}

\begin{quote}
= Total LT liability\textsubscript{t} * (AvgM - 1) / AvgM
\end{quote}

\begin{itemize}
\item
  We group all other liabilities due next year as ``current
  liabilities.''

  \begin{itemize}
  \item
    Includes STDebt, current accrued expenses which include dealer and
    customer allowances, claims and discounts, product warranty,
    payrolls and benefits, and others, to be considered as ``short-term
    loan''

    \begin{itemize}
    \item
      Accrued expenses are like interest-free short-term debt.
    \item
      Short-term loan (STLoan\textsubscript{t}) generates an average ST
      interest (AvgSTInt\textsubscript{t}) owed by the company
    \end{itemize}
  \item
    Also includes current debt which includes principal payment from
    short-term debt and ALSO INCLUDES the current portion of the total
    LT liability

    \begin{itemize}
    \item
      Current liabilities\textsubscript{t} = Total ST
      liability\textsubscript{t} + Total LT liability\textsubscript{t} /
      AvgM
    \end{itemize}
  \end{itemize}
\end{itemize}

\begin{quote}
= Total ST liability\textsubscript{t} + Non-current
liabilities\textsubscript{t} / (AvgM -- 1)
\end{quote}

\begin{itemize}
\item
  Equity

  \begin{itemize}
  \item
    Represented by Stockholders Equity (SE\textsubscript{t})

    \begin{itemize}
    \item
      SE\textsubscript{t} = SE\textsubscript{t-1} + EF\textsubscript{t}
      + NI\textsubscript{t} -- Dividends\textsubscript{t-1} --
      BB\textsubscript{t}

      \begin{itemize}
      \item
        SE\textsubscript{t-1} = stockholder equity last year
      \item
        EF\textsubscript{t} = equity financing to cover portion of CapEx
        spending
      \item
        NI\textsubscript{t} = net income this year
      \item
        Dividends\textsubscript{t-1} = Dividends earned last year, paid
        this year
      \item
        BB\textsubscript{t} = stock buyback this year
      \end{itemize}
    \end{itemize}
  \end{itemize}
\end{itemize}

So, our updated balance sheet, along with subitems for clarity, becomes:

Assets

Non-current assets (\textbf{NCA\textsubscript{t}})

\begin{itemize}
\item
  Net PPE, deferred tax, non-current account receivables, goodwill,
  other assets
\end{itemize}

Advance payments to suppliers for purchases
(\textbf{AdvPP\textsubscript{t}})

Account receivables (\textbf{AR\textsubscript{t}})

Inventory (\textbf{Inv\textsubscript{t}})

Total liquidity (\textbf{TL\textsubscript{t}})

\begin{itemize}
\item
  Investments in market securities (\textbf{IMS\textsubscript{t}})

  \begin{itemize}
  \item
    Short-term and long-term investments
  \end{itemize}
\item
  Cash equivalent (\textbf{Cash\textsubscript{t}})
\end{itemize}

Liabilities

Account payables (\textbf{AP\textsubscript{t}})

Advance payments from customers for sales
(\textbf{AdvPS\textsubscript{t}})

Non-current liabilities (\textbf{NLiab\textsubscript{t}})

\begin{itemize}
\item
  Non-current employee benefits, non-current deferred revenue,
  non-current accrued expenses, long-term debt and capital lease
  obligations MINUS next year's principal, other non-current liabilities
\item
  \textbf{NLiab\textsubscript{t}} = Total LT liability\textsubscript{t}
  * (1 -- 1/\textbf{AvgM})
\end{itemize}

Current liabilities (\textbf{CLiab\textsubscript{t}})

\begin{itemize}
\item
  Consists of the operational short-term obligations (Effective ST Debt)
  plus the principal due this year on long-term loans.
\item
  Effective ST Debt (\textbf{EffSTDebt\textsubscript{t}}): ST debt
  (commercial paper), current accrued expenses, and other current
  liabilities. Two policy variants are implemented:

  \begin{itemize}
  \item
    \textbf{Simple policy} (\texttt{SimpleDebtPolicy}): ST debt is
    deficit-driven: the model borrows short-term only to cover any
    remaining liquidity shortfall after operating and investing
    activities. EffSTDebt\textsubscript{t} = max(0,
    deficit\textsubscript{t}).
  \item
    \textbf{Advanced policy} (\texttt{TrendDebtPolicy}): ST debt is
    modeled as an affine function of sales, capturing the empirical
    pattern that firms like Apple scale short-term borrowing with
    revenue while also maintaining a fixed baseline (e.g., revolving
    facility minimums, notes payable):
    EffSTDebt\textsubscript{t} = EffSTDebt\textsubscript{baseline} +
    Sales\textsubscript{t} $\cdot$ \%STDebt, where both
    EffSTDebt\textsubscript{baseline} and \%STDebt are trainable
    non-negative parameters.
  \end{itemize}
\item
  Current LT Debt (\textbf{CurrLTDebt\textsubscript{t}}): Current
  portion of the long-term debt and capital lease obligations
  principals.

  \begin{itemize}
  \item
    CurrLTDebt\textsubscript{t} = Total LT liability\textsubscript{t} /
    AvgM = NLiab\textsubscript{t} / (AvgM -- 1)
  \end{itemize}
\item
  Total: CLiab\textsubscript{t} =
  EffSTDebt\textsubscript{t} + CurrLTDebt\textsubscript{t}
\end{itemize}

Equity

Stockholder equity (\textbf{SE\textsubscript{t}})

\begin{itemize}
\item
  \textbf{SE\textsubscript{t}} = \textbf{SE\textsubscript{t-1}} +
  \textbf{EF\textsubscript{t}} + \textbf{NI\textsubscript{t}} --
  \textbf{Dividends\textsubscript{t-1}} -- \textbf{BB\textsubscript{t}}

  \begin{itemize}
  \item
    SE\textsubscript{t-1} : total stockholder equity last year
  \item
    EF\textsubscript{t} : equity financing to cover portion of CapEx
    spending
  \item
    NI\textsubscript{t} : net income this year
  \item
    Dividends\textsubscript{t-1} = Dividends earned last year, paid this
    year
  \item
    BB\textsubscript{t} : stock buyback this year
  \end{itemize}
\end{itemize}

\subsection{Time Evolution of Financial Elements}

We refer to \hyperlink{ref:pareja09}{Pareja(09)} for constructing the mathematical equations that
link all the elements on the balance sheet without circularity or plugs.
First, we define a single primary driver that will predict the evolution
of the balance sheet's elements.

\begin{itemize}
\item
  Sales\textsubscript{t} : the total revenue generated by the company in
  year t
\end{itemize}

This is a different approach from \hyperlink{ref:pareja09}{Pareja(09)} where they modeled the unit
price, the unit cost, and the volume sold based on market research data
on volume's price sensitivity. For the scope of this problem set, we
lump these variables into an overarching sales driver whose evolution is
modeled independently based on historical data. Sales serves as the primary driver because it is the most fundamental measure of a firm's economic activity: all other financial statement items are either costs of generating revenue, assets deployed to support revenue, or financing decisions to fund those assets. In the corporate finance literature, revenue-driven forecasting is standard practice \hyperlink{ref:koller20}{[Koller(20)]}. The limitation of this approach is that it does not model the underlying demand drivers (price elasticity, market size, competitive dynamics) that determine revenue growth. The sales forecast itself is generated via linear extrapolation of historical trends, a pragmatic choice given the 4--8 years of available data, though macroeconomic-sensitive models (e.g., LSTM with GDP and inflation inputs) would improve robustness.

\subsubsection{Forecast Flow Overview}

Before diving into the equations for each balance sheet element, it is
helpful to trace how a single period's sales forecast ripples through
all three financial statements to produce next year's Balance Sheet. The
flow is as follows.

\paragraph{Income Statement.}
Revenue equals Sales\textsubscript{t} at the top of the Income
Statement. From Revenue we subtract the operating costs:

\begin{itemize}
\item Cost of Goods Sold (COGS)
\item Operating Expenses (OpEx, e.g., administrative and sales)
\item Depreciation and amortization
\end{itemize}

This yields EBIT (Earnings Before Interest and Taxes). Stated
sequentially:
\[
\text{EBIT} = \text{Revenue} - \text{COGS} - \text{OpEx} - \text{Depreciation}
\]
Next, we account for interest payments on outstanding debt and returns
on short-term investments in market securities (see
\eqref{eq:ms_return_derivation}), producing EBT (Earnings Before Tax):
\[
\text{EBT} = \text{EBIT} - \text{Interest} + \text{MSReturn}
\]
Income tax is assessed as a fixed percentage of EBT, giving Net Income:
\[
\text{Net Income} = \text{EBT} \cdot (1 - \text{\%IT})
\]
From Net Income, a fraction is allocated as next year's dividends, and
the remainder accumulates into Retained Earnings (a component of
Equity).

\paragraph{Cash Budget.}
Net Income is an accrual concept and does not correspond directly to
cash generated. To track actual liquidity, we build a Cash Budget that
records cash movements when they physically occur, independent of
revenue recognition timing. The Cash Budget therefore serves the same
role as the Statement of Cash Flows but is built prospectively: rather
than reconciling Net Income to cash after the fact, it schedules every
expected cash receipt and disbursement directly.

Unlike the standard three-section Statement of Cash Flows, we decompose
financing activities into two separate modules, namely External Financing
(debt) and Transactions with Owners (equity and dividends), and add a
fifth module for Discretionary Transactions in market securities,
reflecting how these are managed as a policy allocation rather than a
residual. Each module produces a Net Liquidity Balance (NLB):

\begin{itemize}
\item \textbf{Operating Activities}: Inflows from sales (including
  collections on receivables and advance payments) minus outflows for
  purchases, OpEx, and income tax.
\item \textbf{Investing Activities}: Capital expenditure outflows
  (CapEx) for fixed asset investment.
\item \textbf{External Financing}: Inflows from new short-term and
  long-term borrowing, minus outflows for principal repayments and
  interest expense.
\item \textbf{Transactions with Owners}: Inflows from equity financing,
  minus outflows for dividends paid and stock buybacks.
\item \textbf{Discretionary Transactions}: Returns from short-term
  investments in market securities. In our model, ST investments are a
  policy allocation of total liquidity rather than an excess-cash
  destination, so this module has no discretionary outflow.
\end{itemize}

The sum of the five NLBs determines the change in total liquidity (Cash
+ Investments in Market Securities) from one period to the next.

\paragraph{Closing the loop.}
With the Income Statement producing Net Income (and hence Retained
Earnings) and the Cash Budget producing the change in total liquidity,
every element of next year's Balance Sheet is determined:

\begin{itemize}
\item \textbf{Assets} evolve via working-capital ratios (AR, AP, Inv,
  AdvPP, AdvPS) applied to next year's Sales and Purchases, the NCA
  roll-forward (CapEx minus Depreciation), and the updated liquidity
  balance from the Cash Budget.
\item \textbf{Liabilities} evolve via the debt rollover schedule (new
  borrowing from External Financing, principal payments, amortization
  of long-term debt into current portion) and working-capital ratios
  for AP and AdvPS.
\item \textbf{Equity} evolves via the Retained Earnings accrual from
  Net Income, minus dividends and buybacks, plus any new equity
  financing.
\end{itemize}

No element is left as a residual plug. That is, no line item is set to
whatever makes the balance sheet balance, which would otherwise mask
errors in the model. Instead, the accounting identity
$\text{Assets} = \text{Liabilities} + \text{Equity}$ emerges
automatically at $t+1$ as a direct consequence of double-entry
bookkeeping: every cash or accrual transaction affects at least two
accounts by equal and opposite amounts, so the identity is structurally
guaranteed, not verified after the fact.

The remainder of this section provides the detailed equations for each
financial element: the individual asset line items, the Cash Budget,
the liability line items, Equity, and the Income Statement.

\subsubsection{Non-Current Assets (NCA)}

The balance of non-current assets evolves sequentially as existing fixed
assets depreciate over their useful lifecycle and new capital
expenditures (CapEx) are capitalized. The fundamental roll-forward
equation is defined as:

\begin{equation}\label{eq:nca_rollforward}
\text{NCA}_t = \text{NCA}_{t-1} - \text{Depreciation}_t + \text{CapEx}_t
\end{equation}

Depreciation is modeled as a percentage of the beginning-of-period
non-current asset base:

\begin{equation}\label{eq:depreciation}
\text{Depreciation}_t = \text{NCA}_{t-1} \cdot \text{\%Depreciation}
\end{equation}

Note: Because our model aggregates all Non-Current Assets, which may
include non-depreciable assets such as land or goodwill, the trained
effective depreciation rate (\%Depr) will naturally be lower than the
strict accounting depreciation rate applied purely to Plant, Property,
and Equipment (PP\&E).

To model CapEx, we diverge from a rigid steady-state assumption where
maintenance perfectly offsets depreciation. Instead, we decompose CapEx
into a maintenance component (scaled by a trainable parameter, \%AM, to
reflect partial, full, or inflationary replacement) and a growth
component driven by sales expansion:

\begin{equation}\label{eq:capex}
\text{CapEx}_t = \text{\%AM} \cdot \text{Depreciation}_t + \text{\%AG} \cdot \text{Sales}_t
\end{equation}

By substituting \eqref{eq:depreciation} and \eqref{eq:capex} into \eqref{eq:nca_rollforward}, we can express
the net change in non-current assets (ΔNCA\textsubscript{t}) purely in
terms of our modeled variables. The growth component (\%AG $\times$ Sales) assumes constant capital intensity of incremental revenue, a reasonable first approximation for asset-light technology firms where growth CapEx scales with business volume. For capital-intensive industries, a non-linear formulation (e.g., step functions for capacity expansion) may be more appropriate. During model training, the optimizer
learns the parameters for the effective depreciation rate (\%Depr), the
maintenance replacement ratio (\%AM), and the incremental asset growth
coefficient (\%AG).

\subsubsection{Cost of Goods Sold (COGS) and Inventory}

Instead of modeling purchases independently and using the Cost of Goods
Sold (COGS) as a plug, we directly model COGS based on historical margin
trends. We define the Cost Ratio (CR) as the proportion of sales
consumed by COGS:

\begin{equation}\label{eq:cost_ratio}
\text{CR}_t = \text{COGS}_t / \text{Sales}_t
\end{equation}

To ensure our predicted cost ratio remains strictly bounded between 0
and 1 (a cost ratio outside this range would imply negative costs or costs exceeding revenue, neither of which is sustainable), we transform the historical cost ratio using the logit function
and model it as a linear function of time (t). The logit-linear parameterization allows for gradual trend shifts. For example, a declining cost ratio may reflect economies of scale, supply chain efficiencies, or pricing power gains. The linear time trend is a parsimonious choice; in practice, cost ratios may exhibit mean reversion as competitive pressure limits sustained margin expansion.

\begin{equation}\label{eq:cost_ratio_logit}
\text{logit}(\text{CR}_t) = \ln \left( \frac{\text{CR}_t}{1 - \text{CR}_t} \right) = \alpha_{\text{CR}} + \beta_{\text{CR}} \cdot t
\end{equation}

The parameters $\alpha$\textsubscript{CR} and $\beta$\textsubscript{CR} are trained
on historical data by minimizing the mean squared error between the
historical logit(CR) and the predicted logit(CR). Once the future Cost
Ratio is predicted, COGS is calculated directly:

\begin{equation}\label{eq:cogs}
\text{COGS}_t = \text{CR}_t \cdot \text{Sales}_t
\end{equation}

Next, inventory can be modeled as a function of last year's inventory,
this year's purchases, and this year's COGS:

\begin{equation}\label{eq:inventory_rollforward}
\text{Inv}_t = \text{Inv}_{t-1} + \text{Purchases}_t - \text{COGS}_t
\end{equation}

We set the inventory target as a percentage of cost of goods sold, rather than sales. Inventory is carried at cost (under FIFO or weighted-average accounting), so it should scale with the cost to produce goods, not the selling price. This is equivalent to modeling Days Sales of Inventory: $\text{\%Inv}_{\text{COGS}} = \text{DSI} / 365$.

\begin{equation}\label{eq:inventory_target}
\text{Inv}_t = \text{\%Inv}_{\text{COGS}} \cdot \text{COGS}_t = \text{\%Inv}_{\text{COGS}} \cdot \text{CR}_t \cdot \text{Sales}_t
\end{equation}

Since Inv\textsubscript{t} and COGS\textsubscript{t} are both independently forecast, Purchases\textsubscript{t} is fully determined by rearranging the inventory identity (7):

\begin{equation}\label{eq:purchases}
\text{Purchases}_t = \text{Inv}_t - \text{Inv}_{t-1} + \text{COGS}_t
\end{equation}

\begin{equation}\label{eq:purchases_expanded}
\text{Purchases}_t = \text{\%Inv} \cdot (\text{Sales}_t - \text{Sales}_{t-1}) + \text{COGS}_t
\end{equation}

Rather than being forecast independently, Purchases\textsubscript{t} is derived directly from the structural accounting identity: it is the quantity of goods that must have been purchased to arrive at the target inventory level given the period's COGS. This is distinct from a balance-sheet plug: no residual is being forced to reconcile a constraint; the identity holds by construction, and Purchases\textsubscript{t}​ is simply its remaining free variable.

We add \%Inventory (\%Inv\textsubscript{COGS}), $\alpha$\textsubscript{CR} and $\beta$\textsubscript{CR}
to the variables to be trained.

\subsubsection{Advance Payments for Purchases}

Advance payments to suppliers for purchases (AdvPP) can be modeled as
strictly proportional to the total purchases made during the current
year. This approach relies on the simplifying assumption that current
purchasing volumes serve as a reliable proxy for near-term future
demand. Therefore, we define:

\begin{equation}\label{eq:advpp}
\text{AdvPP}_t = \text{\%AdvPP} \cdot \text{Purchases}_t
\end{equation}

We add \%Advance payments to suppliers (\%AdvPP) to the variables to
learn in our training. Advance payments are computed as a percentage of the current period's purchases (not next period's), avoiding forward-looking bias in the time-series model.

\subsubsection{Account Receivables}

Next, we have account receivables. A fraction of this year's sales to
customers are on credit and is receivable next year:

\begin{equation}\label{eq:account_receivables}
\text{AR}_t = \text{\%AR} \cdot \text{Sales}_t
\end{equation}

We add \%Account receivables (\%AR) to the variables to learn in our
training. Note that \%AR is equivalent to DSO/365, where DSO (Days Sales Outstanding) measures the average collection period. Similarly, the inventory ratio \%Inv\textsubscript{COGS} corresponds to DSI/365 (Days Sales of Inventory), and the payables ratio \%AP corresponds to DPO/365 (Days Payable Outstanding). Together, these three metrics define the Cash Conversion Cycle: $\text{CCC} = \text{DSO} + \text{DSI} - \text{DPO}$, a standard measure of working capital efficiency in corporate finance.

\subsubsection{Total Liquidity}
\label{sec:total_liquidity_policy}

Finally, the last element in our assets is total liquidity (TL).
Following \hyperlink{ref:pareja09}{Pareja(09)}, we model the total liquidity available as a
function of sales, split between cash and investments in market
securities (IMS). Two policy variants are implemented:

\paragraph{Simple policy (\texttt{SimpleLiquidityPolicy}).}
Total liquidity is a static percentage of sales, with a fixed
cash/liquidity split:

\begin{equation}\label{eq:total_liquidity}
\text{TL}_t = \text{\%TL} \cdot \text{Sales}_t
\end{equation}

\begin{equation}\label{eq:cash}
\text{Cash}_t = \text{\%Cash} \cdot \text{TL}_t
\end{equation}

\begin{equation}\label{eq:ims}
\text{IMS}_t = (1 - \text{\%Cash}) \cdot \text{TL}_t
\end{equation}

The simple policy has a notable limitation: it assumes that total
liquidity scales linearly with sales at a fixed ratio, which is
unrealistic for cash-rich firms whose liquidity trajectory is driven by
strategic capital-return programs rather than by operational needs. Apple
is a clear example. Its cash balance and stock buybacks over the past
decade have been governed primarily by a multi-year buyback program
designed to draw down excess cash, not by any stable relationship to
revenue. A pure sales-proportional model cannot capture this secular
decline in the cash-to-sales ratio.

\paragraph{Advanced policy (\texttt{TrendLiquidityPolicy}).}
To address this, both the total liquidity ratio and the cash allocation
evolve over time via logit-linear trends, with a baseline floor that is
independent of sales:

\begin{equation}\label{eq:total_liquidity_trend}
\text{TL}_t = \beta_{\text{base}} + \text{Sales}_t \cdot \sigma(\alpha_{\text{TL}} + \beta_{\text{TL}} \cdot t)
\end{equation}

\begin{equation}\label{eq:cash_trend}
\text{Cash}_t = \text{TL}_t \cdot \sigma(\alpha_{\text{Cash}} + \beta_{\text{Cash}} \cdot t)
\end{equation}

The baseline floor $\beta_{\text{base}}$ captures the minimum liquidity
a firm maintains regardless of revenue (e.g., strategic cash reserves),
while the logit-linear terms allow the sales-proportional component and
the cash/IMS split to shift gradually over time.

We add the liquidity policy parameters to our variables to be trained.

Before we move onto liabilities, it is helpful to generate a Cash
Budget to model the flow of cash in and out of the liquidity balance
(cash and short-term investments in market securities).

\subsubsection{Cash Budget}

We can break down flows into/out of liquidity balance into net liquidity
balances (NLB) pertaining to the following activities: operating,
capital expense, financing, external investment, and transaction with
owners. The final liquidity is the sum of all NLBs.

\paragraph{Operating NLB}

Operating activity captures all inflows from sales (current year sales
revenue, account receivables from last year, advance payments from
customers for next year) and all outflows (current year purchases,
account payables from last year, advance payments to suppliers for next
year, operating expenses, income tax):

\begin{equation}\label{eq:operating_inflows}
\text{Inflows}_t = \text{AR}_{t-1} + \text{AdvPS}_t + \text{Sales}_t \cdot (1 - \text{\%AR} - \text{\%AdvPS})
\end{equation}

\begin{equation}\label{eq:operating_outflows}
\text{Outflows}_t = \text{AP}_{t-1} + \text{AdvPP}_t + \text{Purchases}_t \cdot (1 - \text{\%AP} - \text{\%AdvPP}) + \text{OpEx}_t + \text{IT}_t
\end{equation}

\begin{equation}\label{eq:operating_nlb}
\text{Operating NLB}_t = \text{Inflows}_t - \text{Outflows}_t
\end{equation}

where OpEx\textsubscript{t} is the operating expense, and
IT\textsubscript{t} is income tax. We can model OpEx\textsubscript{t} as
a function of sales and inflation:

\begin{equation}\label{eq:opex}
\text{OpEx}_t = \text{\%OR} \cdot \text{Sales}_t + \text{OB}_{t-1} \cdot (1 + \text{Inflation}_t)
\end{equation}

where OB is the base operating expense that increases with inflation
each year (rent, insurance), and \%OR is the rate of increase
proportional to this year's sales revenue that impacts the variable
portion of the operating expense. Inflation data can be obtained from
U.S. Bureau of Labor Statistics.

We can model IT\textsubscript{t} as a percentage of net income
(NI\textsubscript{t}):

\begin{equation}\label{eq:income_tax}
\text{IT}_t = \text{\%IT} \cdot \text{NI}_t / (1 - \text{\%IT}) = \text{NI}_t / (1/\text{\%IT} - 1)
\end{equation}

We add \%IT, OB\textsubscript{T\_start} (base operating expense of the
first element in a time series) and \%OR to our variables to be trained.

\paragraph{Capital Expense NLB}

Investment is made into the non-current assets to:

\begin{enumerate}
\def\labelenumi{\alph{enumi})}
\item
  keep non-current assets price the same by accounting for depreciation
\item
  grow non-current assets by some percentage.
\end{enumerate}

We have already done the calculation in \eqref{eq:capex}, which we reproduce
here:

$\text{CapEx}_t = \text{Depreciation}_t + \text{Sales}_t \cdot \text{\%Asset growth}$ from \eqref{eq:capex}

\begin{equation}\label{eq:capex_nlb}
\text{CapEx NLB}_t = -\text{CapEx}_t = -\text{Depreciation}_t - \text{Sales}_t \cdot \text{\%Asset growth}
\end{equation}

\paragraph{Financing NLB}

This is the net liquidity balance after obtaining more liquidity from
new short-term loan (STLoan\textsubscript{t}) and new long-term loan
(LTLoan\textsubscript{t}), and paying all debt obligations this year.

\begin{equation}\label{eq:financing_nlb}
\text{Financing NLB}_t = \text{STLoan}_t + \text{LTLoan}_t - \text{CLiab}_{t-1} - \text{Avg.\ ST interest}_t - \text{Avg.\ LT interest}_t
\end{equation}

where current liabilities\textsubscript{t-1} (CLiab\textsubscript{t-1})
are the principals from short-term liabilities and long-term liabilities
paid this year, represented by:

\begin{equation}\label{eq:cliab_prev}
\text{CLiab}_{t-1} = \text{STLoan}_{t-1} + \text{NLiab}_{t-1} / (\text{AvgM} - 1)
\end{equation}

The average short-term interest (AvgSTInt\textsubscript{t}) and average
long-term interest (AvgLTInt\textsubscript{t}) owed are calculated as
follows:

\begin{equation}\label{eq:avg_st_interest}
\text{AvgSTInt}_t = \text{\%AvgSTInt} \cdot \text{STLoan}_{t-1} = \text{\%AvgSTInt} \cdot (\text{CLiab}_{t-1} - \text{NLiab}_{t-1} / (\text{AvgM} - 1))
\end{equation}

\begin{equation}\label{eq:avg_lt_interest}
\text{AvgLTInt}_t = \text{\%AvgLTInt} \cdot \text{Total LT liability}_{t-1} = \text{\%AvgLTInt} \cdot (\text{NLiab}_{t-1} / (1 - 1/\text{AvgM}))
\end{equation}

We add AvgM, \%AvgSTInt and \%AvgLTInt to our variables to learn. We
derive the expression for STLoan\textsubscript{t} and
LTLoan\textsubscript{t} in the \emph{Liabilities} subsection.

\paragraph{External Investment NLB}

We redeem external investments from last year with an effective
percentage of return (\%MSReturn, derived from
$(\text{EBT} - \text{EBIT} + \text{InterestExpense}) / \text{IMS}$
historically, see \eqref{eq:ms_return_derivation}), and make new investments this year:

\begin{equation}\label{eq:ext_invest_nlb}
\text{External investment NLB}_t = \text{IMS}_{t-1} \cdot (1 + \text{\%MSReturn}) - \text{IMS}_t
\end{equation}

Instead of only allocating the excess cash in market securities, we
invest a proportion of total liquidity balance. Using \eqref{eq:ims}:

\begin{equation}\label{eq:ext_invest_nlb_expanded}
\text{External investment NLB}_t = (1 - \text{\%Cash}) \cdot (\text{TL}_{t-1} \cdot (1 + \text{\%MSReturn}) - \text{TL}_t)
\end{equation}

We add \%MSReturn to our training variables to be learned.

\paragraph{Transaction with Owners NLB}

In \hyperlink{ref:pareja09}{Pareja(09)}'s model, owners can make additional investments to finance a
portion of long-term debt, which counts as an inflow to liquidity. The
company pays dividends to owners and can also engage in stock buybacks:

\begin{equation}\label{eq:owners_nlb}
\text{Owners transaction NLB}_t = \text{Equity financing}_t - \text{Dividends}_{t-1} - \text{Stock buybacks}_t
\end{equation}

where

\begin{equation}\label{eq:equity_financing}
\text{Equity financing}_t = \text{LTLoan}_t \cdot \text{\%EF} / (1 - \text{\%EF})
\end{equation}

The equity financing proportion is trained from historical data, implicitly capturing the firm's revealed preference between debt and equity. This is consistent with the pecking order hypothesis \hyperlink{ref:myers84}{[Myers(84)]}, which predicts firms prefer internal funds, then debt, then equity, resulting in a low and relatively stable \%EF for profitable firms. The trained \%EF thus reflects the firm's historical position on the pecking order rather than an arbitrary split.

Dividends are modeled with two policy variants:\label{sec:dividends}

\paragraph{Simple policy (\texttt{SimpleDividendPolicy}).}
Dividends are a fixed percentage of prior-year net income, with no
memory of past payouts:

\begin{equation}\label{eq:target_dividends}
\text{Dividends}_{t-1} = \text{\%Payout ratio} \cdot \text{Net income}_{t-1}
\end{equation}

\paragraph{Advanced policy (\texttt{LintnerDividendPolicy}).}
Rather than calculating dividends strictly as a fixed percentage of
current net income, we apply the Lintner partial-adjustment model to
account for dividend smoothing. Corporate management typically prefers
stable dividend payouts to avoid signaling financial distress to
investors. Therefore, the modeled dividend is represented as a weighted
average between a target payout (driven by a baseline \%Payout\_Ratio)
and the actual dividends paid in the prior period, regulated by an
optimized Adjustment\_Speed parameter:

\begin{equation}\label{eq:lintner_dividends}
\text{Dividends}_{t-1} = \text{Adj.\ Speed} \cdot (\text{\%PR} \cdot \text{NI}_{t-1}) + (1 - \text{Adj.\ Speed}) \cdot \text{Dividends}_{t-2}
\end{equation}

The \hyperlink{ref:lintner56}{Lintner(56)} partial-adjustment model remains empirically robust for mature dividend-paying firms \hyperlink{ref:brav05}{[Brav(05)]}. The key economic insight is that managers treat dividends as a commitment: cutting dividends signals financial distress, so firms smooth payouts to avoid involuntary future cuts. For mega-cap technology firms, the adjustment speed tends to be low (conservative smoothing), as these firms have ample retained earnings and prefer share buybacks for flexible capital return. The trained Adjustment\_Speed parameter captures this firm-specific conservatism.

Stock buybacks are also modeled with two policy variants:

\paragraph{Simple policy (\texttt{SimpleBuybackPolicy}).}
Buybacks are proportional to depreciation, using the asset base as a
scaling proxy:

\begin{equation}\label{eq:stock_buybacks}
\text{Stock buybacks}_t = \text{\%BB} \cdot \text{Depreciation}_t
\end{equation}

\paragraph{Advanced policy (\texttt{BaselineBuybackPolicy}).}
An affine model that adds a fixed baseline component, capturing the
observation that large firms often maintain a minimum buyback program
independent of asset turnover:

\begin{equation}\label{eq:stock_buybacks_advanced}
\text{Stock buybacks}_t = \beta_{\text{base}} + \beta_{\text{ratio}} \cdot \text{Depreciation}_t
\end{equation}

In both variants, any excess cash remaining after meeting the liquidity
target is also distributed as additional buybacks.

For companies like Apple, their revenue generates tremendous amount of
cash that they sometimes use the excess cash to buy back shares to
spread the earnings to shareholders instead of paying a higher dividend
which has tax implications. We add \%Equity financing (\%EF), \%Payout
ratio (\%PR) and \%Buyback ratio (\%BB) to our variables to learn. We
will define the equations for Net income\textsubscript{t} in the
\emph{Income Statement} subsection.

\paragraph{Final NLB}

We obtain the final net liquidity balance at time $t$ by summing
all the different activities' NLBs:

\begin{multline}\label{eq:final_nlb}
\text{Final NLB}_t = \text{Operating NLB}_t + \text{CapEx NLB}_t + \text{Financing NLB}_t \\
+ \text{External investment NLB}_t + \text{Owners transaction NLB}_t
\end{multline}

By construction, the Final NLB is the period's net cash inflow (the
change in total liquidity from year $t-1$ to year $t$), so it must
equal the difference between this year's and last year's total
liquidity balances:

\begin{equation}\label{eq:final_nlb_tl_closure}
\text{Final NLB}_t = \text{TL}_t - \text{TL}_{t-1}
\end{equation}

This identity is what closes the loop and makes training possible.
Because the liquidity policy (Section~\ref{sec:total_liquidity_policy})
gives us an independent target for $\text{TL}_t$ via
\eqref{eq:total_liquidity_trend}, namely
$\text{TL}_t = \text{TL}_{\text{baseline}} + \text{Sales}_t \cdot
\sigma(\alpha_{\text{TL}} + \beta_{\text{TL}} \cdot t)$. Substituting
this into \eqref{eq:final_nlb_tl_closure} yields a single scalar
constraint that ties together every policy ratio appearing in the
five NLBs. Specifically, the only free variable on the left-hand side
is $\text{LTLoan}_t$ (buried inside $\text{Financing NLB}_t$), which
can therefore be solved for endogenously as the residual financing
needed to bring total liquidity to its policy target. Once
$\text{LTLoan}_t$ is determined, every remaining quantity is fully
specified and the rest of the balance sheet assembles mechanically.
No plug is required; the ``balancing'' is done by a financing
decision that has real economic meaning (how much long-term debt to
raise this period), not by absorbing arbitrary residuals into cash.

\subsubsection{Account Payables}

Account payables are modeled as a percentage of total purchases made
this year, payable next year to the suppliers:

\begin{equation}\label{eq:account_payables}
\text{AP}_t = \text{\%AP} \cdot \text{Purchases}_t
\end{equation}

We add \%AP to our variables to be trained.

\subsubsection{Advance Payments for Future Sales}

Advance payments received from customers for future sales
(AdvPS\textsubscript{t}) are modeled as a direct proportion of the
current year's total sales. This assumes that the volume
of upfront customer payments scales linearly with overall revenue
generation. Therefore, we define:

\begin{equation}\label{eq:advps}
\text{AdvPS}_t = \text{\%AdvPS} \cdot \text{Sales}_t
\end{equation}

We add \%AdvPS to our variables to be trained.

\subsubsection{Current Liabilities}

To model the current liabilities, we first model the effective
short-term loan (STLoan\textsubscript{t}) issued this year. Short-term
loan covers any excess cash needed to cover operational expenses and
last year's short-term loan obligation (principal + interest) due this
year and bring the total liquidity to this year's liquidity target. We
also account for the return from investment in market securities here
because the size of investment is directly proportional to cash,
\textbf{instead of} investing excess cash after ST and LT loan
calculations which is done in \hyperlink{ref:pareja09}{Pareja(09)}:

\begin{multline}\label{eq:st_cash_deficit}
\text{Short-term cash deficit}_t = \text{TL}_t - \text{TL}_{t-1} - \text{Operating NLB}_t \\
+ \text{STLoan}_{t-1} \cdot (1 + \text{\%AvgSTInt}) - \text{\%MSReturn} \cdot \text{IMS}_{t-1}
\end{multline}

\begin{equation}\label{eq:stloan}
\text{STLoan}_t = \max(0,\; \text{Short-term cash deficit}_t)
\end{equation}

Then, we calculate the effective long-term loan
(LTLoan\textsubscript{t}) issued this year to finance
CapEx\textsubscript{t}, last year's current portion of the long-term
liability, interest on long-term liability, and dividends to bring the
liquidity to this year's target after accounting for the equity
financing:

\begin{multline}\label{eq:ltloan}
\text{LTLoan}_t = \max(0,\; \text{Short-term cash deficit}_t - \text{STLoan}_t - \text{CapEx NLB}_t \\
+ \text{NLiab}_{t-1} / (\text{AvgM} - 1) + \text{AvgLTInt}_t + \text{Dividends}_{t-1} + \text{Stock buybacks}_t) \cdot (1 - \text{\%EF})
\end{multline}

Current liabilities are summation of the effective short-term loan and
the current portion of non-current liabilities, and are represented by:

\begin{equation}\label{eq:cliab}
\text{CLiab}_t = \text{STLoan}_t + \text{Total LT liability}_t / \text{AvgM}
\end{equation}

where

\begin{equation}\label{eq:total_lt_liability}
\text{Total LT liability}_t = \text{LTLoan}_t + \text{NLiab}_{t-1}
\end{equation}

\subsubsection{Modeling Effective Short-Term Debt}

The simple deficit-driven formulation of short-term borrowing does not
fit the historical behavior of cash-rich firms like Apple. The
following subsubsection motivates the advanced (logit-linear) policy
variant and derives the metric used throughout the forecasting model.

\paragraph{The Problem}

A strictly deficit-driven approach calculates new short-term loans as a
reaction to cash shortfalls. This forces short-term borrowing to zero
when a company has surplus liquidity. However, in reality, large
corporations (e.g., Apple) strategically carry significant commercial
paper and revolving credit lines (\$45B--\$88B) despite holding massive
cash reserves.

This means the short-term cash deficit is less than zero, i.e., cash surplus, which would result in no short-term borrowing for any year.


\paragraph{Defining the Metric}

Effective short-term debt isolates the operational and short-term
financing liabilities that scale directly with the size of the business,
stripping out strict trade payables and the current portion of long-term
debt.

The balance sheet identity is defined as:

\begin{multline}\label{eq:eff_st_debt_identity}
\text{Effective ST Debt}_t = \text{Current Liabilities}_t - \text{Accounts Payable}_t \\
- \text{Advance Payments Sales}_t - \text{Current LT Debt}_t
\end{multline}

Which equates to the core components we are modeling:

\begin{equation}\label{eq:eff_st_debt_components}
\text{Effective ST Debt}_t = \text{ST Debt}_t + \text{Accrued Expenses}_t + \text{Other Current Liabilities}_t
\end{equation}

\paragraph{The Solution \& Implementation}

To reflect actual corporate financing behavior, effective short-term
debt is now modeled as a policy-driven metric rather than a
deficit-driven reaction. It is projected as an affine function of
sales, consisting of a fixed baseline and a sales-proportional
component:

\begin{equation}\label{eq:eff_st_debt_calc}
\text{Effective ST Debt}_t = \text{EffSTDebt}_{\text{baseline}} + \text{Sales}_t \cdot \text{\%STDebt}
\end{equation}

Both $\text{EffSTDebt}_{\text{baseline}}$ and \%STDebt are trainable
parameters constrained to be non-negative via a softplus bijector. The
baseline component captures fixed short-term obligations that do not
scale from zero with sales (e.g., revolving facility minimums, notes
payable commitments), while the sales-proportional component captures
working-capital-driven borrowing. This affine form was chosen over a
pure sales ratio because the latter would asymptote to zero at zero
sales (unrealistic for firms that carry a standing level of short-term
paper) and over a logit-linear form because the latter saturates as
sales grow, which is inappropriate for firms whose borrowing scales
proportionally with revenue.

\paragraph{Liquidity Cascade} Because short-term debt is now an
  established structural policy, it acts as an \emph{input} to the
  long-term liquidity deficit. The calculated effective short-term debt
  funds a portion of the required capital, shifting the role of the
  purely deficit-driven ``plug'' entirely to the long-term financing
  variable.

\paragraph{Non-current Liabilities}

Non-current liabilities include the new effective
LTLoan\textsubscript{t} issued this year, and last year's non-current
liabilities minus the current portion of this year's non-current
liabilities.

\begin{equation}\label{eq:nliab}
\text{NLiab}_t = \text{Total LT liability}_t \cdot (1 - 1/\text{AvgM})
\end{equation}

\subsubsection{Equity}

This year's equity accounts for the additional equity financing this
year, net income from this year, dividends paid out this year, and stock
buyback this year:

\begin{equation}\label{eq:stockholder_equity}
\text{SE}_t = \text{SE}_{t-1} + \text{EF}_t + \text{NI}_t - \text{Dividends}_{t-1} - \text{BB}_t
\end{equation}

where

\begin{equation}\label{eq:ef}
\text{EF}_t = \text{LTLoan}_t / (1 - \text{\%EF}) \cdot \text{\%EF}
\end{equation}

\begin{equation}\label{eq:dividends}
\text{Dividends}_{t-1} = \text{\%PR} \cdot \text{NI}_{t-1}
\end{equation}

\begin{equation}\label{eq:bb}
\text{BB}_t = \text{\%BB} \cdot \text{\%Depreciation} \cdot \text{NCA}_{t-1} + \text{any excess cash after meeting liquidity target}
\end{equation}

We have most of what we need to learn the training variables except
dividends and income tax, which will come from modeling the net income
in the income statement.

\subsubsection{Income Statement: EBITDA and EBT}

We have all the components needed from the previous sections to
construct the elements related to net income.

EBITDA\textsubscript{t} (earnings before interest, taxes, depreciation,
and amortization) can be calculated by:

\begin{equation}\label{eq:ebitda}
\text{EBITDA}_t = \text{Sales}_t - \text{COGS}_t - \text{OpEx}_t
\end{equation}

Earnings before tax is calculated by subtracting interest payments,
depreciation (and amortization):

\begin{equation}\label{eq:ebt}
\text{EBT}_t = \text{EBITDA}_t - \text{Depreciation}_t - \text{AvgSTInt}_t - \text{AvgLTInt}_t + \text{\%MSReturn} \cdot \text{IMS}_{t-1}
\end{equation}

\subsubsection{Derivation of Historical \%MSReturn}

The term $\text{\%MSReturn} \cdot \text{IMS}_{t-1}$ represents income
from short-term investments in market securities.  In practice, not all
companies separately disclose this line item in their income statement;
many lump it with other non-operating items (foreign exchange
gains/losses, miscellaneous items, etc.) into a single ``Other
income/(expense), net'' (OI\&E) line.  We therefore derive the
historical values from the income statement identity:

\begin{equation}\label{eq:ms_return_derivation}
\text{MSReturn}_t = \text{EBT}_t - \text{EBIT}_t + \text{InterestExpense}_t
\end{equation}

This isolates all non-operating income/expense \emph{except} interest
expense (which the model already handles via \%AvgSTInt and \%AvgLTInt).
For companies like Apple, where investment income dominates OI\&E, this
residual is a close proxy for the true return on the securities
portfolio.  For companies with significant foreign exchange exposure, it
may also absorb FX gains/losses.

When the company's 10-K filing does not provide Interest Expense separately (as is
the case for some fiscal years), we fall back to
$\text{MSReturn}_t = \text{EBT}_t - \text{EBIT}_t$ and set
interest expense to zero for those years, which lumps all non-operating
items into MSReturn but preserves the EBT identity.

\paragraph{Forward projection.}
During forecasting, the model projects $\text{MSReturn}_t$ as a fixed
percentage return on the prior period's IMS balance:
$\text{\%MSReturn} \cdot \text{IMS}_{t-1}$.  The effective rate
\%MSReturn is learned during training as the historical average of
$\text{MSReturn}_t / \text{IMS}_{t-1}$, which absorbs the blended
contribution of investment income and smaller non-operating items into a
single rate parameter. This fixed-percentage return assumption is a simplification: for firms with large treasury portfolios, investment returns depend on duration, credit quality, and market conditions. The trained parameter captures the historical average net return after FX effects and mark-to-market adjustments.

\subsubsection{Net Income and Dividends}

Net income (NI\textsubscript{t}) is determined by subtracting the income
tax:

\begin{equation}\label{eq:net_income}
\text{NI}_t = \text{EBT} \cdot (1 - \text{\%IT})
\end{equation}

Next year's dividends were already modeled in the dividends section (p.~\pageref{sec:dividends}).

We now have everything we need to represent every element in the balance
sheet as a time series.

\subsection{Time Series Representation of the Balance Sheet}

We re-write every element in the balance sheet as a time series that
have been simplified by combining the equations above. The primary
drivers are Sales\textsubscript{t} and Purchases\textsubscript{t}, which
drive every element given the variables to be trained are learned.

\paragraph{Variables to be trained}

\begin{itemize}
\item
  \%AG\, [$>0$] : Asset growth rate as a percentage of non-current assets
\item
  \%AM\, [$>0$] : Asset maintenance ratio as a percentage of depreciation
\item
  \%Depr\, [$>0$] : Depreciation rate as a percentage of non-current assets
\item
  \%AdvPS\, [$>0$] : Advance payments from customers as a percentage of next
  year's sales
\item
  \%AdvPP\, [$>0$] : Advance payments to suppliers as a percentage of next year's
  purchases
\item
  \%AR\, [$>0$] : Account receivables as a percentage of this year's sales
\item
  \%AP\, [$>0$] : Account payables as a percentage of this year's purchases
\item
  \%Inv\, [$>0$] : Inventory target as a percentage of this year's COGS
\item
  $\alpha$\textsubscript{CR}, $\beta$\textsubscript{CR}\, [$\in \mathbb{R}$] :
  Variables that make up the Cost Ratio via logit model

  \begin{itemize}
  \item
    logit(CR\textsubscript{t}) = ln ( CR\textsubscript{t} / (1 --
    CR\textsubscript{t}) ) = $\alpha$\textsubscript{CR} + $\beta$\textsubscript{CR} *
    t
  \end{itemize}
\item
  TL\textsubscript{baseline}, $\alpha_{\text{TL}}$, $\beta_{\text{TL}}$\, [$\in \mathbb{R}$] :
  Baseline floor and logit-linear trend parameters for total liquidity
  with a sales-proportional component:
  $\text{TL}_t = \text{TL}_{\text{baseline}} + \text{Sales}_t \cdot
  \sigma(\alpha_{\text{TL}} + \beta_{\text{TL}} \cdot t)$
\item
  $\alpha_{\text{Cash}}$, $\beta_{\text{Cash}}$\, [$\in \mathbb{R}$] :
  Logit-linear trend parameters for the cash share of total liquidity:
  $\text{Cash}_t = \text{TL}_t \cdot \sigma(\alpha_{\text{Cash}} +
  \beta_{\text{Cash}} \cdot t)$
\item
  \%IT\, [$\in(0,1)$] : Income tax as a percentage of net income this year
\item
  \%OR\, [$\in \mathbb{R}$] : Variable portion of OpEx that increases in proportion to this
  year's sales
\item
  OB\textsubscript{T\_start}\, [$\in \mathbb{R}$] : Baseline OpEx that increases with this
  year's average inflation
\item
  EffSTDebt\textsubscript{baseline}\, [$>0$] and \%STDebt\, [$>0$] : Baseline and
  sales-proportional parameters for calculating effective short-term
  debt (Effective ST Debt\textsubscript{t} = EffSTDebt\textsubscript{baseline}
  + Sales\textsubscript{t} $\cdot$ \%STDebt)
\item
  \%AvgSTInt\, [$>0$] : Average short-term interest on the effective short-term
  loan
\item
  \%AvgLTInt\, [$>0$] : Average interest on total long-term liabilities
\item
  AvgM\, [$>1.001$] : Average maturity of total long-term liability
\item
  \%MSReturn\, [$\in \mathbb{R}$] : Effective percentage return on the investment in market
  securities, derived from historical $(\text{EBT} - \text{EBIT} + \text{InterestExpense}) / \text{IMS}$
\item
  $\alpha_{\text{EF}}$, $\beta_{\text{EF}}$\, [$\in \mathbb{R}$] :
  Logit-linear parameters for the equity financing mix,
  $\text{\%EF}_t = \sigma(\alpha_{\text{EF}}
  + \beta_{\text{EF}} \cdot t)$, where \%EF is the fraction of long-term
  financing raised as equity rather than debt
\item
  \%PR\, [$\in(0,1)$] : Target dividend payout ratio (fraction of net income) in the
  Lintner partial-adjustment model
\item
  Adj.\ Speed\, [$\in(0,1)$] : Lintner partial-adjustment speed controlling how
  quickly actual dividends converge to the target payout
  $\text{\%PR} \cdot \text{NI}_{t-1}$
\item
  BB\textsubscript{baseline}, BB\textsubscript{ratio}\, [$\in \mathbb{R}$] :
  Fixed baseline and depreciation-proportional parameters for stock buyback
  in the affine buyback policy
  ($\text{BB}_t = \text{BB}_{\text{baseline}} + \text{BB}_{\text{ratio}}
  \cdot \text{Depreciation}_t$)
\end{itemize}

\paragraph{Parameter Constraints and Bounds}

To ensure the model generates economically logical and mathematically
stable financial statements, strict constraints are enforced on both the
trainable variables and the underlying equations. These constraints are
implemented natively within the optimizer using TensorFlow Probability
Bijectors.

\subsubsection{Domain Constraints on Trainable Variables}

\begin{itemize}
\item
  \textbf{Strict Positivity ($x > 0$):} Variables that
  cannot logically be negative --- the Depreciation Rate (\%Depr),
  Asset Growth (\%AG), Asset Maintenance (\%AM), Interest Rates
  (\%AvgSTInt, \%AvgLTInt), all working-capital ratios (\%AdvPS,
  \%AdvPP, \%AR, \%AP, \%Inv), and the Effective Short-Term Debt
  parameters (EffSTDebt\textsubscript{baseline}, \%STDebt) --- are
  constrained using a Softplus activation function. Because both
  ST-debt parameters are strictly positive, the linear form
  $\text{EffSTDebt}_{\text{baseline}} + \text{Sales}_t \cdot
  \text{\%STDebt}$ is guaranteed non-negative for any positive sales
  level. This prevents the optimizer from testing negative values
  during gradient descent.
\item
  \textbf{Percentage Bounds ($0 < x < 1$):}
  Variables representing fractions of a whole, such as the Income Tax
  Percentage (\%IT), Dividend Payout Ratio (\%PR), and the Lintner
  Dividend Adjustment Speed ($\alpha$), are constrained
  using a Sigmoid function.
\item
  \textbf{Maturity Boundary ($\text{AvgM} > 1.001$):} Because
  the calculation for the current portion of long-term debt divides by
  $(\text{AvgM} - 1)$, the average maturity must strictly exceed 1 year to
  prevent mathematical singularity and infinite liability generation.
  This is enforced using a chained bijector: a Softplus function shifted
  by a constant of 1.001.
\item
  \textbf{Logit-Linear Trends ($\sigma(\alpha + \beta t) \in (0,1)$):}
  Time-varying policies --- the Cost Ratio (CR), Total Liquidity
  percentage (\%TL), Cash share of Liquidity (\%Cash), and Equity
  Financing mix (\%EF) --- are modeled as logit-linear trends.
  By wrapping the linear trend $(\alpha + \beta t)$ in a sigmoid
  function, the forecasted ratios are strictly bounded between 0 and 1,
  even if extrapolated decades into the future. The underlying
  $\alpha$ and $\beta$ parameters are unconstrained, giving the
  optimizer full freedom in logit space.
\end{itemize}

\subsubsection{Logical Constraints on Equations}

\begin{itemize}
\item
  \textbf{Non-Negative Borrowing:} A company cannot issue ``negative''
  debt to reduce cash. In the Cash Budget equations, the issuance
  of new short-term and long-term loans is constrained using a maximum
  function: $\text{New\_Loan}_t = \max(0,
  \text{Liquidity\_Deficit}_t)$. If there is a cash surplus,
  borrowing falls to exactly zero, and the excess cash cascades into
  stock buybacks.
\item
  \textbf{Structural Identity Check:} The fundamental accounting
  equation ($\text{Assets} = \text{Liabilities} + \text{Equity}$) is
  strictly evaluated at every time step $t$. Rather than enforcing
  this via a loss function penalty (which would allow minor deviations),
  it is enforced structurally by the double-entry routing of the Cash
  Budget. The simulation calculates a delta to verify that the mismatch
  is exactly zero.
\end{itemize}

\paragraph{Variables from External Environment}

Since we assume a flat short-term and long-term interest for loans and
return on market securities, the only external variable is inflation.
This flat-rate assumption is a first-order simplification. In practice, borrowing costs depend on the firm's leverage and credit quality \hyperlink{ref:merton74}{[Merton(74)]}. As the forecast shows rising debt levels, a more complete model would incorporate credit spread widening, a natural extension that would link the cost of debt to the evolving debt-to-equity ratio. The current approach is conservative: it understates the cost of increased leverage, producing optimistic net income projections for high-debt scenarios.

\begin{itemize}
\item
  Inflation\textsubscript{t} : Average inflation at year t
\end{itemize}

\subsubsection{Element-Wise Time-Series Equations}

The following equations express every balance-sheet and income-statement
element as a function of the trained variables,
Sales\textsubscript{t}, Purchases\textsubscript{t}, and the previous
period's state. Together they define the single-period state transition
used by both the simple forward model and the gradient-based training
loop.

\paragraph{Assets}

\begin{itemize}
\item
  Non-current assets

  \begin{itemize}
  \item
    NCA\textsubscript{t} = NCA\textsubscript{t-1} --
    Depreciation\textsubscript{t} + CapEx\textsubscript{t}

    \begin{itemize}
    \item
      Where CapEx\textsubscript{t} = Depreciation\textsubscript{t} +
      Sales\textsubscript{t} * \%AG
    \end{itemize}
  \end{itemize}
\item
  Advance payments to suppliers for purchases

  \begin{itemize}
  \item
    AdvPP\textsubscript{t} = \%AdvPP * Purchases\textsubscript{t}
  \end{itemize}
\item
  Account receivables

  \begin{itemize}
  \item
    AR\textsubscript{t} = \%AR * Sales\textsubscript{t}
  \end{itemize}
\item
  Inventory

  \begin{itemize}
  \item
    Inv\textsubscript{t} = \%Inv * Sales\textsubscript{t}
  \end{itemize}
\item
  Total liquidity (\textbf{TL\textsubscript{t}})

  \begin{itemize}
  \item
    \textit{Simple:}\;
    Target TL\textsubscript{t} = \%TL * Sales\textsubscript{t},\;
    Cash\textsubscript{t} = \%Cash * TL\textsubscript{t},\;
    IMS\textsubscript{t} = (1 -- \%Cash) * TL\textsubscript{t}
  \item
    \textit{Advanced:}\;
    Target TL\textsubscript{t} =
    TL\textsubscript{baseline} + Sales\textsubscript{t} *
    $\sigma(\alpha_{\text{TL}} + \beta_{\text{TL}} \cdot t)$

    \begin{itemize}
    \item
      Cash\textsubscript{t} = TL\textsubscript{t} *
      $\sigma(\alpha_{\text{Cash}} + \beta_{\text{Cash}} \cdot t)$
    \item
      IMS\textsubscript{t} = TL\textsubscript{t} --
      Cash\textsubscript{t}
    \end{itemize}
  \item
    Should be equivalent to: Final TL\textsubscript{t} = Operating
    NLB\textsubscript{t} + CapEx NLB\textsubscript{t} + Financing
    NLB\textsubscript{t} + External investment NLB\textsubscript{t} +
    Owners transactions NLB\textsubscript{t}
  \end{itemize}
\end{itemize}

\paragraph{Liabilities}

\begin{itemize}
\item
  Account payables

  \begin{itemize}
  \item
    AP\textsubscript{t} = \%AP * Purchases\textsubscript{t}
  \end{itemize}
\item
  Advance payments from customers for sales

  \begin{itemize}
  \item
    AdvPS\textsubscript{t} = \%AdvPS * Sales\textsubscript{t}
  \end{itemize}
\item
  Effective short-term debt

  \begin{itemize}
  \item
    \textit{Simple (deficit-driven):}\;
    EffSTDebt\textsubscript{t} = max(0,
    LiqDeficit\textsubscript{ST,t})
  \item
    \textit{Advanced (policy-driven):}\;
    EffSTDebt\textsubscript{t} = EffSTDebt\textsubscript{baseline} +
    Sales\textsubscript{t} * \%STDebt
  \end{itemize}
\item
  Current portion of long-term debt

  \begin{itemize}
  \item
    CurLTDebt\textsubscript{t} = (NewLTLoan\textsubscript{t} +
    NCL\textsubscript{t-1}) / AvgM
  \end{itemize}
\end{itemize}

The short-term liquidity deficit drives financing. In the simple model
the deficit is absorbed by ST borrowing; in the advanced model it flows
entirely to long-term financing:

\begin{itemize}
\item
  LiqDeficit\textsubscript{ST,t} = TL\textsubscript{t} --
  (Cash\textsubscript{t-1} + IMS\textsubscript{t-1}) -- Operating
  NLB\textsubscript{t} + EffSTDebt\textsubscript{t-1} * (1 +
  \%AvgSTInt)
\item
  Operating NLB\textsubscript{t} = AR\textsubscript{t-1} +
  AdvPS\textsubscript{t} + Sales\textsubscript{t} * (1 - \%AR - \%AdvPS)
  -- (AP\textsubscript{t-1} + AdvPP\textsubscript{t} +
  Purchases\textsubscript{t} * (1 - \%AP - \%AdvPP) +
  OpEx\textsubscript{t} + IT\textsubscript{t})
\item
  AP\textsubscript{t-1} = \%AP * Purchases\textsubscript{t-1}
\item
  AdvPP\textsubscript{t} = \%AdvPP * Purchases\textsubscript{t}
\item
  OpEx\textsubscript{t} = \%OR * Sales\textsubscript{t} +
  OB\textsubscript{t-1} * (1 + Inflation\textsubscript{t})
\item
  IT\textsubscript{t} = NI\textsubscript{t} / (1/\%IT - 1)
\item
  AR\textsubscript{t-1} = \%AR * Sales\textsubscript{t-1}
\item
  AdvPS\textsubscript{t} = \%AdvPS * Sales\textsubscript{t}
\end{itemize}

\begin{itemize}
\item
  Non-current liabilities

  \begin{itemize}
  \item
    NCL\textsubscript{t} = (NewLTLoan\textsubscript{t} +
    NCL\textsubscript{t-1}) * (AvgM -- 1) / AvgM
  \end{itemize}
\end{itemize}

\begin{quote}
where
\end{quote}

\begin{itemize}
\item
  LiqDeficit\textsubscript{LT,t} = LiqDeficit\textsubscript{ST,t} --
  EffSTDebt\textsubscript{t} -- MSReturn\textsubscript{t} +
  CapEx\textsubscript{t} + CurLTDebt\textsubscript{t-1} +
  InterestLT\textsubscript{t} + Dividends\textsubscript{t} + Stock
  buybacks\textsubscript{t}
\item
  NewLTLoan\textsubscript{t} = max(0,
  LiqDeficit\textsubscript{LT,t}) * (1 -- \%EF\textsubscript{t})

  \begin{itemize}
  \item
    \textit{Simple:}\; \%EF\textsubscript{t} = \%EF \quad (static)
  \item
    \textit{Advanced:}\; \%EF\textsubscript{t} =
    $\sigma(\alpha_{\text{EF}} + \beta_{\text{EF}} \cdot t)$
  \end{itemize}
\item
  CapEx\textsubscript{t} = NCA\textsubscript{t-1} * \%Depr + \%AG *
  Sales\textsubscript{t}
\item
  Dividends\textsubscript{t}:

  \begin{itemize}
  \item
    \textit{Simple:}\;
    Dividends\textsubscript{t} = \%PR * NI\textsubscript{t-1}
  \item
    \textit{Advanced (Lintner):}\;
    Dividends\textsubscript{t} = Adj.\ Speed * (\%PR *
    NI\textsubscript{t-1}) + (1 -- Adj.\ Speed) *
    Dividends\textsubscript{t-1}
  \end{itemize}
\item
  NI\textsubscript{t-1} : net income from last year
\end{itemize}

\begin{quote}
= (1 -- \%IT) * (Sales\textsubscript{t-1} -- COGS\textsubscript{t-1} --
OpEx\textsubscript{t-1} -- Depreciation\textsubscript{t-1} --
InterestST\textsubscript{t-1} -- InterestLT\textsubscript{t-1} +
MSReturn\textsubscript{t-1})
\end{quote}

\begin{itemize}
\item
  MSReturn\textsubscript{t} = \%MSReturn * IMS\textsubscript{t-1}
\item
  COGS\textsubscript{t-1} = \%Inv * (Sales\textsubscript{t-2} --
  Sales\textsubscript{t-1}) + Purchases\textsubscript{t-1}
\item
  OpEx\textsubscript{t-1} = \%OR * Sales\textsubscript{t-1} +
  OB\textsubscript{t-2} * (1 + Inflation\textsubscript{t-1})
\item
  Depreciation\textsubscript{t} = NCA\textsubscript{t-1} * \%Depr
\item
  InterestST\textsubscript{t} = \%AvgSTInt *
  EffSTDebt\textsubscript{t-1}
\item
  InterestLT\textsubscript{t} = \%AvgLTInt * (NCL\textsubscript{t-1} +
  CurLTDebt\textsubscript{t-1})
\item
  Stock buybacks\textsubscript{t}:

  \begin{itemize}
  \item
    \textit{Simple:}\;
    \%BB * Depreciation\textsubscript{t} + excess cash after meeting
    liquidity target
  \item
    \textit{Advanced:}\;
    BB\textsubscript{baseline} + BB\textsubscript{ratio} *
    Depreciation\textsubscript{t} + excess cash after meeting liquidity
    target
  \end{itemize}
\end{itemize}

\paragraph{Equity}

\begin{itemize}
\item
  Stockholder equity

  \begin{itemize}
  \item
    SE\textsubscript{t} = SE\textsubscript{t-1} +
    EquityFinancing\textsubscript{t} + NI\textsubscript{t} --
    Dividends\textsubscript{t} -- Stock buybacks\textsubscript{t}
  \end{itemize}
\end{itemize}

\begin{quote}
where EquityFinancing\textsubscript{t} = max(0,
LiqDeficit\textsubscript{LT,t}) * \%EF\textsubscript{t}
(static or logit-linear; see NewLTLoan above)
\end{quote}

We see in all elements' equations that they are driven by Sales and
Purchases and elements from previous years, indicating no circularity or
plugs. The net liquidity balance calculation and the new short-term and
long-term borrowing equations ensure that the assets and liabilities +
equity are increased/decreased the same. Thus, assets = liabilities +
equity will always be satisfied. We can check with test values in our
simple TensorFlow model.

\subsection{Forward Model Implementation}

This forward model has been implemented in Python in
\texttt{run\_simple\_model\_forecast.py}, which composes
the modular policy classes from the \texttt{financial\_forecast/} package
using data from Apple's financial
statements. The demo below uses the \textbf{simple} policy variants
throughout (static ratios, deficit-driven ST debt, fixed payout
dividends); the advanced policies are exercised in the trained models of
Section~2.5. The financial statements were first obtained from Yahoo
finance as suggested by Question 4 (we wrote the script in
\texttt{yahoofinance.py}), but we could only access past 4 years of
data, which is not sufficient for training in the following section. We
managed to obtain older financial data by directly pulling from Apple's
10-K filings. The script that we wrote and used for the data extraction
is found in \texttt{run\_build\_financial\_statements\_sec.py}.

The following is the output of the code when run with
\texttt{python run\_simple\allowbreak{}\_model\allowbreak{}\_forecast\allowbreak{}\_demo.py},
which uses Apple's FY2024 balance sheet as the initial state,
approximates every policy variable from the most recent year's ratios,
holds them constant, and simply iterates four annual forecast steps:


\begin{table}[H]
\centering
\small
\begin{tabular}{l rrrr}
\toprule
\textbf{Line Item} & \textbf{FY2025} & \textbf{FY2026} & \textbf{FY2027} & \textbf{FY2028} \\
\midrule
\multicolumn{5}{l}{\textbf{Assets}} \\
Non-Current Assets      & 211.2 & 211.0 & 210.9 & 210.8 \\
Adv Payments (Purch)    &  17.5 &  18.4 &  19.3 &  20.2 \\
Accounts Receivable     &  67.4 &  70.9 &  74.4 &  77.9 \\
Inventory               &   6.5 &   6.9 &   7.2 &   7.5 \\
Cash                    &  42.7 &  44.9 &  47.1 &  49.3 \\
Invest in Mkt Sec       &   0.0 &   0.0 &   0.0 &   0.0 \\
\textbf{Total Assets}   & \textbf{345.3} & \textbf{352.1} & \textbf{358.9} & \textbf{365.7} \\
\addlinespace
\multicolumn{5}{l}{\textbf{Liabilities}} \\
Accounts Payable        &  71.2 &  74.7 &  78.4 &  82.1 \\
Adv Payments (Sales)    &   9.0 &   9.4 &   9.9 &  10.4 \\
Effective ST Debt       &   0.0 &   0.0 &   0.0 &   0.0 \\
Current LT Debt         &  13.5 &  14.5 &  15.4 &  16.2 \\
Non-Current Liabilities & 193.4 & 207.7 & 220.3 & 230.9 \\
\textbf{Total Liabilities} & \textbf{287.1} & \textbf{306.4} & \textbf{324.1} & \textbf{339.5} \\
\addlinespace
\textbf{Equity}         & \textbf{58.2} & \textbf{45.7} & \textbf{34.8} & \textbf{26.2} \\
\addlinespace
\textbf{Total Liab + Equity} & \textbf{345.3} & \textbf{352.1} & \textbf{358.9} & \textbf{365.7} \\
\midrule
\multicolumn{5}{l}{\textbf{Income Statement}} \\
Net Income              &  74.6 &  80.8 &  83.7 &  86.7 \\
\addlinespace
CHECK: Assets$-$(L+E)   &   0.0 &   0.0 &   0.0 &   0.0 \\
\bottomrule
\end{tabular}
\par\smallskip
{\footnotesize All figures in \$ billions.}
\caption{Four-year balance sheet and income statement forecast (FY\,2025--2028) for Apple, using constant single-year policy ratios derived from FY\,2024.}
\label{tab:simple_forecast_demo}
\end{table}

Even without using any plugs, we see that the assets = liabilities +
equity identity is satisfied because our model follows the double-entry
principle with any inflows and outflows on the balance sheet as done by
\hyperlink{ref:pareja09}{Pareja(09)}.

\subsection{Training the Model}

The model's trainable parameters are split into two groups and trained
in two sequential phases. \textbf{Policy-phase} parameters directly
govern individual working-capital, liquidity, OpEx, dividend, and
buyback line items, and can each be anchored to a corresponding
observable series in the historical financial statements; they are
trained independently by minimizing mean-squared error between
predicted and observed values. \textbf{Structural-phase} parameters
(interest rates, debt maturity, market securities return, equity
financing mix) affect Net Income, Equity, and the liability side of
the Balance Sheet only through the full multi-year forecast rollout
via \texttt{forecast\_step} transitions; they are therefore trained
jointly after the policy phase, with the policy parameters held fixed.

The primary forecast driver for both phases is:

\begin{itemize}
\item Sales\textsubscript{t} : Total revenue (Income Statement)
\end{itemize}

\paragraph{Policy-phase variables.}
Grouped by the policy module that owns them:

\begin{itemize}
\item \textbf{Non-current assets (\texttt{CapexPolicy}).}
  \%Depr (effective depreciation rate as a fraction of beginning NCA),
  \%AM (maintenance CapEx multiplier on Depreciation), and \%AG
  (growth CapEx as a fraction of Sales). Learned from historical NCA,
  Depreciation, and Sales.

\item \textbf{Cost ratio (\texttt{TrendCostRatioPolicy}).}
  $\alpha_{\text{CR}}$, $\beta_{\text{CR}}$ --- logit-linear cost ratio
  parameters such that
  $\text{CR}_t = \sigma(\alpha_{\text{CR}} + \beta_{\text{CR}} \cdot t)$.
  Learned from historical COGS and Sales.

\item \textbf{Working capital (\texttt{WorkingCapitalPolicy}).}
  \%AR (AR / Sales), \%AP (AP / Purchases), \%AdvPS (advance payments
  from customers / Sales), \%AdvPP (advance payments to suppliers /
  Purchases), and \%Inv (Inventory / COGS). Each learned from its
  respective historical ratio.

\item \textbf{Total liquidity (\texttt{TrendLiquidityPolicy}).}
  $\text{TL}_{\text{baseline}}$, $\alpha_{\text{TL}}$,
  $\beta_{\text{TL}}$ for total liquidity with a baseline floor and a
  logit-linear sales-proportional component:
  $\text{TL}_t = \text{TL}_{\text{baseline}} + \text{Sales}_t \cdot
  \sigma(\alpha_{\text{TL}} + \beta_{\text{TL}} \cdot t)$. Plus
  $\alpha_{\text{Cash}}$, $\beta_{\text{Cash}}$ for the cash share of
  liquidity:
  $\text{Cash}_t = \text{TL}_t \cdot
  \sigma(\alpha_{\text{Cash}} + \beta_{\text{Cash}} \cdot t)$. Learned
  from historical Cash, short-term investments, and Sales.

\item \textbf{Effective short-term debt (\texttt{TrendDebtPolicy}).}
  $\text{EffSTDebt}_{\text{baseline}}$ and \%STDebt for the affine ST
  debt policy:
  $\text{EffSTDebt}_t = \text{EffSTDebt}_{\text{baseline}} +
  \text{Sales}_t \cdot \text{\%STDebt}$. Learned from historical
  Effective ST Debt and Sales.

\item \textbf{Operating expense (\texttt{SimpleOpEx}).}
  \%OR (variable OpEx rate as a slope on sales deviation from the
  historical mean) and $\text{OB}_{T\_\text{start}}$ (baseline OpEx,
  accumulated with cumulative inflation). Learned from historical
  OpEx, Sales, and inflation.

\item \textbf{Income tax (\texttt{SimpleTax}).}
  \%IT --- effective income tax rate. Learned from historical Net
  Income and income tax provision.

\item \textbf{Dividends (\texttt{LintnerDividendPolicy}).}
  \%PR (target payout ratio as a fraction of Net Income) and
  Adj.\ Speed (Lintner partial-adjustment speed: how quickly actual
  dividends converge to the target). Learned from historical dividend
  payouts and Net Income.

\item \textbf{Stock buyback (\texttt{BaselineBuybackPolicy}).}
  $\text{BB}_{\text{baseline}}$ (fixed baseline buyback amount) and
  $\text{BB}_{\text{ratio}}$ (multiplier on Depreciation for the
  variable component), giving an affine buyback model
  $\text{BB}_t = \text{BB}_{\text{baseline}} + \text{BB}_{\text{ratio}}
  \cdot \text{Depreciation}_t$. Learned from historical repurchases
  and Depreciation.
\end{itemize}

Each variable in this group can be trained independently against its
corresponding observable series using mean-squared error minimization
under Adam via TensorFlow's \texttt{GradientTape}.

\paragraph{Structural-phase variables.}
Trained jointly after the policy phase converges:

\begin{itemize}
\item \%AvgSTInt : Average interest rate on Effective Short-Term Debt.

\item \%AvgLTInt : Average interest rate on total long-term liabilities.

\item AvgM : Average maturity of the long-term debt portfolio.

\item \%MSReturn : Effective return on Investments in Market
  Securities, derived from historical
  $(\text{EBT} - \text{EBIT} + \text{InterestExpense}) / \text{IMS}$.

\item $\alpha_{\text{EF}}$, $\beta_{\text{EF}}$ : Logit-linear equity
  financing mix such that
  $\text{\%EF}_t = \sigma(\alpha_{\text{EF}} + \beta_{\text{EF}} \cdot t)$,
  where \%EF is the fraction of long-term financing raised as equity
  rather than debt.
\end{itemize}

These parameters affect Net Income, Equity, and the liability side of
the Balance Sheet only through the forward rollout, so they are trained
concurrently by rolling out the full forecast and minimizing MSE
between predicted and observed Net Income, Stockholders' Equity,
Current Liabilities, and Non-current Liabilities.

The code has been implemented in \texttt{run\_trainable\allowbreak{}\_model\allowbreak{}\_forecast\allowbreak{}\_simple\allowbreak{}\_policies.py},
which composes the \texttt{TrainableFinancialModel} with simple policy modules
and runs the \texttt{PolicyTrainer} and \texttt{StructuralTrainer}. For
this simple training model, \%AM = 100\% and Purchases are also used
together with Sales as dual driver.

In order to test the model's prediction capability, we train the
parameters from Apple's 2018\textasciitilde2024 data and leave 2025 for backtesting, and forecast 9 more years to observe the trend. The following is the output after
running \texttt{run\_trainable\allowbreak{}\_model\allowbreak{}\_forecast\allowbreak{}\_adv\allowbreak{}\_policies.py}:


\begin{lstlisting}[style=console, caption={Trained policy parameters for the advanced-policy model}, label=lst:params_simple_adv]
Simple parameters:
Final %AG: 0.00000
Final %AM: 0.94058
Final %Depr: 0.05552
Final %AdvPS: 0.02126
Final %AdvPP: 0.07840
Final %AR: 0.15866
Final %AP: 0.30796
Final %Inv (COGS): 0.02939
Total Liquidity (baseline + logit-linear): baseline=0.6012, alpha=-2.3463, beta=-0.411515
  => %TL at t=0: 0.0874, %TL at t=6: 0.0080
Cash % of Liquidity (logit-linear): alpha=-0.2235, beta=0.032375
  => %Cash at t=0: 0.4444, %Cash at t=6: 0.4927
Final %IT: 0.17023
Final %PR: 0.99973
Final DivAdjSpeed: 0.00370
Stock Buyback (baseline + ratio*depr): baseline=1.5612, ratio=-6.674733
Effective ST Debt (baseline + % of sales): baseline=0.0000, pct=0.17869
Cost Ratio (logit-linear): alpha=0.3459, beta=-0.047524
  => CR at t=0: 0.5856, CR at t=6: 0.5152
OpEx Variable %: 0.0939
OpEx Baseline (USD): 4.02e+10
\end{lstlisting}

\noindent\rule{\textwidth}{0.4pt}

\begin{lstlisting}[style=console, caption={Trained structural parameters for the advanced-policy model}, label=lst:params_structural_adv]
Structural parameters:

Final %AvgSTInt: 0.00003
Final %AvgLTInt: 0.01687
Final %MSReturn: 0.07588
Equity Financing (logit-linear): alpha=0.4894, beta=-0.006361
Avg Maturity Years: 15.3410
\end{lstlisting}

\noindent\rule{\textwidth}{0.4pt}


\begin{table}[H]
\centering
\small
\vspace{0.5em}

\begin{tabular}{l rrr}
\toprule
\textbf{Line Item} & \textbf{Actual} & \textbf{Predicted} & \textbf{Error \%} \\
\midrule
\multicolumn{4}{l}{\textbf{Assets}} \\
Non-Current Assets        & 211.3 & 211.3 & $+$0.0\% \\
Adv Payments (Purchases)  &  14.6 &  16.3 & $+$12.0\% \\
Accounts Receivable       &  73.0 &  66.0 & $-$9.5\% \\
Inventory                 &   5.7 &   6.2 & $+$7.7\% \\
Cash                      &  35.9 &  31.2 & $-$13.1\% \\
Invest in Mkt Securities  &  18.8 &  31.1 & $+$65.9\% \\
\textbf{Total Assets}     & \textbf{359.2} & \textbf{362.2} & $+$0.8\% \\
\addlinespace
\multicolumn{4}{l}{\textbf{Liabilities}} \\
Accounts Payable          &  69.9 &  64.2 & $-$8.2\% \\
Adv Payments (Sales)      &   9.1 &   8.8 & $-$2.3\% \\
Effective ST Debt         &  74.4 &  74.4 & $+$0.0\% \\
Non-Current Liabilities   & 119.9 & 125.5 & $+$4.6\% \\
\textbf{Equity}           & \textbf{73.7} & \textbf{80.6} & $+$9.3\% \\
\midrule
\multicolumn{4}{l}{\textbf{Income Statement}} \\
COGS                      & 209.3 & 209.5 & $+$0.1\% \\
OpEx                      &  62.2 &  59.5 & $-$4.3\% \\
Depreciation              &  11.7 &  11.8 & $+$0.6\% \\
Interest Payments         &   0.0 &   2.4 & N/A \\
ST Investment Returns     &  $-$0.3 &   2.7 & $+$932.8\% \\
Tax                       &  20.7 &  23.1 & $+$11.5\% \\
\textbf{Net Income}       & \textbf{112.0} & \textbf{112.6} & $+$0.6\% \\
\addlinespace
\multicolumn{4}{l}{\textbf{Cash Flow}} \\
Dividends                 &  15.4 &  15.5 & $+$0.6\% \\
Stock Buyback             &  90.7 &  77.6 & $-$14.5\% \\
\bottomrule
\end{tabular}
\par\smallskip
{\footnotesize All figures in \$ billions.}
\caption{Backtest of the trained advanced-policy model against Apple's actual FY\,2025 results (out-of-sample). Error percentages are relative to actual values.}
\label{tab:backtest_advpolicies}
\end{table}

\begin{figure}[H]\centering
\includegraphics[width=\textwidth]{report_latex_media/media/advpolicies_result_1.png}
\caption{Ten-year forecast from the trained advanced-policy model (1 of 3).}
\label{fig:advpolicies_forecast}
\end{figure}

\begin{figure}[H]\ContinuedFloat\centering
\includegraphics[width=\textwidth]{report_latex_media/media/advpolicies_result_2.png}
\caption[]{Advanced-policy forecast (2 of 3).}
\end{figure}

\begin{figure}[H]\ContinuedFloat\centering
\includegraphics[width=\textwidth]{report_latex_media/media/advpolicies_result_3.png}
\caption[]{Advanced-policy forecast (3 of 3).}
\end{figure}

The backtest shows mixed but informative results. On the income
statement side, the model predicts Net Income within 0.6\% of the
actual figure, and COGS and Depreciation are nearly exact ($+$0.1\%
and $+$0.6\%), suggesting that the trained cost ratio and CapEx
parameters generalize well out of sample. Balance sheet aggregates
are also close: Total Assets err by only $+$0.8\%.

The larger item-level errors reflect known modelling simplifications.
We have lumped various balance sheet elements together to be more
consistent with the models described in
\hyperlink{ref:pareja07}{Pareja(07)} and
\hyperlink{ref:pareja09}{Pareja(09)}, as explained in the
\emph{Balance Sheet Elements} subsection, so individual line items
such as Accounts Receivable ($-$9.5\%) and Accounts Payable
($-$8.2\%) absorb the aggregation mismatch even when the totals
are accurate.

The Interest Payments actual value is reported as zero because
Apple's FY2025 10-K does not disclose a single comparable interest
expense figure in the format our extraction pipeline expects, so the
error percentage is not meaningful. Similarly, the ST Investment
Returns actual value is $-$\$0.3B (a small net loss on the
portfolio), while the model predicts \$2.7B, yielding a $+$932.8\%
error. This is driven by the same liquidity-split limitation
discussed below.

The largest outlier is Investment in Market Securities ($+$65.9\%).
This is a direct consequence of the logit-linear liquidity policy:
the trend model captures the smooth, secular decline in Apple's
total liquidity ratio but cannot reproduce the jagged year-to-year
reallocation between cash and market securities, which is driven by
discrete treasury decisions (e.g., shifting maturities in the bond
portfolio) rather than a smooth function of sales or time. Cash
itself is underestimated by 13.1\% for the same reason. Because
the model overestimates the IMS balance, it also overestimates
the return generated on those securities, explaining the ST
Investment Returns miss. A richer liquidity model that
distinguishes active rebalancing from trend-driven allocation
would improve these individual estimates, though the total
liquidity error remains modest.


\subsection{Applying Machine Learning}

As our model is rather a simple model based on \hyperlink{ref:pareja09}{Pareja(09)}, there is a lot
of improvements that could make the model more performant. A short list
of potential improvements is as follows:

\begin{itemize}
\item
  Predicting purchases and sales of future years using Long Short-Term
  Memory and macroeconomic indicators (e.g., inflation, interest rate)
\item
  Implementing Bayesian model to account for probability distribution of
  accounting variables (e.g., company policy for capital expenditure,
  operating expense prediction, cash balance targets) to account for the
  noise term, n(t).
\item
  Implementing neural networks for each accounting variables that take
  into consideration macroeconomic signals (e.g., incorporate consumer
  price index, GDP growth, and interest rates to predict advance
  purchases percentage)
\item
  Incorporating the company's stock buyback target if it has been
  disclosed publicly into training the buyback ratio with a neural
  network
\item
  Propagating uncertainty across all key drivers (sales growth, cost ratios, interest rates, CapEx intensity) via full Monte Carlo simulation, not just OpEx. Adding scenario analysis (base case, downside with revenue contraction, upside with margin expansion) would strengthen the forecast's decision-support value.
\item
  Computing projected financial ratios (debt-to-equity, interest coverage, current ratio) from the forecast output and flagging years where they breach investment-grade thresholds, enabling proactive capital structure management.
\end{itemize}

For our model, we choose to implement Bayesian learning with Variational
Inference method \hyperlink{ref:durr20}{[Dürr(20)]} to train the operational expense because

\begin{enumerate}
\def\labelenumi{\alph{enumi}.}
\item
  OpEx is influential in the net income value
\item
  OpEx, in our model, is heavily influenced by sales revenue, which is
  one of the main drivers in our model, up to 23\%
\end{enumerate}

Applying Variational Inference allows us to model the aleatoric and
epistemic uncertainties of OpEx by calculating the posterior
distribution of OpEx's linear regression parameters given the past OpEx,
sales, and inflation data instead of a singular set of parameters (i.e.,
slope and offset) that a simple linear regression gives you. It allows
us to predict the OpEx, and therefore the earnings and other balance
sheet elements, with a confidence interval.

Most of the code from the simple trainable model remains unchanged
except how we model the OpEx during training. The OpEx module is swapped
from \texttt{SimpleOpEx} to \texttt{BayesianOpEx}, and the trajectory
simulator from \texttt{DeterministicSimulator} to
\texttt{MonteCarloSimulator}. The following updates are
made:

\begin{enumerate}
\def\labelenumi{\arabic{enumi}.}
\item
  Bayesian Conversion: Replaced the deterministic \%OR and
  OB\textsubscript{T\_start} variables with Variational Posteriors.
\end{enumerate}

\begin{itemize}
\item
  We learn a Mean ($\mu$) and a Standard Deviation($\sigma$) for each of these
  variables.
\item
  During training/inference, we sample their values from their
  distributions:

  \begin{itemize}
  \item
    Value = $\mu$ + $\sigma$ * $\epsilon$, where $\epsilon$ $\sim$$\mathcal{N}$(0, 1)
  \end{itemize}
\end{itemize}

\begin{enumerate}
\def\labelenumi{\arabic{enumi}.}
\setcounter{enumi}{1}
\item
  Loss Function: Replaced the simple Mean Squared Error (MSE) for OpEx
  with Negative Log Likelihood + KL Divergence.
\end{enumerate}

\begin{itemize}
\item
  This balances fitting the data (Likelihood) with keeping the
  distribution close to our prior beliefs (KL).
\end{itemize}

\begin{enumerate}
\def\labelenumi{\arabic{enumi}.}
\setcounter{enumi}{2}
\item
  Monte Carlo Simulation: Added a run\_monte\_carlo\_forecast function.
  Instead of predicting one single Net Income line, we run the forecast
  1,000 times, sampling different OpEx values each time, to generate
  Confidence Intervals.
\end{enumerate}

So, to answer the `Hint', we are improving the model such that when we
predict the earnings (i.e. net income), we can now model n(t) with the
confidence intervals given by the Monte Carlo simulation. After we have
trained all the variables, they become the very function that take in
this year's y(t) and input x(t) to predict y(t+1). The input x(t)
includes the primary drivers, which are sales and purchases time series,
and inflation, which is the external environment variable that we
include in the model for predicting the operating expense.

We implemented the updated Bayesian model in \texttt{run\_trainable\allowbreak{}\_model\allowbreak{}\_forecast\allowbreak{}\_adv\allowbreak{}\_policies\allowbreak{}\_w\allowbreak{}\_bayesianopex.py},
which composes the \texttt{TrainableFinancialModel} with \texttt{BayesianOpEx},
\texttt{MonteCarloSimulator}, and all advanced trend-based policies
(\texttt{TrendLiquidityPolicy}, \texttt{LintnerDividendPolicy}, etc.)
but with a \texttt{SimpleTax} module (no tax anomalies). The output is reproduced
below:

\noindent\rule{\textwidth}{0.4pt}

\begin{lstlisting}[style=console, caption={Trained policy parameters with Bayesian OpEx}, label=lst:params_simple_bayesian]
Simple parameters:

Final %AG: 0.00000
Final %AM: 0.94058
Final %Depr: 0.05552
Final %AdvPS: 0.02126
Final %AdvPP: 0.07840
Final %AR: 0.15866
Final %AP: 0.30796
Final %Inv (COGS): 0.02939
Total Liquidity (baseline + logit-linear): baseline=0.6012, alpha=-2.3463, beta=-0.411515
  => %TL at t=0: 0.0874, %TL at t=6: 0.0080
Cash % of Liquidity (logit-linear): alpha=-0.2235, beta=0.032375
  => %Cash at t=0: 0.4444, %Cash at t=6: 0.4927
Final %IT: 0.17023
Final %PR: 0.99973
Final DivAdjSpeed: 0.00370
Stock Buyback (baseline + ratio*depr): baseline=1.5612, ratio=-6.674733
Effective ST Debt (baseline + % of sales): baseline=0.0000, pct=0.17869
Cost Ratio (logit-linear): alpha=0.3459, beta=-0.047524
  => CR at t=0: 0.5856, CR at t=6: 0.5152
Bayesian OpEx Variable %: Mean=0.0930, Std=0.0161
Bayesian OpEx Baseline (USD):   Mean=3.93e+10, Std=8.27e+08
OpEx aleatoric uncertainty (USD): 2.45e+09
\end{lstlisting}

\noindent\rule{\textwidth}{0.4pt}

\begin{lstlisting}[style=console, caption={Trained structural parameters with Bayesian OpEx}, label=lst:params_structural_bayesian]
Structural parameters:

Final %AvgSTInt: 0.00003
Final %AvgLTInt: 0.01687
Final %MSReturn: 0.07587
Equity Financing (logit-linear): alpha=0.4917, beta=-0.005617
Avg Maturity Years: 15.3390
\end{lstlisting}

\noindent\rule{\textwidth}{0.4pt}


\begin{table}[H]
\centering
\footnotesize

\begin{tabular}{l rrrrr}
\toprule
\textbf{Line Item} & \textbf{Actual} & \textbf{Pred.\ (Mean)} & \textbf{Lower 2.5\%} & \textbf{Upper 97.5\%} & \textbf{Error \%} \\
\midrule
\multicolumn{6}{l}{\textbf{Assets}} \\
Non-Current Assets        & 211.3 & 211.3 & 211.3 & 211.3 & $+$0.0\% \\
Adv Payments (Purch.)     &  14.6 &  16.3 &  16.3 &  16.3 & $+$12.0\% \\
Accounts Receivable       &  73.0 &  66.0 &  66.0 &  66.0 & $-$9.5\% \\
Inventory                 &   5.7 &   6.2 &   6.2 &   6.2 & $+$7.7\% \\
Cash                      &  35.9 &  31.2 &  31.2 &  31.2 & $-$13.1\% \\
Invest in Mkt Sec         &  18.8 &  31.1 &  31.1 &  31.1 & $+$65.9\% \\
\textbf{Total Assets}     & \textbf{359.2} & \textbf{362.2} & \textbf{362.2} & \textbf{362.2} & $+$0.8\% \\
\addlinespace
\multicolumn{6}{l}{\textbf{Liabilities}} \\
Accounts Payable          &  69.9 &  64.2 &  64.2 &  64.2 & $-$8.2\% \\
Adv Payments (Sales)      &   9.1 &   8.8 &   8.8 &   8.8 & $-$2.3\% \\
Effective ST Debt         &  74.4 &  74.4 &  74.4 &  74.4 & $+$0.0\% \\
Non-Current Liabilities   & 119.9 & 125.4 & 123.6 & 127.2 & $+$4.6\% \\
\textbf{Equity}           & \textbf{73.7} & \textbf{80.7} & \textbf{78.7} & \textbf{82.6} & $+$9.4\% \\
\midrule
\multicolumn{6}{l}{\textbf{Income Statement}} \\
COGS                      & 209.3 & 209.5 & 209.5 & 209.5 & $+$0.1\% \\
OpEx                      &  62.2 &  59.3 &  53.2 &  65.2 & $-$4.6\% \\
Depreciation              &  11.7 &  11.8 &  11.8 &  11.8 & $+$0.6\% \\
Interest Payments         &   0.0 &   2.4 &   2.4 &   2.4 & N/A \\
ST Investment Returns     & $-$0.3 &   2.7 &   2.7 &   2.7 & $+$932.7\% \\
Tax                       &  20.7 &  23.1 &  22.1 &  24.2 & $+$11.7\% \\
\textbf{Net Income}       & \textbf{112.0} & \textbf{112.8} & \textbf{107.9} & \textbf{117.8} & $+$0.7\% \\
\addlinespace
\multicolumn{6}{l}{\textbf{Cash Flow}} \\
Dividends                 &  15.4 &  15.5 &  15.5 &  15.5 & $+$0.6\% \\
Stock Buyback             &  90.7 &  77.6 &  77.6 &  77.6 & $-$14.5\% \\
\bottomrule
\end{tabular}
\par\smallskip
{\scriptsize All figures in \$ billions. Confidence intervals reflect Bayesian OpEx uncertainty only; deterministic line items show identical lower/upper bounds.}
\caption{Bayesian backtest of the advanced-policy model with Bayesian OpEx against Apple's actual FY\,2025 results, including 95\% confidence intervals from Monte Carlo sampling.}
\label{tab:backtest_bayesian}
\end{table}

The Operation Expense learned with Bayesian Variational Inference is as
follows:

\begin{figure}[H]\centering
\includegraphics[width=\textwidth]{report_latex_media/media/image4.png}
\caption{Bayesian variational inference fit of operating expenses to historical data, showing the posterior mean and 95\% credible band.}
\label{fig:bayesian_opex_fit}
\end{figure}

And the net income now has a confidence level due to the noise term, and
we see that the uncertainty increases the further the forecast year as
expected for Variational Inference model's increasing epistemic
uncertainty outside of the training range:

\begin{table}[H]
\centering
\small
\begin{tabular}{l rrr}
\toprule
\textbf{Year} & \textbf{Mean} & \textbf{2.5\% CI} & \textbf{97.5\% CI} \\
\midrule
2025 & 1.13e+11 & 1.08e+11 & 1.18e+11 \\
2026 & 1.23e+11 & 1.18e+11 & 1.28e+11 \\
2027 & 1.34e+11 & 1.28e+11 & 1.39e+11 \\
2028 & 1.45e+11 & 1.39e+11 & 1.51e+11 \\
2029 & 1.56e+11 & 1.50e+11 & 1.62e+11 \\
2030 & 1.68e+11 & 1.61e+11 & 1.75e+11 \\
2031 & 1.80e+11 & 1.73e+11 & 1.88e+11 \\
2032 & 1.93e+11 & 1.86e+11 & 2.00e+11 \\
2033 & 2.05e+11 & 1.98e+11 & 2.13e+11 \\
2034 & 2.19e+11 & 2.10e+11 & 2.27e+11 \\
\bottomrule
\end{tabular}
\par\smallskip
{\footnotesize All figures in USD.}
\caption{Forecasted Apple net income (FY\,2025--2034) with 95\% Bayesian confidence intervals from the Bayesian OpEx model.}
\label{tab:net_income_ci}
\end{table}

\begin{figure}[H]\centering
\includegraphics[width=\textwidth]{report_latex_media/media/advpolicies_bayesian_result_1.png}
\caption{Ten-year Bayesian forecast with 95\% Monte Carlo confidence intervals (1 of 3).}
\label{fig:bayesian_forecast}
\end{figure}

\begin{figure}[H]\ContinuedFloat\centering
\includegraphics[width=\textwidth]{report_latex_media/media/advpolicies_bayesian_result_2.png}
\caption[]{Bayesian forecast (2 of 3).}
\end{figure}

\begin{figure}[H]\ContinuedFloat\centering
\includegraphics[width=\textwidth]{report_latex_media/media/advpolicies_bayesian_result_3.png}
\caption[]{Bayesian forecast (3 of 3).}
\end{figure}


\subsection{Forecast Results}

The forecasted 10 years (with inflation = 3\%, sales growing linearly at
the historical rate of training data) of Balance Sheet, Income
Statement, and Cash Budget is attached below:

\begin{figure}[H]\centering
\includegraphics[width=\textwidth]{report_latex_media/media/image8.png}
\caption{Final ten-year forecast with tax anomaly adjustment, Bayesian OpEx, and advanced policies (1 of 4).}
\label{fig:final_forecast}
\end{figure}

\begin{figure}[H]\ContinuedFloat\centering
\includegraphics[width=\textwidth]{report_latex_media/media/image9.png}
\caption[]{Final forecast (2 of 4).}
\end{figure}

\begin{figure}[H]\ContinuedFloat\centering
\includegraphics[width=\textwidth]{report_latex_media/media/image10.png}
\caption[]{Final forecast (3 of 4).}
\end{figure}

\begin{figure}[H]\ContinuedFloat\centering
\includegraphics[width=\textwidth]{report_latex_media/media/image11.png}
\caption[]{Final forecast (4 of 4).}
\end{figure}


\subsection{Cross-Company Validation}\label{sec:cross_company}

To demonstrate that the model generalizes beyond a single company, we applied
the identical architecture---\texttt{Trainable\-Financial\-Model} with
\texttt{BayesianOpEx}, \texttt{Monte\-Carlo\-Simulator}, and all advanced
policies---to Microsoft\,(MSFT) and Pfizer\,(PFE). Only the trainable
parameters were re-fitted on each company's FY\,2018--2024 financial
statements; no structural changes were made to the model. FY\,2025 actuals
serve as the one-step-ahead backtest in both cases. Full model-fit
visualizations and ten-year projection figures are provided in
Appendix~\ref{app:cross_company}.

\subsubsection{Microsoft (MSFT)}

Microsoft exhibits steady, high-growth financials similar to Apple, making it a
natural validation candidate. Table~\ref{tab:backtest_msft} compares the
one-step-ahead FY\,2025 prediction against actuals.

\begin{table}[H]
\centering
\footnotesize

\begin{tabular}{l rrr}
\toprule
\textbf{Line Item} & \textbf{Actual} & \textbf{Predicted} & \textbf{Error \%} \\
\midrule
\multicolumn{4}{l}{\textbf{Assets}} \\
Non-Current Assets        & 427.9 & 427.3 & $-$0.1\% \\
Adv Payments (Purch.)     &  25.7 &  25.3 & $-$1.6\% \\
Accounts Receivable       &  69.9 &  64.6 & $-$7.6\% \\
Inventory                 &   0.9 &   3.5 & $+$271.9\%$^{*}$ \\
Cash                      &  30.2 &  25.4 & $-$15.9\% \\
Invest in Mkt Sec         &  64.3 &  44.8 & $-$30.3\% \\
\textbf{Total Assets}     & \textbf{619.0} & \textbf{590.9} & $-$4.5\% \\
\addlinespace
\multicolumn{4}{l}{\textbf{Liabilities}} \\
Accounts Payable          &  27.7 &  24.4 & $-$11.9\% \\
Adv Payments (Sales)      &  64.6 &  68.0 & $+$5.3\% \\
Non-Current Liabilities   & 134.3 & 113.7 & $-$15.3\% \\
\textbf{Equity}           & \textbf{343.5} & \textbf{336.7} & $-$2.0\% \\
\midrule
\multicolumn{4}{l}{\textbf{Income Statement}} \\
COGS                      &  65.8 &  67.7 & $+$2.8\% \\
OpEx                      &  65.4 &  66.1 & $+$1.2\% \\
Depreciation              &  22.0 &  25.6 & $+$16.2\% \\
Income Tax                &  21.8 &  20.4 & $-$6.3\% \\
\textbf{Net Income}       & \textbf{101.8} & \textbf{100.5} & $-$1.3\% \\
\addlinespace
\multicolumn{4}{l}{\textbf{Cash Flow}} \\
Dividends                 &  24.1 &  24.2 & $+$0.7\% \\
Stock Buyback             &  18.4 &  33.4 & $+$81.1\% \\
\bottomrule
\end{tabular}
\par\smallskip
{\scriptsize All figures in \$ billions.
$^{*}$\,Inventory error is large in percentage terms but represents only
\$2.6\,B absolute difference on a \$619\,B balance sheet.}
\caption{One-step-ahead backtest of the advanced-policy model with Bayesian OpEx
against Microsoft's actual FY\,2025 results.}
\label{tab:backtest_msft}
\end{table}

The model captures Microsoft's core income-statement dynamics well: net income
is predicted within 1.3\%, COGS within 2.8\%, and operating expenses within
1.2\%. Balance-sheet totals are within 5\%, with equity error at only 2.0\%.
The largest deviations appear in investment in market securities ($-$30.3\%)
and stock buyback ($+$81.1\%), reflecting a shift in Microsoft's
capital-allocation strategy during FY\,2025 that the trend-based policies
cannot fully capture from the training window alone.

\subsubsection{Pfizer (PFE) --- Stress Test}

Pfizer presents a deliberately challenging stress test for the model. Revenue
surged from \$42\,B (FY\,2018) to \$101\,B (FY\,2022) during the COVID-19
vaccine cycle, then contracted to \$63\,B by FY\,2024. This boom--bust pattern
violates the model's assumption of smooth, trend-driven growth.

\begin{table}[H]
\centering
\footnotesize

\begin{tabular}{l rrr}
\toprule
\textbf{Line Item} & \textbf{Actual} & \textbf{Predicted} & \textbf{Error \%} \\
\midrule
\multicolumn{4}{l}{\textbf{Assets}} \\
Non-Current Assets        & 165.3 & 171.9 & $+$4.0\% \\
Adv Payments (Purch.)     &   2.8 &   4.8 & $+$72.5\% \\
Accounts Receivable       &  11.9 &   9.3 & $-$21.5\% \\
Inventory                 &  10.7 &  10.6 & $-$0.1\% \\
Cash                      &   1.1 &   0.6 & $-$43.4\% \\
Invest in Mkt Sec         &   9.2 &   5.8 & $-$36.4\% \\
\textbf{Total Assets}     & \textbf{201.0} & \textbf{203.2} & $+$1.1\% \\
\addlinespace
\multicolumn{4}{l}{\textbf{Liabilities}} \\
Accounts Payable          &   5.2 &   6.7 & $+$27.9\% \\
Adv Payments (Sales)      &   0.8 &   1.7 & $+$118.4\% \\
Non-Current Liabilities   &  84.4 &  91.3 & $+$8.2\% \\
\textbf{Equity}           & \textbf{86.5} & \textbf{80.0} & $-$7.5\% \\
\midrule
\multicolumn{4}{l}{\textbf{Income Statement}} \\
COGS                      &   9.5 &  24.5 & $+$159.2\% \\
OpEx                      &  13.8 &  15.3 & $+$11.2\% \\
Depreciation              &   6.6 &   6.7 & $+$1.9\% \\
Income Tax                & $-$0.3 &   0.0 & N/A \\
\textbf{Net Income}       & \textbf{7.8} & \textbf{0.5} & $-$93.2\% \\
\addlinespace
\multicolumn{4}{l}{\textbf{Cash Flow}} \\
Dividends                 &   9.8 &   9.5 & $-$2.9\% \\
Stock Buyback             &   0.0 &   4.6 & N/A \\
\bottomrule
\end{tabular}
\par\smallskip
{\scriptsize All figures in \$ billions.}
\caption{Stress-test backtest of the advanced-policy model with Bayesian OpEx
against Pfizer's actual FY\,2025 results, illustrating model limitations under
regime-change conditions.}
\label{tab:backtest_pfe}
\end{table}

As expected, the model struggles with Pfizer's post-pandemic normalization.
COGS is overestimated by 159\% because the training window includes the
high-manufacturing COVID years, biasing the learned cost ratio upward. Net
income is consequently underestimated by 93\%. However, balance-sheet aggregates
remain surprisingly robust: total assets are within 1.1\% and equity within
7.5\%, suggesting that the structural accounting constraints---the identity
$A = L + E$ and the debt amortization schedules---provide a stabilizing anchor
even when income-statement parameters are poorly calibrated.

This result highlights an important limitation: the model assumes historical
trends are representative of future dynamics. For companies experiencing regime
changes---such as Pfizer's COVID cycle---the training window must be carefully
curated, or regime-aware extensions (e.g.\ change-point detection, windowed
training) should be incorporated.


\section{Part 2}

\subsection{a) LLM Selection and Infrastructure Setup}

Our favorite LLM is DeepSeek V3.2 for reasoning which has 128k context
window and gpt-5-nano for non-reasoning task which has 400k context
window. From our daily usage, DeepSeekV3.2 follows instructions well with
less hallucinations than other models from competitors. Gpt-5-nano, from
our experience, is an inexpensive model that performs well for pretty
complicated tasks but hallucinates more. We use Microsoft's Azure Foundry
to access these models via Azure API. We have set up an AWS Lambda
function that makes the API calls so that we can hide the API keys for
our Azure API. The AWS Lambda address is saved as an environment
variable in a .env file, with its content having the format of:

{\small AZURE\_DEEPSEEK\_ENDPOINT=\url{https://*****.lambda-url.us-west-1.on.aws/}}

We have included the .env file in the zipped file attached with our
report submission.

\subsection{b) LLM-Only Balance Sheet Forecasting}

We can use the same input data from part 1 and feed it to DeepSeekV3.2
to forecast the balance sheet, as implemented in
\texttt{financial\_forecast/models/llm\_forecaster.py}.

\begin{quote}
We use the following prompt with the identity check that assets =
liabilities + equity:
\end{quote}

\begin{lstlisting}[style=python, caption={Prompt template for LLM-only balance sheet forecasting}, label=lst:prompt_llm_forecast]
prompt: Dict[str, Any] = {
    "task": "Forecast all required financial elements for future years. "
            "Return ONLY valid JSON (no markdown).",
    "historical_facts": historical_rows,
    "value_units": "trillions_of_usd",
    "value_scale_to_usd": LLM_AMOUNT_SCALE,
    "required_elements": ELEMENT_KEYS,
    "accounting_identity": (
        "inventory + nca + accounts_receivable + cash + "
        "investment_in_market_securities + advance_payments_purchases = "
        "accounts_payable + advance_payments_sales + current_liabilities + "
        "non_current_liabilities + equity"
    ),
    "output_schema": {
        "forecast": [
            {
                "inventory": "float",
                "nca": "float",
                "accounts_receivable": "float",
                "cash": "float",
                "investment_in_market_securities": "float",
                "advance_payments_purchases": "float",
                "accounts_payable": "float",
                "advance_payments_sales": "float",
                "current_liabilities": "float",
                "non_current_liabilities": "float",
                "equity": "float",
                "dividends": "float",
                "net_income": "float",
                "sales": "float",
                "cogs": "float",
                "depreciation": "float",
                "opex": "float",
                "tax": "float",
                "stock_buyback": "float",
            }
        ]
    },
}
\end{lstlisting}

The returned fields are further processed within
\texttt{financial\_forecast/models/llm\_forecaster.py} to ensure the balance sheet
identity holds. The forecasted values are plotted below:

\begin{figure}[H]\centering
\includegraphics[width=\textwidth]{report_latex_media/media/image14.png}
\caption{LLM-only balance sheet forecast for Apple (1 of 2).}
\label{fig:llm_only_forecast}
\end{figure}

\begin{figure}[H]\ContinuedFloat\centering
\includegraphics[width=\textwidth]{report_latex_media/media/image15.png}
\caption[]{LLM-only balance sheet forecast (2 of 2). The negative Stockholders' Equity results from enforcing $A = L + E$ on the LLM's inconsistent outputs.}
\end{figure}

We see in ``Stockholders' Equity'' graph that it drops to negative. This
is because the LLM actually failed to return the projected balance sheet
fields that abide by the accounting identity of assets = liabilities +
equity, so the script's procedure that ensures the identity results in
negative equity value. It is clear that LLM-only forecasting of balance
sheet fields is severely inadequate.

\subsection{c) Combining LLM with the Balance Sheet Forecasting Model (Tax Anomalies)}

Is it possible to combine LLM with the balance sheet forecasting model?
Yes. When we look at the forecasted model, we see that one of the major
contributors for predicted income discrepancy from data is the
calculated tax (income underestimated by \textasciitilde5 billion, and
tax overestimated by \textasciitilde2 billion). Net income
(\textasciitilde100 billion dollars) is sales -- cogs -- opex --
depreciation -- debt interest + market securities return -- tax, and we
actually see good predicted values for COGS, opex, and depreciation
(market securities return is derived as $\text{EBT} - \text{EBIT} + \text{InterestExpense}$
per \eqref{eq:ms_return_derivation}, and their combined amount is an order of magnitude smaller
than tax) all show discrepancy from data of less than 1 billion, while
tax, which takes \textasciitilde30\% of net income, is off from the data
by \textasciitilde2 billion, almost 10\% higher than predicted which
lowers the predicted net income.

In the case of Apple, there was a spike in effective tax rate in 2024
due to a one-time income tax charge resulting from the final resolution
of the European Commission State Aid Decision. This skewed the learned
effective tax rate upwards during training, as evidenced by the
consistent over-estimation of tax for one-step forecast model fit.

\begin{figure}[H]\centering
\noindent\begin{minipage}{0.48\textwidth}\centering\includegraphics[width=\linewidth]{report_latex_media/media/tax_withoutanomaly.png}\end{minipage}\hfill\begin{minipage}{0.48\textwidth}\centering\includegraphics[width=\linewidth]{report_latex_media/media/netincome_withoutanomaly.png}\end{minipage}
\caption{One-step-ahead model fit before accounting for tax anomalies: effective tax rate (left) and net income (right). The consistent over-estimation of tax reflects upward bias from the unmodelled one-time charge.}
\label{fig:fit_without_anomaly}
\end{figure}

To avoid this issue, we can ask LLM to notify us of any one-time tax
payments like this so that we can appropriately handle it during
training. When calculating the MSE between predicted and actual tax
owed, the ground truth data can be offsetted by the one-time amount to
better learn the baseline effective tax rate.

We can feed the 10-K to LLM and ask it to identify pages that contain
the one-off tax information, and extract this data in JSON format for
model consumption:

Script: \texttt{run\_tax\_anomaly\_extraction.py} (using \texttt{TaxAnomalyExtractor} from \texttt{financial\_forecast/extraction/})

\begin{lstlisting}[style=python, caption={Prompt for identifying 10-K pages containing tax anomaly disclosures}, label=lst:prompt_tax_selection]
DEFAULT_SELECTION_PROMPT = (
    "You are given pages from a Form 10-K annual report.\n"
    "Identify all pages that contain any of the following:\n"
    "1) Item 7: Management's Discussion and Analysis (MD&A)\n"
    "2) Item 8: Financial Statements and Supplementary Data notes "
    "for:\n"
    " - Income Taxes\n"
    "Include pages with headings, substantive discussion, tables, "
    "and continuations of "
    "these sections.\n"
    'Return ONLY valid JSON with this exact shape: '
    '{{"pages": [<page_number>, ...]}}.\n'
    'If none match, return {{"pages": []}}.\n\n'
    "Pages:\n{pages}"
)
\end{lstlisting}

\begin{lstlisting}[style=python, caption={Prompt for extracting one-time tax anomaly amounts from 10-K pages}, label=lst:prompt_tax_extraction]
DEFAULT_EXTRACTION_PROMPT = (
    "You are an expert financial analyst. Your task is to analyze the provided "
    "excerpts from a company's Form 10-K (MD&A and Financial Footnotes) and "
    "extract data regarding one-time tax anomalies.\n\n"
    "Carefully evaluate the text for the following:\n\n"
    "Current Year Anomalies: Identify any massive, non-recurring, discrete tax "
    "charges or benefits ASSESSED THIS YEAR (not last year or next year) that skewed the current year's net income (e.g., "
    "finalized state aid decisions, sudden impacts from new tax legislation). Positive for tax paid by company, negative for tax reduction or benefit for the company.\n\n"
    "You must respond ONLY with a valid JSON object using the exact schema below. "
    "Do not include any markdown formatting, preamble, or postscript. If a specific "
    "data point is not explicitly mentioned or cannot be reliably quantified from "
    "the text, output null for that field.\n\n"
    "JSON Schema:\n"
    "{{\n"
    '  "current_tax_year": <number>,\n'
    '  "tax_onetime_amount": <number in billions or null>,\n'
    '  "tax_onetime_note": "<string explaining the anomaly or null>",\n'
    "}}\n\n"
    "Pages:\n{pages}"
)
\end{lstlisting}

LLM was able to extract the one-time tax payment from this portion of
10-K:

2024 10-K page 24 and page 39:

\begin{figure}[H]\centering
\includegraphics[width=\textwidth]{report_latex_media/media/image18.png}
\caption{Apple FY\,2024 Form 10-K, page~24: tax anomaly disclosure.}
\label{fig:10k_tax_pages}
\end{figure}

\begin{figure}[H]\ContinuedFloat\centering
\includegraphics[width=\textwidth]{report_latex_media/media/image19.png}
\caption[]{Apple FY\,2024 Form 10-K, page~39: tax anomaly disclosure (continued).}
\end{figure}

JSON output:

\begin{lstlisting}[style=json, caption={LLM-extracted tax anomaly data from Apple's FY\,2024 Form 10-K}, label=lst:tax_anomaly_output]
{
    "selected_pages": [
        23, 24, 25, 26, 27, 41, 42, 43, 45, 46, 47
    ],
    "extraction": {
        "tax_onetime_amount": 10.2,
        "tax_onetime_note": "One-time income tax charge of $10.2 billion net related to the European Commission State Aid Decision."
    }
}
\end{lstlisting}

\begin{sloppypar}
We update our balance sheet forecasting script to include the
tax\_onetime\_amount for better learning of the baseline tax rate. We
run
\texttt{run\_trainable\allowbreak{}\_model\allowbreak{}\_forecast\allowbreak{}\_adv\allowbreak{}\_policies\allowbreak{}\_w\allowbreak{}\_bayesianopex\allowbreak{}\_taxanomaly.py},
which injects a \texttt{TaxWithAnomalies(data.tax\_onetime\_payments)}
module into the model instead of \texttt{SimpleTax()}.
\end{sloppypar}

Effective tax rate drops from 17\% to 15.5\%! And the learning loss is
smaller for the effective tax rate \%IT:

\paragraph{Before}

\begin{figure}[H]\centering
\includegraphics[width=\textwidth]{report_latex_media/media/image20.png}
\caption{Training loss curves before accounting for tax anomalies (1 of 2).}
\label{fig:training_loss_before_tax}
\end{figure}

\begin{figure}[H]\ContinuedFloat\centering
\includegraphics[width=\textwidth]{report_latex_media/media/image21.png}
\caption[]{Training loss before tax anomalies (2 of 2).}
\end{figure}

Before accounting for tax anomalies, the training losses are:
tax loss = 5.45e$-$01, net income loss = 3.37e$-$01.

\paragraph{After}

\begin{figure}[H]\centering
\includegraphics[width=\textwidth]{report_latex_media/media/image22.png}
\caption{Training loss curves after accounting for tax anomalies (1 of 2).}
\label{fig:training_loss_after_tax}
\end{figure}

\begin{figure}[H]\ContinuedFloat\centering
\includegraphics[width=\textwidth]{report_latex_media/media/image23.png}
\caption[]{Training loss after tax anomalies (2 of 2).}
\end{figure}

\begin{figure}[H]\centering
\noindent\begin{minipage}{0.48\textwidth}\centering\includegraphics[width=\linewidth]{report_latex_media/media/tax_withanomaly.png}\end{minipage}\hfill\begin{minipage}{0.48\textwidth}\centering\includegraphics[width=\linewidth]{report_latex_media/media/netincome_withanomaly.png}\end{minipage}
\caption{One-step-ahead model fit after accounting for tax anomalies: effective tax rate (left) and net income (right).}
\label{fig:fit_with_anomaly}
\end{figure}

After accounting for tax anomalies, the training losses drop to:
tax loss = 5.99e$-$02, net income loss = 1.69e$-$01.
The tax loss falls by nearly an order of magnitude (from 0.545 to 0.060),
because the model no longer needs to distort its baseline tax rate to
accommodate one-time charges. The net income loss roughly halves
(from 0.337 to 0.169), since the cleaner tax estimates flow directly
into the income statement.

\subsection{d) Strategic Recommendations for CEO/CFO}

For recommendation to CEO/CFO given our result, we feed the following
prompt, and attach the historical and 10-year forecast of balance sheet
and income statement:

\begin{lstlisting}[style=python, caption={Prompt template for strategic CEO/CFO recommendations from forecast data}, label=lst:prompt_ceo_advisor]
CEO_ADVISOR_PROMPT = (
    "You are an elite Corporate Finance Advisor and Strategic Capital "
    "Expert at a top-tier investment bank. You specialize in advising "
    "the CEOs and CFOs of mature, cash-generative mega-cap technology "
    "firms on capital structure optimization, liquidity management, "
    "and capital allocation.\n\n"

    "Context & Task:\n"
    "I will provide you with two sets of financial data for a target "
    "company:\n"
    "  Historical Baseline: Actual financial performance from recent "
    "years.\n"
    "  ML Forecast: A multi-year financial forecast (Balance Sheet, "
    "Income Statement, and Cash Budget) generated by an advanced "
    "Monte Carlo machine learning ensemble.\n"
    "Your task is to analyze the historical data against the forecast "
    "to identify the three most critical financial leverage points "
    "(risks or opportunities) over the forecasted period.\n\n"

    "STRICT GUARDRAILS (CRITICAL):\n"
    "  NO OPERATIONAL ADVICE: Do not recommend operational, marketing,"
    " or business strategy changes.\n"
    "  STRICTLY FINANCIAL DOMAIN: Recommendations must be confined to "
    "Capital Structure (debt/equity mix, WACC), Liquidity & Treasury "
    "(working capital, commercial paper), and Capital Allocation "
    "(dividends, share repurchases, CapEx).\n"
    "  MANDATORY EVIDENCE: You cannot make a claim without citing "
    "specific numbers from the provided tables.\n\n"

    "Analytical Heuristics (The Mega-Cap Tech Lens):\n"
    "  Historical Benchmarking: Always contextualize the forecast by "
    "comparing it to the historical baseline. Identify structural "
    "breaks or trend accelerations.\n"
    "  Strategic Leverage vs. Distress: Do not penalize shrinking "
    "Equity from aggressive buybacks -- treat as WACC optimization, "
    "not distress, if FCF remains strong.\n"
    "  Net Debt Focus: Evaluate Cash + Marketable Securities against "
    "Liabilities; assess if debt is easily serviced.\n"
    "  Operating Leverage: Note if Revenue/NI grows faster or slower "
    "than historical OpEx.\n\n"

    "Output Format:\n"
    "Provide exactly 3 strategic priorities. For each:\n"
    "  Observation: [Quantitative trend, forecast vs. historical]\n"
    "  Implication: [Financial risk or WACC opportunity]\n"
    "  Recommendation: [Precise corporate finance action]\n\n"

    "Data Input:\n"
    "{historical_tables}\n\n"
    "{forecast_tables}"
)
\end{lstlisting}

The prompt, together with the historical and forecast markdown tables
from the JSON report artifact exported by \texttt{ForecastPipeline}, is
sent to our reasoning model via \texttt{run\_recommendation\_to\_ceo.py}.
This script loads the self-contained \texttt{forecast\_report.json},
constructs the prompt using the \texttt{DeepseekCEOAdvisor} class (which
extends the \texttt{Advisor} ABC from \texttt{financial\_forecast/reporting/advisor.py}),
and calls the Azure DeepSeek endpoint with greedy decoding
(\texttt{temperature=0}, \texttt{top\_k=1}) for minimal hallucination

And here is the recommendation that our reasoning model generated:

\begin{llmoutput}[LLM-Generated Recommendation: Analysis of Financial Leverage Points]

\paragraph{Priority 1: Capital Return Sustainability}\mbox{}\\
\textit{Observation.}~The forecast projects a significant and
accelerating decline in Stockholders'{} Equity, from
\$70.7B in FY2025 to \$40.8B in FY2034, while Net Income grows from
\$111.1B to \$215.4B over the same period. This is a structural break
from the historical baseline (2017--2023), where Equity averaged \$78.5B
and showed no consistent downward trend, even as Net Income grew.

\textit{Implication.}~This trajectory indicates an extremely aggressive
capital return policy is embedded in the forecast, far exceeding
historical levels. While this optimizes WACC by shrinking the equity
base, it creates a critical dependency on continuous, high levels of
free cash flow. The balance sheet is being deliberately levered,
shifting from an equity-heavy to a more debt-dependent structure.

\textit{Recommendation.}~The CFO must formally establish and communicate
a~\textbf{Sustainable Capital Return Framework}. This framework should
define the maximum allowable annual capital return (dividends +
buybacks) as a percentage of forward-looking Free Cash Flow (e.g.,
80--90\%), ensuring it does not jeopardize the credit rating or the
ability to service the forecasted rising debt load. This moves the
policy from implicit to explicit, managing investor expectations and
credit agency concerns.

\noindent\rule{\linewidth}{0.3pt}

\paragraph{Priority 2: Refinancing and Liquidity Risk}\mbox{}\\
\textit{Observation.}~The forecast shows a dramatic increase
in~\textbf{Effective Short-Term Debt}, which balloons from \$87.2B in
FY2025 to \$190.3B by FY2034. Concurrently, highly liquid assets (Cash +
Marketable Securities) are projected to decline from a combined \$62.3B
to \$60.2B. This creates a rising~\textbf{Net Debt}~position.
Historically (2020--2023), the company maintained a robust liquidity
buffer, with Cash + Securities averaging \$55.8B against much lower debt
levels.

\textit{Implication.}~The company is forecasted to become significantly
more reliant on short-term wholesale funding. This introduces
substantial refinancing and interest rate risk, especially as interest
payments are projected to rise from \$4.5B to \$7.5B. The shrinking
liquidity coverage ratio could pressure credit spreads and commercial
paper rates during market stress.

\textit{Recommendation.}~Execute a~\textbf{Strategic Liability
Management Exercise}. The Treasury must term out a significant portion
of the forecasted short-term debt into longer-dated fixed-rate bonds.
This locks in historically low(er) long-term rates for a mega-cap tech
issuer, reduces refinancing risk, and better matches the duration of the
company's long-lived assets. Concurrently, establish a
minimum liquidity covenant (e.g., Cash + Securities $\geq$ 1.0x ST Debt) for
internal risk management.

\noindent\rule{\linewidth}{0.3pt}

\paragraph{Priority 3: Operating Leverage and Capital Allocation}\mbox{}\\
\textit{Observation.}~The forecast predicts a powerful~\textbf{positive
operating leverage}~trend. From FY2025 to FY2034, Revenue grows by
46.5\% (\$416B to \$610B), while the combined growth of COGS and OpEx is
only 39.5\%. This divergence accelerates EBIT margin expansion. This is
a continuation but acceleration of a historical trend; from 2017--2023,
Revenue grew 47.2\% while COGS+OpEx grew 45.8\%.

\textit{Implication.}~This trajectory will generate super-linear growth
in Free Cash Flow, the fundamental driver for all capital allocation
decisions (debt service, buybacks, dividends). It provides immense
balance sheet flexibility and debt capacity. However, it also raises the
strategic opportunity cost of holding excess, low-yielding cash on the
balance sheet.

\textit{Recommendation.}~Initiate a~\textbf{Dynamic Share Repurchase
Program}~tied to Free Cash Flow Surplus. Instead of a fixed buyback
amount, the program should automatically allocate a percentage of
quarterly FCF in excess of a pre-determined threshold (e.g., FCF above
levels needed for debt service and a baseline dividend) to repurchases.
This mechanism directly links the benefit of operating leverage to
per-share earnings accretion and WACC optimization, while maintaining a
cash safety net.

\end{llmoutput}

\subsection{e, f, g) Financial Ratios and Automated Extraction Pipeline}

We need to calculate cost-to-income, quick ratio, debt-to-equity ratio,
debt-to-assets ratio, debt-to-capital ratio, debt-to-EBITDA ratio, and
interest coverage ratio:

\begin{enumerate}
\def\labelenumi{\arabic{enumi}.}
\item
  Cost-to-Income ratio = Operating Expense / Revenue
\end{enumerate}

\begin{quote}
= Describes company's operating efficiency. Operating Expense includes
cost of revenue.

\begin{itemize}
\item
  LLM to extract \textbf{Revenue} and \textbf{Operating Expense} from
  \textbf{Income Statement}. Operating Expense may not be a simple line
  item, may need to combine a few lines, e.g., GM has line for
  ``Automotive and other cost of sales'', ``GM Financial interest,
  operating and other expenses'' and ``Automotive and other selling,
  general and administrative expense''
\end{itemize}
\end{quote}

\begin{enumerate}
\def\labelenumi{\arabic{enumi}.}
\setcounter{enumi}{1}
\item
  Quick Ratio = (Cash \& Cash Equivalents + Short-Term Market Securities
  + Accounts Receivable) / Current Liabilities
\end{enumerate}

\begin{quote}
= Ability to meet short-term liabilities with the most liquid assets

\begin{itemize}
\item
  LLM to extract \textbf{Cash \& Cash Equivalents}, \textbf{Short-Term
  Market Securities}, \textbf{Accounts Receivable}, and \textbf{Current
  Liabilities} from \textbf{Balance Sheet}
\end{itemize}
\end{quote}

\begin{enumerate}
\def\labelenumi{\arabic{enumi}.}
\setcounter{enumi}{2}
\item
  Debt-to-Equity ratio = (Short-Term Debt + Long-Term Debt) / Total
  Equity
\end{enumerate}

\begin{quote}
\begin{itemize}
\item
  LLM to extract \textbf{Total Debt} (including short-term and long-term
  debts) and \textbf{Total Equity} from \textbf{Balance Sheet}
\end{itemize}
\end{quote}

\begin{enumerate}
\def\labelenumi{\arabic{enumi}.}
\setcounter{enumi}{3}
\item
  Debt-to-Assets ratio = (Short-Term Debt + Long-Term Debt) / Total
  Assets
\end{enumerate}

\begin{quote}
\begin{itemize}
\item
  Total Debt already extracted. LLM to extract \textbf{Total Assets}
  from \textbf{Balance Sheet}
\end{itemize}
\end{quote}

\begin{enumerate}
\def\labelenumi{\arabic{enumi}.}
\setcounter{enumi}{4}
\item
  Debt-to-Capital ratio = (Short-Term Debt + Long-Term Debt) /
  (Short-Term Debt + Long-Term Debt + Total Equity)
\end{enumerate}

\begin{quote}
\begin{itemize}
\item
  Total Debt and Total Equity already extracted.
\end{itemize}
\end{quote}

\begin{enumerate}
\def\labelenumi{\arabic{enumi}.}
\setcounter{enumi}{5}
\item
  Debt-to-EBITDA ratio = (Short-Term Debt + Long-Term Debt) / EBITDA
\end{enumerate}

\begin{quote}
= How many years it would take for a company to pay back its debt using
its EBITDA

\begin{itemize}
\item
  Total Debt already extracted. LLM to extract \textbf{Net Income},
  \textbf{Taxes}, and \textbf{Interest Expenses} from \textbf{Income
  Statement}; and \textbf{Depreciation and Amortization} from
  \textbf{Cash Flows} to calculate EBITDA = Net Income + Interest
  Expenses + Taxes + Depreciation + Amortization
\end{itemize}
\end{quote}

\begin{enumerate}
\def\labelenumi{\arabic{enumi}.}
\setcounter{enumi}{6}
\item
  Interest Coverage ratio = EBIT / Interest Expense
\end{enumerate}

\begin{quote}
= How easily a company can pay interest on its outstanding debt

\begin{itemize}
\item
  Interest Expense already extracted. Calculate EBIT = Net Income +
  Interest Expenses + Taxes
\end{itemize}
\end{quote}

\subsubsection{List of Fields to Be Extracted}

\paragraph{Balance Sheet}

\begin{itemize}
\item
  Cash \& Cash Equivalents
\item
  Short-Term Market Securities
\item
  Accounts Receivables
\item
  Total Current Liabilities
\item
  Total Debt (Short-Term and Long-Term)
\item
  Total Equity
\item
  Total Assets
\end{itemize}

\paragraph{Income Statement}

\begin{itemize}
\item
  Revenue
\item
  Operating Expenses
\item
  Net Income
\item
  Taxes
\item
  Interest Expenses
\end{itemize}

\paragraph{Cash Flows}

\begin{itemize}
\item
  Depreciation and Amortization
\end{itemize}

When extracting these values, also give LLM context in which we are
obtaining these values, such that in case there isn't a 1:1 mapping of
these variables to line items in these statements, LLM can assign the
relevant fields and/or combine fields to return the field of interest
(e.g., operating expenses in ExxonMobile include:

\begin{itemize}
\item
  Crude oil and product purchases
\item
  Production and manufacturing expenses
\item
  Selling, general and administrative expenses
\item
  Depreciation and depletion
\item
  Exploration expenses, including dry holes
\item
  Non-service pension and postretirement benefit expense
\item
  Other taxes and duties
\end{itemize}

When extracting the financial statements, also we had to make sure that
we extract any supplementary tables. For example, Google's Income
Statement had a line for `Other Income' which combined interest income,
interest expenses, and others, without listing out individual
components. Only in a separate `Other Income (Expenses), Net' table do
we see the line item `Interest expense' which we need for calculating
Interest Coverage ratio and Debt-to-EBITDA ratio. So, we had to update
the LLM query to also extract supplementary tables.

The extraction step is broken down into three LLM agents for extraction
and a Python script for actual calculation:

\begin{figure}[H]\centering
\includegraphics[width=\textwidth]{report_latex_media/media/image26.jpg}
\caption{End-to-end extraction pipeline: from annual report PDF through LLM-based page identification and table extraction to Python ratio calculation.}
\label{fig:extraction_pipeline}
\end{figure}

Input annual report pdf -\textgreater{} pdfplumber to extract raw text
-\textgreater{} Feed raw text to efficient non-reasoning model to
identify pages that contain financial statements and supplementary
tables of interest -\textgreater{} Feed the relevant pages to the more
resource-intensive reasoning model to extract the text of the primary
table and supplementary tables -\textgreater{} Feed the extracted text
data to the reasoning model to identify all the relevant financial line
items for fields of interest in a rigid json format ready for
consumption by a Python script -\textgreater{} Python script consumes
the json data to calculate the financial ratios

We have tried using embeddings as done in \hyperlink{ref:zhang25}{[Zhang(25)]}, but we
noticed that it is not effective in 2 ways:

\begin{enumerate}
\def\labelenumi{\arabic{enumi}.}
\item
  Vectorization of each page fails to encapsulate all the semantic
  content. For example, if a page contains a supplementary table of
  interest along with other supplementary tables and texts, the
  vectorization of that page may not align closely with our query vector
  for that supplementary table of interest, and we cannot retrieve that
  information.
\item
  Embeddings are more sensitive than LLM to use of different synonyms or
  languages. We have observed from our attempt to use embedding models
  on Alibaba's annual report that our query for `Consolidated Income
  Statement' failed to retrieve the page with ``合併利潤表''
\end{enumerate}

LLMs excel in this area because they are incredibly good at semantic
understanding and most of them are trained on multi-lingual data, giving
them multi-lingual ability as an emergent property. However, LLMs are
susceptible to hallucination. If you try to get it to do too many things
at once, or perform a highly quantitative task, it will hallucinate and
produce wrong answers. It is essential to break the task down into a
step-by-step process that reduces hallucination and response quality,
similar to how chain-of-thought prompting improves output quality
\hyperlink{ref:wei22}{[Wei(22)]}. Also, by first extracting the relevant text before
feeding it to the LLM to extract the financial fields, we prevent
context dilution which is known to degrade output quality \hyperlink{ref:liu23}{[Liu(23)]}.

\subsubsection{run\_statement\_extraction.py}

Using \texttt{FinancialStatementExtractor} from \texttt{financial\_forecast/extraction/}.

\paragraph{Primary financial statement extraction}

\paragraph{Step 1: Page number identifying step}

\begin{lstlisting}[style=python, caption={Step~1 prompt: identifying pages containing the target financial statement}, label=lst:prompt_step1]
prompt = (
    "You are given multiple pages from an annual report. "
    f"Identify which pages contain the table "
    "for the requested query. "
    'Return ONLY JSON in the shape: {"pages": '
    '[<page_number>, ...]}. '
    f"Include all pages that contain the table or line items. "
    'If none, return {"pages": []}.\n\n'
    f"Query: {query}\n\n"
    f"Pages:\n{pages_text}"
)
\end{lstlisting}

Where:

\begin{itemize}
\item
  \hl{query} is one of {[}``Consolidated Balance Sheet'', ``Consolidated
  Income Statement'', ``Consolidated Cash Flow Statement''{]}
\item
  \hl{pages\_text} is the entire text from annual report directly
  extracted with pdfplumber
\end{itemize}

\paragraph{Step 2: Table extraction step}

\begin{lstlisting}[style=python, caption={Step~2 prompt: extracting the primary financial statement table}, label=lst:prompt_step2]
prompt = (
    f"Find the table that corresponds to {query} and output it "
    f"in a nice tabular form from the following pages data.\n"
    f"Pages:\n{pages_text}"
)
\end{lstlisting}

Where:

\begin{itemize}
\item
  \hl{query} is one of {[}``Consolidated Balance Sheet'', ``Consolidated
  Income Statement'', ``Consolidated Cash Flow Statement''{]}
\item
  \hl{pages\_text} is the portion of the pdfplumber-extracted text from
  only the pages identified in \ul{Step 1}
\end{itemize}

Output from Step 2:

\begin{figure}[H]\centering
\includegraphics[width=3.94631in,height=4.18705in]{report_latex_media/media/image27.png}
\caption{Example output from Step~2: the LLM reproduces the Consolidated Balance Sheet in tabular form from the raw page text.}
\label{fig:step2_output}
\end{figure}

\paragraph{Supplementary tables extraction}

\paragraph{Step 3: Page number identifying step}

\begin{lstlisting}[style=python, caption={Step~3 prompt: identifying pages with supplementary disclosures}, label=lst:prompt_step3]
supplementary_query = (
    f"Supplementary disclosures, breakdowns, expansions, "
    f"and note tables for {query}. "
    f"Use this extracted primary statement as context to "
    f"infer relevant line items and "
    f"find related supplementary pages: {primary_context}"
)

prompt = (
    "You are given multiple pages from an annual report. "
    f"Identify which pages contain the table "
    "for the requested query. "
    'Return ONLY JSON in the shape: {"pages": '
    '[<page_number>, ...]}. '
    f"Include all pages that contain the table or line items. "
    'If none, return {"pages": []}.\n\n'
    f"Query: {supplementary_query}\n\n"
    f"Pages:\n{pages_text}"
)
\end{lstlisting}

Where:

\begin{itemize}
\item
  \hl{query} is one of {[}``Consolidated Balance Sheet'', ``Consolidated
  Income Statement'', ``Consolidated Cash Flow Statement''{]}
\item
  \hl{primary\_context} is the extracted table from \ul{Step 2}
\item
  \hl{pages\_text} is the entire text from annual report directly
  extracted with pdfplumber
\end{itemize}

\paragraph{Step 4: Table extraction step}

\begin{lstlisting}[style=python, caption={Step~4 prompt: extracting supplementary tables from identified pages}, label=lst:prompt_step4]
prompt = (
    "You are extracting supplementary financial disclosures "
    "from annual report pages.\n"
    f"Primary statement: {query}\n"
    "Primary statement extraction context:\n"
    f"{primary_context}\n\n"
    "Infer the line items from the primary statement context "
    "above. "
    "Extract supplementary tables related to those inferred "
    "line items "
    "(for example, breakdowns/expansions/schedules/notes such "
    "as an expanded "
    "'Other Income' table). "
    "You can be generous with the supplementary tables you "
    "extract since it is better to have more than not have "
    "necessary information.\n"
    "Return in a nice tabular format.\n"
    'If none are found, return "No supplementary tables '
    'found".\n\n'
    f"Pages:\n{pages_text}"
)
\end{lstlisting}

Where:

\begin{itemize}
\item
  \hl{query} is one of {[}``Consolidated Balance Sheet'', ``Consolidated
  Income Statement'', ``Consolidated Cash Flow Statement''{]}
\item
  \hl{primary\_context} is the extracted table from \ul{Step 2}
\item
  \hl{pages\_text} is the portion of the pdfplumber-extracted text from
  only the pages identified in \ul{Step 3}
\end{itemize}

Output from Step 4:

\begin{figure}[H]\centering
\includegraphics[width=\textwidth]{report_latex_media/media/image28.png}
\caption{Example output from Step~4: supplementary table extractions (1 of 6).}
\label{fig:step4_output}
\end{figure}

\begin{figure}[H]\ContinuedFloat\centering
\includegraphics[width=\textwidth]{report_latex_media/media/image29.png}
\caption[]{Step~4 output (continued, 2 of 6).}
\end{figure}

\begin{figure}[H]\ContinuedFloat\centering
\includegraphics[width=\textwidth]{report_latex_media/media/image30.png}
\caption[]{Step~4 output (continued, 3 of 6).}
\end{figure}

\begin{figure}[H]\ContinuedFloat\centering
\includegraphics[width=\textwidth]{report_latex_media/media/image31.png}
\caption[]{Step~4 output (continued, 4 of 6).}
\end{figure}

\begin{figure}[H]\ContinuedFloat\centering
\includegraphics[width=\textwidth]{report_latex_media/media/image32.png}
\caption[]{Step~4 output (continued, 5 of 6).}
\end{figure}

\begin{figure}[H]\ContinuedFloat\centering
\includegraphics[width=\textwidth]{report_latex_media/media/image33.png}
\caption[]{Step~4 output (continued, 6 of 6).}
\end{figure}

\paragraph{Step 5: Extracting relevant fields as JSON for ratio calculation}

Then, from this final extracted information containing primary and
supplementary tables, we extract the relevant elements from each
financial statement (Balance Sheet, Income Statement, or Cash Flows).
When crafting the prompt, it is important to be careful so that it
prevents LLM from hallucinating as much as possible. LLM is quite bad at
doing math on the spot, but is much better at copying numbers. So, for
retrieving fields that often requires summation of multiple elements, we
ask LLM to copy the elements into an array instead of asking it to sum
them up and return a single value. We use the following prompt:

\begin{lstlisting}[style=python, caption={Step~5 prompt (original): extracting normalized financial fields as JSON}, label=lst:prompt_step5_original]
STATEMENT_CONFIGS: Dict[str, Dict[str, Any]] = {
    "consolidated-balance-sheet": {
        "suffix": ".consolidated-balance-sheet.llm.json",
        "output_suffix": ".consolidated-balance-sheet.normalized.json",
        "fields": [
            "cash_and_cash_equivalents",
            "marketable_securities",
            "accounts_receivable",
            "total_current_liabilities",
            "total_debt_short_term_and_long_term",
            "total_equity",
            "total_assets",
        ],
    },
    "consolidated-income-statement": {
        "suffix": ".consolidated-income-statement.llm.json",
        "output_suffix": ".consolidated-income-statement.normalized.json",
        "fields": [
            "revenue",
            "total_operating_cost",
            "net_income",
            "income_tax_expense",
            "interest_expenses",
        ],
    },
    "consolidated-cash-flow-statement": {
        "suffix": ".consolidated-cash-flow-statement.llm.json",
        "output_suffix": ".consolidated-cash-flow-statement.normalized.json",
        "fields": [
            "depreciation_and_amortization",
        ],
    },
}

fields_schema = ",\n".join(
    [f'  "{field}": number[]'
     for field in required_fields]
)

prompt = (
    "You are a financial statement extraction engine.\n"
    f"Extract normalized {statement_type} values from the "
    f"statement and supplementary tables.\n\n"
    "Return ONLY valid JSON with this exact schema:\n"
    "{\n"
    '  "company_id": "string",\n'
    '  "periods": [\n'
    "    {\n"
    '      "year": "string",\n'
    '      "currency": "string",\n'
    '      "values": {\n'
    f"        {fields_schema}\n"
    "      }\n"
    "    }\n"
    "  ],\n"
    '  "notes": string\n'
    "}\n\n"
    "Rules:\n"
    "1) Use every period/column available in the statement.\n"
    "2) Return numeric array values only (no commas, no currency "
    "symbols, no percent signs).\n"
    "3) If a value cannot be directly mapped to a field, try to "
    "map number(s) to the corresponding field by taking into "
    "account the industry the company is in, and note what you "
    "did in notes field. If you cannot find a match, return an "
    "empty array.\n"
    "4) If multiple numbers make up a field, return them in an "
    "array without adding them up.\n"
    "5) If there are multiple close synonyms, use best "
    "accounting match.\n"
    "6) For parentheses negatives, return negative numbers.\n"
    "7) For year, report the year of the period only.\n"
    "8) For currency, report its formal 3-letter acronym.\n"
    "9) For total_operating_cost, list all elements that should "
    "be included in standard practice for the industry the "
    "company is in. Elements that should be subtracted should "
    "be negative. The sign convention should be such that if an "
    "element is a cost, it should be negative, and if it is an "
    "income, it should be positive.\n"
    "10) Return JSON only.\n\n"
    f"company_id: {company_id}\n\n"
    "statement and supplementary tables:\n"
    f"{statement_and_supplementary_tables}\n"
)
\end{lstlisting}

Where

\begin{itemize}
\item
  \hl{required\_fields} are selected from \hl{STATEMENT\_CONFIGS} and
  are dependent on the type of financial statement we are feeding the
  LLM
\item
  \hl{statement\_type} is one of {[}``Consolidated Balance Sheet'',
  ``Consolidated Income Statement'', ``Consolidated Cash Flow
  Statement''{]}
\item
  \hl{statement\_text} is the extracted primary and supplementary tables
  from Step 4
\end{itemize}

The LLM struggled to identify the correct elements for marketable
securities. So, we add the definition of marketable securities according
to definition:

\begin{itemize}
\item
  short\_term\_market\_securities: these are liquid, unrestricted debt
  or equity investments intended to be sold in the short term, e.g.,
  U.S. Treasuries, commercial paper, money market funds, publicly traded
  equities, short-term investments.
\end{itemize}

The LLM kept wanting to add up the components before reporting them in a
single field, e.g., it kept calculating the total debt by adding up the
short-term and long-term borrowings often with hallucination and
incorrect summed results.

So, we had to explicitly tell LLM to return the original source label
and its value as an array of maps:

\begin{lstlisting}[style=python, caption={Step~5 prompt (improved): requiring source labels to prevent hallucination}, label=lst:prompt_step5_improved]
fields_schema = ",\n".join(
    [
        f'  "{field}": [{{"source_label": number}}]'
        for field in required_fields
    ]
)

prompt = (
    "You are a financial statement extraction engine.\n"
    f"Extract normalized {statement_type} values from the "
    f"statement and supplementary tables.\n\n"
    "Return ONLY valid JSON with this exact schema:\n"
    "{\n"
    '  "company_id": "string",\n'
    '  "periods": [\n'
    "    {\n"
    '      "year": "string",\n'
    '      "currency": "string",\n'
    '      "scale": number,\n'
    '      "values": {\n'
    f"        {fields_schema}\n"
    "      }\n"
    "    }\n"
    "  ],\n"
    '  "notes": string\n'
    "}\n\n"
    "Rules:\n"
    "1) Use every period/column available in the statement.\n"
    "2) Return an array of maps that make up each field where "
    "each map is {source_label: number} (source_label = actual "
    "label from the statement; no commas, no currency symbols, "
    "no percent signs in numbers).\n"
    "3) Ensure that the elements in the array correctly make up "
    "the total value of the field.\n"
    "4) If multiple elements need to be combined to make up a "
    "field, return all of them individually in an array.\n"
    "5) If a value(s) cannot be directly mapped to a field, "
    "try to map element(s) to the corresponding field by "
    "taking into account the industry the company is in, and "
    "note your reasoning in the notes field. If you cannot "
    "find a match, return an empty array.\n"
    "6) If there are multiple close synonyms, use best "
    "accounting match.\n"
    "7) For year, report the year of the period only.\n"
    "8) For currency, report its formal 3-letter acronym.\n"
    "9) For scale, report the scale of the values in the "
    "statement. For example, if the values are in millions, "
    "the scale should be 1E6.\n"
    "10) For total_operating_cost, list all elements that "
    "reflect the total operational costs incurred in "
    "generating income and are included in standard practice "
    "for the industry the company is in.\n"
    "11) Generally, numbers in parentheses are negative.\n"
    "12) For income_tax_expense, interest_expenses, and "
    "total_operating_cost, the sign convention is the "
    "opposite: if an element reduces income, it should be "
    "positive; and if it increases income, it should be "
    "negative.\n"
    "13) For short_term_market_securities, these are liquid, "
    "unrestricted debt or equity investments intended to be "
    "sold in the short term, e.g., U.S. Treasuries, "
    "commercial paper, money market funds, publicly traded "
    "equities, short-term investments.\n"
    "14) Return JSON only.\n\n"
    f"company_id: {company_id}\n\n"
    "statement and supplementary tables:\n"
    f"{statement_and_supplementary_tables}\n"
)
\end{lstlisting}

Output example for GM's 2023 annual report:

\paragraph{Before: LLM attempts to add the total debt by itself}

\begin{lstlisting}[style=json, caption={GM 2023 extraction (original): LLM attempts to sum debt components}, label=lst:extraction_before]
{
    "year": "2023",
    "currency": "USD",
    "values": {
        "cash_and_cash_equivalents": [18853.0],
        "short_term_market_securities": [7613.0],
        "accounts_receivable": [12378.0],
        "total_current_liabilities": [94445.0],
        "total_debt_short_term_and_long_term": [16413.0, 105327.0],
        "total_equity": [68189.0],
        "total_assets": [273064.0]
    }
},
\end{lstlisting}

\paragraph{After: LLM lists out all the original elements without summation}

\begin{lstlisting}[style=json, caption={GM 2023 extraction (improved): source labels prevent hallucinated summation}, label=lst:extraction_after]
{
    "year": "2023",
    "currency": "USD",
    "scale": 1000000,
    "values": {
        "cash_and_cash_equivalents": [
            {"Cash and cash equivalents": 18853.0}
        ],
        "short_term_market_securities": [
            {"Marketable debt securities": 7613.0}
        ],
        "accounts_receivable": [
            {"Accounts and notes receivable, net": 12378.0},
            {"GM Financial receivables, net (current)": 39076.0}
        ],
        "total_current_liabilities": [
            {"Total current liabilities": 94445.0}
        ],
        "total_debt_short_term_and_long_term": [
            {"Short-term debt and current portion of long-term debt - Automotive": 428.0},
            {"Short-term debt and current portion of long-term debt - GM Financial": 38540.0},
            {"Long-term debt - Automotive": 15985.0},
            {"Long-term debt - GM Financial": 66788.0}
        ],
        "total_equity": [
            {"Total Equity": 68189.0}
        ],
        "total_assets": [
            {"Total Assets": 273064.0}
        ]
    }
},
...
\end{lstlisting}

Even with clever prompt engineering, LLM still hallucinates. To measure
its hallucination rate, we repeated Step 5 (i.e., extracting json from
financial statements for financial ratio calculations) for 101 runs and
plotted each field value vs run using the following scripts:

Run multi-run normalization via \texttt{run\_statement\allowbreak{}\_normalization\allowbreak{}\_multi.py},
which invokes the \texttt{StatementNormalizer} and
\texttt{MedianRatioCalculator} from \texttt{financial\_forecast/extraction/}:

\begin{lstlisting}[style=console]
python run_statement_normalization_multi.py
\end{lstlisting}

This performs multiple LLM normalization passes, computes median-aggregated
ratios, and generates hallucination rate analysis via the
\texttt{statement\_hallucination} module:

The results are reproduced below for GM, LVMH, Tencent, Alibaba,
JPMorgan Chase, Exxon Mobil, Volkswagen, Microsoft, and Google:

\begin{figure}[H]\centering
\includegraphics[width=\textwidth]{report_latex_media/media/image34.png}
\caption{Hallucination analysis across 101 LLM extraction runs --- GM (1 of 3).}
\label{fig:hallucination_1}
\end{figure}

\begin{figure}[H]\ContinuedFloat\centering
\includegraphics[width=\textwidth]{report_latex_media/media/image35.png}
\caption[]{Hallucination analysis --- LVMH (2 of 3).}
\end{figure}

\begin{figure}[H]\ContinuedFloat\centering
\includegraphics[width=\textwidth]{report_latex_media/media/image36.png}
\caption[]{Hallucination analysis --- Tencent (3 of 3).}
\end{figure}

\begin{figure}[H]\centering
\includegraphics[width=\textwidth]{report_latex_media/media/image37.png}
\caption{Hallucination analysis --- Alibaba (1 of 3).}
\label{fig:hallucination_2}
\end{figure}

\begin{figure}[H]\ContinuedFloat\centering
\includegraphics[width=\textwidth]{report_latex_media/media/image38.png}
\caption[]{Hallucination analysis --- JPMorgan Chase (2 of 3).}
\end{figure}

\begin{figure}[H]\ContinuedFloat\centering
\includegraphics[width=\textwidth]{report_latex_media/media/image39.png}
\caption[]{Hallucination analysis --- ExxonMobil (3 of 3).}
\end{figure}

\begin{figure}[H]\centering
\includegraphics[width=\textwidth]{report_latex_media/media/image40.png}
\caption{Hallucination analysis --- Volkswagen (1 of 3).}
\label{fig:hallucination_3}
\end{figure}

\begin{figure}[H]\ContinuedFloat\centering
\includegraphics[width=\textwidth]{report_latex_media/media/image41.png}
\caption[]{Hallucination analysis --- Microsoft (2 of 3).}
\end{figure}

\begin{figure}[H]\ContinuedFloat\centering
\includegraphics[width=\textwidth]{report_latex_media/media/image42.png}
\caption[]{Hallucination analysis --- Google (3 of 3).}
\end{figure}

We also calculate the top 15 hallucinated fields as below:

\begin{lstlisting}[style=console, caption={Top 15 most hallucinated financial fields across 101 extraction runs}, label=lst:top_hallucinated_fields]
- exxon_2024 | total_operating_cost | rate_non_null=0.2970 (90/303)
- microsoft_2025 | depreciation_and_amortization | rate_non_null=0.1980 (60/303)
- volkswagen_2024 | net_accounts_receivable | rate_non_null=0.1881 (38/202)
- jpmorgan_2024 | depreciation_and_amortization | rate_non_null=0.1683 (51/303)
- volkswagen_2024 | depreciation_and_amortization | rate_non_null=0.1683 (34/202)
- volkswagen_2024 | total_operating_cost | rate_non_null=0.1287 (26/202)
- alibaba_2025 | cash_and_cash_equivalents | rate_non_null=0.0990 (20/202)
- alibaba_2025 | short_term_market_securities | rate_non_null=0.0990 (20/202)
- alibaba_2025 | net_accounts_receivable | rate_non_null=0.0990 (20/202)
- alibaba_2025 | total_assets | rate_non_null=0.0990 (20/202)
- alibaba_2025 | total_current_liabilities | rate_non_null=0.0990 (20/202)
- alibaba_2025 | total_debt_short_term_and_long_term | rate_non_null=0.0990 (20/202)
- alibaba_2025 | total_equity | rate_non_null=0.0990 (20/202)
- tencent_2024 | short_term_market_securities | rate_non_null=0.0990 (20/202)
- jpmorgan_2024 | total_operating_cost | rate_non_null=0.0792 (24/303)
\end{lstlisting}

As you can see, the hallucination is much better than a coin-toss; LLM
retrieves the correct values more than 70\% of the time. Thus, it is
important to repeat the extraction step multiple times to retrieve the
correct values to do any further quantitative analysis.

The most hallucinated value was ExxonMobile's operating cost because LLM
was not sure whether to include ``Other taxes and duties'' as part of
operating cost. It reasoned more often than not that it is part of the
operating cost. There are several ways to reduce this hallucination
rate. One is to give direct examples of what we want for this specific
item for this company. Another is to reword the name or the definition
of ``Operating Cost'' in a way that there are less ambiguities for LLM
in determining the elements that belong to this field.

As an example of this second point, we were able to reduce the
hallucination rate in extracting the Short-Term Market Securities from
JPMorgan Chase's annual report. Previously, we used the following name
and definition in the LLM prompt:

For marketable\_securities, these are most liquid, unrestricted debt or
equity investments intended to be sold in the near term (e.g., U.S.
Treasuries, commercial paper, money market funds, publicly traded
equities)

And we observed high hallucination rate for this field because LLM
outputted different array of elements that made up marketable
securities, as evidenced by the following 3 extraction runs:

Run 1:

\begin{lstlisting}[style=json, caption={JPMorgan Chase marketable securities extraction --- Run 1}, label=lst:hallucination_run1]
"marketable_securities": [
  {"Federal funds sold and securities purchased under resale agreements": 295001.0},
  {"Securities borrowed": 219546.0},
  {"Available-for-sale securities": 406852.0}
],
\end{lstlisting}

Run 2:

\begin{lstlisting}[style=json, caption={JPMorgan Chase marketable securities extraction --- Run 2}, label=lst:hallucination_run2]
"marketable_securities": [
  {"Federal funds sold and securities purchased under resale agreements": 295001.0},
  {"Securities borrowed": 219546.0},
  {"Trading assets": 637784.0},
  {"Available-for-sale securities": 406852.0}
],
\end{lstlisting}

Run 5:

\begin{lstlisting}[style=json, caption={JPMorgan Chase marketable securities extraction --- Run 5}, label=lst:hallucination_run5]
"marketable_securities": [
  {"Federal funds sold and securities purchased under resale agreements": 295001.0},
  {"Securities borrowed": 219546.0}
],
\end{lstlisting}

Then we updated our naming and definition from
``marketable\_securities'' to ``short\_term\_market\_securities'':

For short\_term\_market\_securities, these are liquid, unrestricted debt
or equity investments intended to be sold in the short term, e.g., U.S.
Treasuries, commercial paper, money market funds, publicly traded
equities, short-term investments.

We see that the hallucination rate practically goes to zero:

\begin{figure}[H]\centering
\noindent\begin{minipage}{0.48\textwidth}\centering\includegraphics[width=\linewidth]{report_latex_media/media/image43.png}\end{minipage}\hfill\begin{minipage}{0.48\textwidth}\centering\includegraphics[width=\linewidth]{report_latex_media/media/image44.png}\end{minipage}
\caption{Hallucination rate for JPMorgan Chase short-term market securities before (left) and after (right) renaming the extraction field. The improved naming eliminates extraction variability.}
\label{fig:hallucination_improvement}
\end{figure}

After obtaining the correct field values, calculating the various
financial ratios is straightforward. We write a Python script for the
calculation, which gives us definite results. The ratio calculation
runs automatically after normalization via
\texttt{run\_ratio\_calculation.py} (using \texttt{RatioCalculator} from
\texttt{financial\_forecast/extraction/statement\_ratios.py})

Calculated fields and financial ratios for GM are as follows for 2022
and 2023 in the 2023 annual report:

\begin{lstlisting}[style=json, caption={Calculated financial ratios for GM (FY\,2022--2023)}, label=lst:ratios_gm]
{
  "year": "2022",
  "currency": "USD",
  "field_values": {
    "cash_and_cash_equivalents": 19153000000.0,
    "short_term_market_securities": 12150000000.0,
    "net_accounts_receivable": 46956000000.0,
    "total_current_liabilities": 91173000000.0,
    "total_debt_short_term_and_long_term": 114699000000.0,
    "total_equity": 71927000000.0,
    "total_assets": 264037000000.0,
    "total_revenue": 156735000000.0,
    "total_operating_cost": 146421000000.0,
    "net_income": 9708000000.0,
    "income_tax_expense": 1888000000.0,
    "interest_expenses": 987000000.0,
    "depreciation_and_amortization": 11290000000.0
  },
  "derived_values": {
    "ebit": 12583000000.0,
    "ebitda": 23873000000.0,
    "quick_assets": 78259000000.0,
    "debt_plus_equity": 186626000000.0
  },
  "ratios": {
    "cost_to_income_ratio": 0.9341946597760551,
    "quick_ratio": 0.8583571890801005,
    "debt_to_equity_ratio": 1.5946584731741904,
    "debt_to_assets_ratio": 0.43440502656824614,
    "debt_to_capital_ratio": 0.6145928220076516,
    "debt_to_ebitda_ratio": 4.804549072173585,
    "interest_coverage_ratio": 12.74873353596758
  }
},
{
  "year": "2023",
  "currency": "USD",
  "field_values": {
    "cash_and_cash_equivalents": 18853000000.0,
    "short_term_market_securities": 7613000000.0,
    "net_accounts_receivable": 51454000000.0,
    "total_current_liabilities": 94445000000.0,
    "total_debt_short_term_and_long_term": 121741000000.0,
    "total_equity": 68189000000.0,
    "total_assets": 273064000000.0,
    "total_revenue": 171842000000.0,
    "total_operating_cost": 162544000000.0,
    "net_income": 9840000000.0,
    "income_tax_expense": 563000000.0,
    "interest_expenses": 911000000.0,
    "depreciation_and_amortization": 11888000000.0
  },
  "derived_values": {
    "ebit": 11314000000.0,
    "ebitda": 23202000000.0,
    "quick_assets": 77920000000.0,
    "debt_plus_equity": 189930000000.0
  },
  "ratios": {
    "cost_to_income_ratio": 0.9458921567486412,
    "quick_ratio": 0.8250304409974059,
    "debt_to_equity_ratio": 1.7853466101570634,
    "debt_to_assets_ratio": 0.44583321126182873,
    "debt_to_capital_ratio": 0.640978255146633,
    "debt_to_ebitda_ratio": 5.247004568571675,
    "interest_coverage_ratio": 12.419319429198683
  }
}
\end{lstlisting}

\subsection{h, i) Applicability to Other Annual Reports}

Our script works for other annual reports as well, including LVMH,
Tencent, Alibaba, JPMorgan Chase, ExxonMobil, Volkswagen, Microsoft, and
Google:

LVMH:

\begin{lstlisting}[style=json, caption={Calculated financial ratios for LVMH (FY\,2022--2024)}, label=lst:ratios_lvmh]
{
  "year": "2022",
  "currency": "EUR",
  "field_values": {
    "cash_and_cash_equivalents": 7301000000.0,
    "short_term_market_securities": 3552000000.0,
    "net_accounts_receivable": 4258000000.0,
    "total_current_liabilities": 31543000000.0,
    "total_debt_short_term_and_long_term": 19739000000.0,
    "total_equity": 56604000000.0,
    "total_assets": 134646000000.0,
    "total_revenue": 79184000000.0,
    "total_operating_cost": 58220000000.0,
    "net_income": 14084000000.0,
    "income_tax_expense": 5362000000.0,
    "interest_expenses": 271000000.0,
    "depreciation_and_amortization": 6226000000.0
  },
  "derived_values": {
    "ebit": 19717000000.0,
    "ebitda": 25943000000.0,
    "quick_assets": 15111000000.0,
    "debt_plus_equity": 76343000000.0
  },
  "ratios": {
    "cost_to_income_ratio": 0.7352495453626995,
    "quick_ratio": 0.4790603303427068,
    "debt_to_equity_ratio": 0.3487209384495795,
    "debt_to_assets_ratio": 0.14659923057498922,
    "debt_to_capital_ratio": 0.2585567766527383,
    "debt_to_ebitda_ratio": 0.7608603476853101,
    "interest_coverage_ratio": 72.75645756457564
  }
},
{
  "year": "2023",
  "currency": "EUR",
  "field_values": {
    "cash_and_cash_equivalents": 7774000000.0,
    "short_term_market_securities": 3490000000.0,
    "net_accounts_receivable": 4728000000.0,
    "total_current_liabilities": 33145000000.0,
    "total_debt_short_term_and_long_term": 21907000000.0,
    "total_equity": 62701000000.0,
    "total_assets": 143694000000.0,
    "total_revenue": 86153000000.0,
    "total_operating_cost": 63600000000.0,
    "net_income": 15174000000.0,
    "income_tax_expense": 5673000000.0,
    "interest_expenses": 760000000.0,
    "depreciation_and_amortization": 7177000000.0
  },
  "derived_values": {
    "ebit": 21607000000.0,
    "ebitda": 28784000000.0,
    "quick_assets": 15992000000.0,
    "debt_plus_equity": 84608000000.0
  },
  "ratios": {
    "cost_to_income_ratio": 0.7382215361043725,
    "quick_ratio": 0.4824860461608086,
    "debt_to_equity_ratio": 0.3493883670116904,
    "debt_to_assets_ratio": 0.1524559132601222,
    "debt_to_capital_ratio": 0.2589235060514372,
    "debt_to_ebitda_ratio": 0.7610825458588104,
    "interest_coverage_ratio": 28.430263157894736
  }
},
{
  "year": "2024",
  "currency": "EUR",
  "field_values": {
    "cash_and_cash_equivalents": 9631000000.0,
    "short_term_market_securities": 3956000000.0,
    "net_accounts_receivable": 4731000000.0,
    "total_current_liabilities": 33696000000.0,
    "total_debt_short_term_and_long_term": 22942000000.0,
    "total_equity": 69287000000.0,
    "total_assets": 149190000000.0,
    "total_revenue": 84683000000.0,
    "total_operating_cost": 65804000000.0,
    "net_income": 12550000000.0,
    "income_tax_expense": 5157000000.0,
    "interest_expenses": 952000000.0,
    "depreciation_and_amortization": 7796000000.0
  },
  "derived_values": {
    "ebit": 18659000000.0,
    "ebitda": 26455000000.0,
    "quick_assets": 18318000000.0,
    "debt_plus_equity": 92229000000.0
  },
  "ratios": {
    "cost_to_income_ratio": 0.7770626926301619,
    "quick_ratio": 0.5436253561253561,
    "debt_to_equity_ratio": 0.3311155050731017,
    "debt_to_assets_ratio": 0.15377706280581807,
    "debt_to_capital_ratio": 0.24875039304340282,
    "debt_to_ebitda_ratio": 0.8672084672084672,
    "interest_coverage_ratio": 19.599789915966387
  }
}
\end{lstlisting}

Tencent:

\begin{lstlisting}[style=json, caption={Calculated financial ratios for Tencent (FY\,2023--2024)}, label=lst:ratios_tencent]
{
  "year": "2023",
  "currency": "CNY",
  "field_values": {
    "cash_and_cash_equivalents": 172320000000.0,
    "short_term_market_securities": 14903000000.0,
    "net_accounts_receivable": 46606000000.0,
    "total_current_liabilities": 352157000000.0,
    "total_debt_short_term_and_long_term": 348618000000.0,
    "total_equity": 873681000000.0,
    "total_assets": 1577246000000.0,
    "total_revenue": 609015000000.0,
    "total_operating_cost": 453642000000.0,
    "net_income": 118048000000.0,
    "income_tax_expense": 43276000000.0,
    "interest_expenses": 12268000000.0,
    "depreciation_and_amortization": 59008000000.0
  },
  "derived_values": {
    "ebit": 173592000000.0,
    "ebitda": 232600000000.0,
    "quick_assets": 233829000000.0,
    "debt_plus_equity": 1222299000000.0
  },
  "ratios": {
    "cost_to_income_ratio": 0.744878204970321,
    "quick_ratio": 0.6639907768410113,
    "debt_to_equity_ratio": 0.3990220686955536,
    "debt_to_assets_ratio": 0.22102956672579926,
    "debt_to_capital_ratio": 0.285214992403659,
    "debt_to_ebitda_ratio": 1.4987876182287188,
    "interest_coverage_ratio": 14.149983697424194
  }
},
{
  "year": "2024",
  "currency": "CNY",
  "field_values": {
    "cash_and_cash_equivalents": 132519000000.0,
    "short_term_market_securities": 12913000000.0,
    "net_accounts_receivable": 48203000000.0,
    "total_current_liabilities": 396909000000.0,
    "total_debt_short_term_and_long_term": 338615000000.0,
    "total_equity": 1053896000000.0,
    "total_assets": 1780995000000.0,
    "total_revenue": 660257000000.0,
    "total_operating_cost": 460160000000.0,
    "net_income": 196467000000.0,
    "income_tax_expense": 45018000000.0,
    "interest_expenses": 11981000000.0,
    "depreciation_and_amortization": 56213000000.0
  },
  "derived_values": {
    "ebit": 253466000000.0,
    "ebitda": 309679000000.0,
    "quick_assets": 193635000000.0,
    "debt_plus_equity": 1392511000000.0
  },
  "ratios": {
    "cost_to_income_ratio": 0.6969407367131284,
    "quick_ratio": 0.4878574181991338,
    "debt_to_equity_ratio": 0.3212983064742631,
    "debt_to_assets_ratio": 0.19012686728486042,
    "debt_to_capital_ratio": 0.24316863565171118,
    "debt_to_ebitda_ratio": 1.093438689740021,
    "interest_coverage_ratio": 21.155663133294382
  }
}
\end{lstlisting}

Alibaba:

\begin{lstlisting}[style=json, caption={Calculated financial ratios for Alibaba (FY\,2024--2025)}, label=lst:ratios_alibaba]
{
  "year": "2024",
  "currency": "CNY",
  "field_values": {
    "cash_and_cash_equivalents": 248125000000.0,
    "short_term_market_securities": 262955000000.0,
    "net_accounts_receivable": 29362000000.0,
    "total_current_liabilities": 421507000000.0,
    "total_debt_short_term_and_long_term": 170776000000.0,
    "total_equity": 1101871000000.0,
    "total_assets": 1764829000000.0,
    "total_revenue": 941168000000.0,
    "total_operating_cost": 827818000000.0,
    "net_income": 71332000000.0,
    "income_tax_expense": 22529000000.0,
    "interest_expenses": 7947000000.0,
    "depreciation_and_amortization": 44504000000.0
  },
  "derived_values": {
    "ebit": 101808000000.0,
    "ebitda": 146312000000.0,
    "quick_assets": 540442000000.0,
    "debt_plus_equity": 1272647000000.0
  },
  "ratios": {
    "cost_to_income_ratio": 0.8795645410808697,
    "quick_ratio": 1.282166132472296,
    "debt_to_equity_ratio": 0.1549872898007117,
    "debt_to_assets_ratio": 0.0967663156033814,
    "debt_to_capital_ratio": 0.13418960638731714,
    "debt_to_ebitda_ratio": 1.1672043304718684,
    "interest_coverage_ratio": 12.810872027180068
  }
},
{
  "year": "2025",
  "currency": "CNY",
  "field_values": {
    "cash_and_cash_equivalents": 145487000000.0,
    "short_term_market_securities": 228826000000.0,
    "net_accounts_receivable": 30652000000.0,
    "total_current_liabilities": 435346000000.0,
    "total_debt_short_term_and_long_term": 230703000000.0,
    "total_equity": 1078393000000.0,
    "total_assets": 1804227000000.0,
    "total_revenue": 996347000000.0,
    "total_operating_cost": 856203000000.0,
    "net_income": 125976000000.0,
    "income_tax_expense": 35445000000.0,
    "interest_expenses": 9596000000.0,
    "depreciation_and_amortization": 42459000000.0
  },
  "derived_values": {
    "ebit": 171017000000.0,
    "ebitda": 213476000000.0,
    "quick_assets": 404965000000.0,
    "debt_plus_equity": 1309096000000.0
  },
  "ratios": {
    "cost_to_income_ratio": 0.8593421769724805,
    "quick_ratio": 0.9302141285322479,
    "debt_to_equity_ratio": 0.21393221209707408,
    "debt_to_assets_ratio": 0.12786805651395305,
    "debt_to_capital_ratio": 0.17623077299143836,
    "debt_to_ebitda_ratio": 1.0806975959826866,
    "interest_coverage_ratio": 17.821696540225094
  }
}
\end{lstlisting}

JPMorgan Chase:

\begin{lstlisting}[style=json, caption={Calculated financial ratios for JPMorgan Chase (FY\,2023--2024)}, label=lst:ratios_jpmorgan]
{
  "year": "2023",
  "currency": "USD",
  "field_values": {
    "cash_and_cash_equivalents": 624151000000.0,
    "short_term_market_securities": 476588000000.0,
    "net_accounts_receivable": 107363000000.0,
    "total_current_liabilities": 3155690000000.0,
    "total_debt_short_term_and_long_term": 436537000000.0,
    "total_equity": 327878000000.0,
    "total_assets": 3875393000000.0,
    "total_revenue": 239425000000.0,
    "total_operating_cost": 177813000000.0,
    "net_income": 49552000000.0,
    "income_tax_expense": 12060000000.0,
    "interest_expenses": 81321000000.0,
    "depreciation_and_amortization": 7512000000.0
  },
  "derived_values": {
    "ebit": 142933000000.0,
    "ebitda": 150445000000.0,
    "quick_assets": 1208102000000.0,
    "debt_plus_equity": 764415000000.0
  },
  "ratios": {
    "cost_to_income_ratio": 0.7426668058891093,
    "quick_ratio": 0.38283291451314927,
    "debt_to_equity_ratio": 1.3314007039203606,
    "debt_to_assets_ratio": 0.11264328546808026,
    "debt_to_capital_ratio": 0.5710733044223361,
    "debt_to_ebitda_ratio": 2.90163847253149,
    "interest_coverage_ratio": 1.7576394781175835
  }
},
{
  "year": "2024",
  "currency": "USD",
  "field_values": {
    "cash_and_cash_equivalents": 469317000000.0,
    "short_term_market_securities": 514547000000.0,
    "net_accounts_receivable": 101223000000.0,
    "total_current_liabilities": 3256638000000.0,
    "total_debt_short_term_and_long_term": 454311000000.0,
    "total_equity": 344758000000.0,
    "total_assets": 4002814000000.0,
    "total_revenue": 278906000000.0,
    "total_operating_cost": 203825000000.0,
    "net_income": 58471000000.0,
    "income_tax_expense": 16610000000.0,
    "interest_expenses": 101350000000.0,
    "depreciation_and_amortization": 7938000000.0
  },
  "derived_values": {
    "ebit": 176431000000.0,
    "ebitda": 184369000000.0,
    "quick_assets": 1085087000000.0,
    "debt_plus_equity": 799069000000.0
  },
  "ratios": {
    "cost_to_income_ratio": 0.7308017755085943,
    "quick_ratio": 0.33319239043455245,
    "debt_to_equity_ratio": 1.3177678255471954,
    "debt_to_assets_ratio": 0.11349790422437815,
    "debt_to_capital_ratio": 0.5685504005286152,
    "debt_to_ebitda_ratio": 2.464139849974779,
    "interest_coverage_ratio": 1.740809077454366
  }
}
\end{lstlisting}

ExxonMobil:

\begin{lstlisting}[style=json, caption={Calculated financial ratios for ExxonMobil (FY\,2023--2024)}, label=lst:ratios_exxonmobil]
{
  "year": "2023",
  "currency": "USD",
  "field_values": {
    "cash_and_cash_equivalents": 31568000000.0,
    "short_term_market_securities": 0.0,
    "net_accounts_receivable": 38015000000.0,
    "total_current_liabilities": 65316000000.0,
    "total_debt_short_term_and_long_term": 41573000000.0,
    "total_equity": 212538000000.0,
    "total_assets": 376317000000.0,
    "total_revenue": 344582000000.0,
    "total_operating_cost": 290950000000.0,
    "net_income": 36010000000.0,
    "income_tax_expense": 15429000000.0,
    "interest_expenses": 849000000.0,
    "depreciation_and_amortization": 20641000000.0
  },
  "derived_values": {
    "ebit": 52288000000.0,
    "ebitda": 72929000000.0,
    "quick_assets": 69583000000.0,
    "debt_plus_equity": 254111000000.0
  },
  "ratios": {
    "cost_to_income_ratio": 0.8443563505928923,
    "quick_ratio": 1.065328556555821,
    "debt_to_equity_ratio": 0.19560266869924436,
    "debt_to_assets_ratio": 0.11047335092488514,
    "debt_to_capital_ratio": 0.163601733100889,
    "debt_to_ebitda_ratio": 0.5700475805235229,
    "interest_coverage_ratio": 61.58775029446407
  }
},
{
  "year": "2024",
  "currency": "USD",
  "field_values": {
    "cash_and_cash_equivalents": 23187000000.0,
    "short_term_market_securities": 0.0,
    "net_accounts_receivable": 43681000000.0,
    "total_current_liabilities": 70307000000.0,
    "total_debt_short_term_and_long_term": 41710000000.0,
    "total_equity": 270606000000.0,
    "total_assets": 453475000000.0,
    "total_revenue": 349585000000.0,
    "total_operating_cost": 299716000000.0,
    "net_income": 33680000000.0,
    "income_tax_expense": 13810000000.0,
    "interest_expenses": 996000000.0,
    "depreciation_and_amortization": 23442000000.0
  },
  "derived_values": {
    "ebit": 48486000000.0,
    "ebitda": 71928000000.0,
    "quick_assets": 66868000000.0,
    "debt_plus_equity": 312316000000.0
  },
  "ratios": {
    "cost_to_income_ratio": 0.8573479983408899,
    "quick_ratio": 0.951085951612215,
    "debt_to_equity_ratio": 0.1541355328411048,
    "debt_to_assets_ratio": 0.09197860962566845,
    "debt_to_capital_ratio": 0.1335506346136605,
    "debt_to_ebitda_ratio": 0.5798854409965521,
    "interest_coverage_ratio": 48.68072289156626
  }
}
\end{lstlisting}

Volkswagen:

\begin{lstlisting}[style=json, caption={Calculated financial ratios for Volkswagen (FY\,2023--2024)}, label=lst:ratios_volkswagen]
{
  "year": "2023",
  "currency": "EUR",
  "field_values": {
    "cash_and_cash_equivalents": 43449000000.0,
    "short_term_market_securities": 26772000000.0,
    "net_accounts_receivable": 88230000000.0,
    "total_current_liabilities": 206036000000.0,
    "total_debt_short_term_and_long_term": 232799000000.0,
    "total_equity": 189186000000.0,
    "total_assets": 600649000000.0,
    "total_revenue": 322284000000.0,
    "total_operating_cost": 314907000000.0,
    "net_income": 17861000000.0,
    "income_tax_expense": 5237000000.0,
    "interest_expenses": 3640000000.0,
    "depreciation_and_amortization": 27566000000.0
  },
  "derived_values": {
    "ebit": 26738000000.0,
    "ebitda": 54304000000.0,
    "quick_assets": 158451000000.0,
    "debt_plus_equity": 421985000000.0
  },
  "ratios": {
    "cost_to_income_ratio": 0.9771102505864393,
    "quick_ratio": 0.7690452153992506,
    "debt_to_equity_ratio": 1.2305297432156714,
    "debt_to_assets_ratio": 0.3875791019380703,
    "debt_to_capital_ratio": 0.5516760074410227,
    "debt_to_ebitda_ratio": 4.286958603417796,
    "interest_coverage_ratio": 7.345604395604395
  }
},
{
  "year": "2024",
  "currency": "EUR",
  "field_values": {
    "cash_and_cash_equivalents": 40296000000.0,
    "short_term_market_securities": 27326000000.0,
    "net_accounts_receivable": 89985000000.0,
    "total_current_liabilities": 217039000000.0,
    "total_debt_short_term_and_long_term": 254081000000.0,
    "total_equity": 196731000000.0,
    "total_assets": 632905000000.0,
    "total_revenue": 324656000000.0,
    "total_operating_cost": 320570000000.0,
    "net_income": 12394000000.0,
    "income_tax_expense": 4411000000.0,
    "interest_expenses": 3446000000.0,
    "depreciation_and_amortization": 31346000000.0
  },
  "derived_values": {
    "ebit": 20251000000.0,
    "ebitda": 51597000000.0,
    "quick_assets": 157607000000.0,
    "debt_plus_equity": 450812000000.0
  },
  "ratios": {
    "cost_to_income_ratio": 0.9874143709033562,
    "quick_ratio": 0.7261690295292551,
    "debt_to_equity_ratio": 1.2915148095622957,
    "debt_to_assets_ratio": 0.40145203466555013,
    "debt_to_capital_ratio": 0.5636074461194467,
    "debt_to_ebitda_ratio": 4.924336686241448,
    "interest_coverage_ratio": 5.876668601276843
  }
}
\end{lstlisting}

Microsoft:

\begin{lstlisting}[style=json, caption={Calculated financial ratios for Microsoft (FY\,2024--2025)}, label=lst:ratios_microsoft]
{
  "year": "2024",
  "currency": "USD",
  "field_values": {
    "cash_and_cash_equivalents": 18315000000.0,
    "short_term_market_securities": 57228000000.0,
    "net_accounts_receivable": 56924000000.0,
    "total_current_liabilities": 125286000000.0,
    "total_debt_short_term_and_long_term": 51630000000.0,
    "total_equity": 268477000000.0,
    "total_assets": 512163000000.0,
    "total_revenue": 245122000000.0,
    "total_operating_cost": 135689000000.0,
    "net_income": 88136000000.0,
    "income_tax_expense": 19651000000.0,
    "interest_expenses": 2935000000.0,
    "depreciation_and_amortization": 22287000000.0
  },
  "derived_values": {
    "ebit": 110722000000.0,
    "ebitda": 133009000000.0,
    "quick_assets": 132467000000.0,
    "debt_plus_equity": 320107000000.0
  },
  "ratios": {
    "cost_to_income_ratio": 0.5535570042672628,
    "quick_ratio": 1.0573168590265472,
    "debt_to_equity_ratio": 0.19230697601656752,
    "debt_to_assets_ratio": 0.10080775065750552,
    "debt_to_capital_ratio": 0.16128981871686654,
    "debt_to_ebitda_ratio": 0.38816922163161893,
    "interest_coverage_ratio": 37.72470187393527
  }
},
{
  "year": "2025",
  "currency": "USD",
  "field_values": {
    "cash_and_cash_equivalents": 30242000000.0,
    "short_term_market_securities": 64323000000.0,
    "net_accounts_receivable": 69905000000.0,
    "total_current_liabilities": 141218000000.0,
    "total_debt_short_term_and_long_term": 43151000000.0,
    "total_equity": 343479000000.0,
    "total_assets": 619003000000.0,
    "total_revenue": 281724000000.0,
    "total_operating_cost": 153196000000.0,
    "net_income": 101832000000.0,
    "income_tax_expense": 21795000000.0,
    "interest_expenses": 2385000000.0,
    "depreciation_and_amortization": 34153000000.0
  },
  "derived_values": {
    "ebit": 126012000000.0,
    "ebitda": 160165000000.0,
    "quick_assets": 164470000000.0,
    "debt_plus_equity": 386630000000.0
  },
  "ratios": {
    "cost_to_income_ratio": 0.5437804375914015,
    "quick_ratio": 1.1646532311744962,
    "debt_to_equity_ratio": 0.12562922332951942,
    "debt_to_assets_ratio": 0.06971048605580264,
    "debt_to_capital_ratio": 0.11160799731008975,
    "debt_to_ebitda_ratio": 0.2694159148378235,
    "interest_coverage_ratio": 52.835220125786165
  }
}
\end{lstlisting}

Google:

\begin{lstlisting}[style=json, caption={Calculated financial ratios for Google (FY\,2023--2024)}, label=lst:ratios_google]
{
  "year": "2023",
  "currency": "USD",
  "field_values": {
    "cash_and_cash_equivalents": 24048000000.0,
    "short_term_market_securities": 86868000000.0,
    "net_accounts_receivable": 47964000000.0,
    "total_current_liabilities": 81814000000.0,
    "total_debt_short_term_and_long_term": 11870000000.0,
    "total_equity": 283379000000.0,
    "total_assets": 402392000000.0,
    "total_revenue": 307394000000.0,
    "total_operating_cost": 223101000000.0,
    "net_income": 73795000000.0,
    "income_tax_expense": 11922000000.0,
    "interest_expenses": 308000000.0,
    "depreciation_and_amortization": 11946000000.0
  },
  "derived_values": {
    "ebit": 86025000000.0,
    "ebitda": 97971000000.0,
    "quick_assets": 158880000000.0,
    "debt_plus_equity": 295249000000.0
  },
  "ratios": {
    "cost_to_income_ratio": 0.7257818955477335,
    "quick_ratio": 1.941965922702716,
    "debt_to_equity_ratio": 0.04188736638918198,
    "debt_to_assets_ratio": 0.02949859838167757,
    "debt_to_capital_ratio": 0.040203353779352344,
    "debt_to_ebitda_ratio": 0.12115830194649437,
    "interest_coverage_ratio": 279.30194805194805
  }
},
{
  "year": "2024",
  "currency": "USD",
  "field_values": {
    "cash_and_cash_equivalents": 23466000000.0,
    "short_term_market_securities": 72191000000.0,
    "net_accounts_receivable": 52340000000.0,
    "total_current_liabilities": 89122000000.0,
    "total_debt_short_term_and_long_term": 10883000000.0,
    "total_equity": 325084000000.0,
    "total_assets": 450256000000.0,
    "total_revenue": 350018000000.0,
    "total_operating_cost": 237628000000.0,
    "net_income": 100118000000.0,
    "income_tax_expense": 19697000000.0,
    "interest_expenses": 268000000.0,
    "depreciation_and_amortization": 15311000000.0
  },
  "derived_values": {
    "ebit": 120083000000.0,
    "ebitda": 135394000000.0,
    "quick_assets": 147997000000.0,
    "debt_plus_equity": 335967000000.0
  },
  "ratios": {
    "cost_to_income_ratio": 0.678902227885423,
    "quick_ratio": 1.6606112968739482,
    "debt_to_equity_ratio": 0.033477501199689924,
    "debt_to_assets_ratio": 0.02417069400518816,
    "debt_to_capital_ratio": 0.03239306241386803,
    "debt_to_ebitda_ratio": 0.08038022364358834,
    "interest_coverage_ratio": 448.07089552238807
  }
}
\end{lstlisting}

\section{Software Architecture}

The codebase follows a modular, object-oriented architecture. All
forecast logic lives in the \texttt{financial\_forecast/} Python package,
organized into sub-packages:

\begin{itemize}
\item \texttt{models/} -- composable financial model with pluggable policy
  modules (\texttt{BaseFinancialModel}, \texttt{TrainableFinancialModel}),
  each module typed against its abstract base class (e.g.,
  \texttt{OpExModule}, \texttt{LiquidityPolicy}, \texttt{DividendPolicy},
  \texttt{DebtPolicy}, \texttt{SimpleTax})
\item \texttt{training/} -- \texttt{ForecastPipeline} orchestrator,
  \texttt{PolicyTrainer}, \texttt{StructuralTrainer}
\item \texttt{inference/} -- \texttt{DeterministicSimulator},
  \texttt{MonteCarloSimulator}, forecast driver models, plotting
\item \texttt{extraction/} -- LLM-based PDF extraction, normalization,
  ratio calculation, hallucination analysis
\item \texttt{reporting/} -- \texttt{TableFormatter} ABC with
  \texttt{MarkdownTableFormatter}, and \texttt{Advisor} ABC with
  \texttt{DeepseekCEOAdvisor} for LLM-based CEO recommendations
\item \texttt{data/} -- \texttt{HistoricalDataLoader} and per-company
  data modules
\item \texttt{serialization/} -- \texttt{.npz} parameter save/load
\item \texttt{clients/} -- \texttt{AzureLLMClient} for Azure-hosted LLM
  endpoints
\end{itemize}

The \texttt{ForecastPipeline} class automatically exports a self-contained
JSON report (\texttt{forecast\_report.json}) containing historical and
forecast markdown tables, model metadata, and a reference to the trained
parameters file. This JSON artifact is consumed by
\texttt{run\_recommendation\_to\_ceo.py} for downstream LLM analysis,
decoupling the training/forecast stage from the LLM recommendation stage.

All method signatures throughout the package use Python type hints,
and each entry-point \texttt{run\_*.py} script is a thin orchestrator
that wires together the modular components.

\section{Testing and Validation}
\label{sec:testing}

This section describes the validation strategy used to verify the
correctness of both the structural financial forecast model (Part~1) and
the LLM-based extraction and analysis pipelines (Part~2). The test suite
is implemented in Python using \texttt{pytest} and comprises 22 test
files organized into a layered architecture of unit, integration, and
end-to-end tests.

\subsection{Balance Sheet Identity as Core Invariant}

The ``no-plug'' design of the forecast model means the fundamental
accounting identity (\emph{Assets = Liabilities + Equity}) is not
enforced by a balancing variable; it must \emph{emerge} from the
correctness of every cash flow, accrual, and valuation equation. This
makes the identity a powerful test oracle: at every forecast step, we
compute

\[
\Delta_t = \lvert \text{Total Assets}_t - (\text{Total Liabilities}_t + \text{Equity}_t) \rvert
\]

and assert $\Delta_t < 10^{-4}$. A nonzero delta pinpoints a specific
structural error (e.g., an expense recorded without a corresponding cash
deduction), rather than being silently absorbed by a plug. This check is
executed across all forecast horizons and all model variants
(deterministic, SimpleOpEx, BayesianOpEx) as the primary integration
test.

\subsection{Test Suite Architecture}

The test suite is organized by the module under test, with each layer
building on the one below it:

\begin{itemize}
  \item \textbf{Type contract tests} (\texttt{test\_types.py}):
    Validate the \texttt{RecurrentState} (14-key) and
    \texttt{ForecastInputs} (3-key) TypedDict schemas, ensuring missing,
    extra, or misspelled keys raise \texttt{KeyError} with descriptive
    context labels.

  \item \textbf{Decomposed unit tests}
    (\texttt{test\_cash\_budget\_decomposed.py},
    \texttt{test\_policy\_trainer\_decomposed.py}): Verify individual
    computations against hand-calculated values using concrete tensors
    (e.g., operating net liquidity balance, loss scale computation under
    STD and NONE modes).

  \item \textbf{Model variant tests}
    (\texttt{test\_simple\_model\_forecast.py},
    \texttt{test\_trainable\_financial\_model\_simpleopex.py},
    \texttt{test\_trainable\_financial\_model\_bayesianopex.py}):
    Each model configuration is tested for balance sheet identity,
    finite output values, determinism (for non-stochastic variants), and
    serialization round-trip correctness.

  \item \textbf{Simulation interface tests}
    (\texttt{test\_simulator\_interface.py}): Verify that
    \texttt{DeterministicSimulator} and \texttt{MonteCarloSimulator}
    expose consistent interfaces and correct default parameters.

  \item \textbf{Regression tests}
    (\texttt{test\_inflation\_alignment.py}): Guard against a previously
    discovered bug where inflation data of different length than
    financial history caused silent misalignment. Tests verify
    tail-alignment, exact-length pass-through, and \texttt{ValueError}
    on insufficient data.

  \item \textbf{Extraction pipeline tests}
    (\texttt{test\_extraction.py}, \texttt{test\_tax\_anomalies.py},
    \texttt{test\_statement\_normalization.py},
    \texttt{test\_llm\_forecast.py}): Test the LLM-based financial
    statement extraction, tax anomaly detection, statement normalization,
    and LLM balance sheet forecasting pipelines. These use
    \texttt{unittest.mock} to isolate tests from external API calls.

  \item \textbf{Risk warning pipeline tests}
    (\texttt{tests/risk/}---8 files): End-to-end integration tests for
    the risk extraction pipeline, including PDF page flagging, category
    extraction, risk memo synthesis, entity anonymization, and Pydantic
    model validation for domain types.
\end{itemize}

\subsection{Testing Strategies}

The following strategies are applied across the test suite:

\begin{itemize}
  \item \textbf{Balance sheet identity invariant}: Every model test
    asserts $\lvert \text{Assets} - \text{Claims} \rvert < 10^{-4}$ at
    each forecast step, across all time horizons and model variants.

  \item \textbf{Finite-value enforcement}: All TensorFlow tensor outputs
    are checked with \texttt{tf.debugging.assert\_all\_finite} or
    equivalent assertions to detect NaN or Inf propagation.

  \item \textbf{Determinism checks}: Non-stochastic model variants are
    called multiple times with identical inputs; outputs must be
    bit-identical across runs.

  \item \textbf{Stochastic correctness}: The BayesianOpEx variant is
    tested for variational sampling sanity, bijector-implied economic
    bounds (ensuring sampled parameters remain in economically meaningful
    ranges), and Monte Carlo trajectory finiteness across all samples.

  \item \textbf{Numerical precision}: All financial computations use
    \texttt{tf.float64} (enforced by dtype assertions in tests).
    Floating-point comparisons use \texttt{pytest.approx()} with
    explicit tolerances. Loss scale computations handle degenerate cases
    with $\epsilon = 10^{-12}$.

  \item \textbf{Mock isolation}: LLM and external API interactions are
    replaced with \texttt{unittest.mock} patches, ensuring tests run
    deterministically without network access. Control flow, JSON output
    structure, and file naming conventions are verified against mocked
    responses.

  \item \textbf{Serialization round-trips}: Trainable model parameters
    are saved and reloaded, with assertions that restored values match
    originals within floating-point tolerance.

  \item \textbf{Pydantic validation}: Risk domain models are tested for
    type coercion (e.g., string-to-int page numbers), optional field
    handling, union type acceptance, and rejection of invalid inputs.
\end{itemize}

\subsection{Parameter Recovery (Identifiability Testing)}

The preceding tests verify \emph{structural} correctness: that the
accounting identity holds, outputs are finite, and parameters survive
serialization.  However, they do not verify that the gradient-based
training pipeline converges to the \emph{correct} parameter values.  A
model can produce finite, balanced forecasts while its learned
parameters sit at a spurious local minimum far from the true
data-generating process.

\textbf{Parameter recovery testing} (also known as
\emph{simulation-based calibration}) closes this gap.  The procedure
is:
\begin{enumerate}
  \item Define a set of known ground-truth parameter values.
  \item Create a ``generator'' model and assign those values.
  \item Run \texttt{forecast\_step} in a loop from a cash-poor initial
    balance sheet, producing perfectly consistent synthetic financial
    statements, including the full cash budget circularity (new ST/LT
    borrowing, equity financing, interest feedback).
  \item Train a fresh ``recovery'' model on the synthetic data.
  \item Assert that recovered parameters match ground truth within
    tolerance.
\end{enumerate}

We test four scenarios of increasing complexity: (1)~all simple
(static-ratio) policies, (2)~advanced policies with logit-linear
trends and Lintner dividends, (3)~advanced policies with Bayesian
variational OpEx, and (4)~advanced policies with Bayesian OpEx and
one-time tax anomalies.  The most involved scenario~(4) uses the
following OpEx and tax formulations:

\begin{equation}\label{eq:opex_bayesian_recovery}
\text{OpEx}_t = \underbrace{\beta_0 \cdot \Pi_t^{\text{cum}}}_{\text{baseline}} + \underbrace{\beta_1 \cdot (\text{Sales}_t - \bar{S})}_{\text{variable}} + \varepsilon_t, \qquad \varepsilon_t \sim \mathcal{N}(0,\,\sigma^2)
\end{equation}

\begin{equation}\label{eq:tax_anomaly_recovery}
\text{Tax}_t = \text{EBT}_t \cdot \tau + \delta_t, \qquad \delta_t = \begin{cases} a_{\text{onetime}} & \text{if year } t \text{ has anomaly} \\ 0 & \text{otherwise} \end{cases}
\end{equation}

where $\Pi_t^{\text{cum}} = \prod_{s=1}^{t}(1 + \text{Inflation}_s)$ is
the cumulative inflation factor (so that $\beta_0 \cdot \Pi_t^{\text{cum}}$
is the closed-form equivalent of the recursive baseline
$\text{OB}_{t-1} \cdot (1 + \text{Inflation}_t)$ from \eqref{eq:opex}),
$\bar{S}$ is the mean historical sales (used to center the variable
component), $\beta_0, \beta_1$ have Normal variational posteriors with
trainable location and scale, $\sigma$ is the learned aleatoric noise,
$\tau$ is the effective tax rate, and $\delta_t$ are known one-time
payments (e.g., a \$2B settlement in 2020).

\begin{lstlisting}[style=console, caption={Parameter recovery results for Scenario~4: all parameters recovered within tolerance}, label=lst:param_recovery_s4]
Parameter Recovery — Scenario 4 (BayesianOpEx + Tax Anomalies)
---------------------------------------------------------------
Parameter                     True     Recovered   Status
---------------------------------------------------------------
Depreciation Rate             0.0550   ~0.055      PASS (rel 10%)
Asset Growth                  0.0080   ~0.008      PASS (rel 15%)
Accounts Receivable %         0.1600   ~0.160      PASS (rel 10%)
Cost Ratio Alpha              0.3200   ~0.320      PASS (abs 0.10)
Cost Ratio Beta              -0.0200   ~-0.020     PASS (abs 0.02)
BayesianOpEx Var. Loc         0.1000   ~0.100      PASS (rel 20%)
BayesianOpEx Base Scale       1.0(i)   < 1.0       PASS (shrinks)
BayesianOpEx Var. Scale       1.0(i)   < 1.0       PASS (shrinks)
Noise Sigma                   0.0050   < 0.1       PASS
Income Tax Rate               0.1500   ~0.150      PASS (rel 10%)
Lintner Payout Ratio          0.1600   ~0.160      PASS (rel 10%)
Lintner Adj. Speed            0.3000   ~0.300      PASS (rel 20%)
Avg LT Interest Rate          0.0350   ~0.035      PASS (rel 20%)
Avg Maturity Years            5.0000   ~5.000      PASS (rel 20%)
---------------------------------------------------------------
All 82 tests pass across 4 scenarios (simple, advanced, Bayesian,
Bayesian + tax anomalies). (i) = initial value before training.
\end{lstlisting}

The results confirm three properties: (a)~the model is
\emph{identifiable}: ground-truth parameters are uniquely recoverable
from data generated by those same parameters; (b)~the Adam optimizer
converges to the correct minimum rather than a spurious local optimum;
and (c)~the no-plug accounting framework correctly propagates all
financial flows through the cash budget circularity, even under
stochastic OpEx and one-time tax anomalies.  Parameters trained
indirectly via \texttt{forecast\_step} transitions (interest rates,
debt maturity, equity financing mix) require wider tolerances due to
the smaller effective sample size of state transitions, consistent with
the theoretical expectation that indirect estimation paths have lower
Fisher information.

Beyond the numerical tolerance checks above, we also verify the fit
\emph{visually} by taking the trained recovery model and applying its
own one-step-ahead prediction (\texttt{forecast\_step}) to every
historical transition, then overlaying the resulting fitted values on
the synthetic observations.  This is the same one-step-ahead diagnostic
that \texttt{ForecastPipeline} uses for real historical data.
Figure~\ref{fig:param_recovery_fit_s4} shows the result for
Scenario~4 (the most involved scenario, with Bayesian OpEx and tax
anomalies).  Across twelve key metrics---including the net income,
cost of goods sold, operating expenses (with injected noise),
tax (with the 2020 anomaly spike and 2022 dip), non-current assets,
cash, working-capital components, and the full capital structure
(effective short-term debt, non-current liabilities, equity)---the
red fitted series tracks the black synthetic truth essentially
perfectly.  This confirms that the recovered parameters reproduce the
data-generating process not only in isolation (passing the per-parameter
tolerance tests) but also jointly through the full forward model
evolution, including the cash budget circularity.

\begin{figure}[H]
\centering
\includegraphics[width=\textwidth]{report_latex_media/media/parameter_recovery_s4_fit.png}
\caption{Parameter recovery, Scenario~4 (advanced policies + Bayesian
OpEx + tax anomalies): \emph{deterministic} one-step-ahead fit of the
trained recovery model against the synthetic ground truth it was
trained on.  Black markers are the synthetic observations produced by
the generator model; red markers are the trained model's
\texttt{forecast\_step} predictions with \texttt{use\_mean\_opex=True}
(posterior mean, no sampling), initialized from each preceding state.
The 2020 spike and 2022 dip in the Tax panel correspond to the known
one-time anomalies $\delta_{2020}=\$2\text{B}$ and
$\delta_{2022}=-\$0.5\text{B}$, which the \texttt{TaxWithAnomalies}
module reproduces exactly.  The slight noise in the OpEx panel reflects
the injected aleatoric noise $\varepsilon_t\sim\mathcal{N}(0,\sigma^2)$.}
\label{fig:param_recovery_fit_s4}
\end{figure}

The deterministic fit in Figure~\ref{fig:param_recovery_fit_s4} uses
the posterior \emph{mean} of the Bayesian OpEx variational parameters
and suppresses aleatoric noise.  To also visualize the \emph{predictive
uncertainty} implied by the trained posterior, we repeat the same
one-step-ahead fit but instead draw $N=500$ Monte Carlo samples from
the OpEx posterior and noise distributions at each transition, then
plot the sample mean together with a 95\% credible band.  The only
code difference from the deterministic variant is passing
\texttt{use\_mean\_opex=False} to \texttt{forecast\_step\_compiled} on
a state tiled to $[N, 14]$ after a single
\texttt{opex\_module.prepare\_mc} call.  Figure~\ref{fig:param_recovery_fit_s4_mc}
shows the result.  As expected, the band is widest on OpEx itself and
propagates into Net Income, Tax, Non-Current Liabilities, and Equity
(which depend on OpEx through the income statement and cash budget),
while COGS, cash, working-capital items, and non-current assets---which
are deterministic functions of sales, cost ratio, and working-capital
ratios at the 1-step horizon---collapse to a near-zero band.  This is
the expected topology of uncertainty in the model and further confirms
the recovered posterior behaves correctly.

\begin{figure}[H]
\centering
\includegraphics[width=\textwidth]{report_latex_media/media/parameter_recovery_s4_fit_mc.png}
\caption{Parameter recovery, Scenario~4: \emph{Monte Carlo}
one-step-ahead fit.  Same setup as Figure~\ref{fig:param_recovery_fit_s4}
except the trained model is evaluated with
\texttt{use\_mean\_opex=False} on a state tiled to $N=500$ samples,
drawing from the Bayesian OpEx posterior
$q(\beta_0),q(\beta_1)$ and the aleatoric noise
$\mathcal{N}(0,\sigma^2)$.  The red line is the MC mean and the shaded
red region is the 95\% credible band (2.5--97.5 percentiles across
samples).  Panels whose value depends on OpEx at this 1-step horizon
(OpEx, Net Income, Tax, Non-Current Liabilities, Equity) display the
propagated uncertainty; panels that do not (COGS, NCA, Cash, AR, AP,
Inventory, Effective ST Debt) have near-zero bands, as they are
deterministic functions of the (exact) state and exogenous sales.}
\label{fig:param_recovery_fit_s4_mc}
\end{figure}

\section{References}

\hypertarget{ref:almeida14}{}\hyperlink{ref:almeida14}{[Almeida(14)]} Almeida, Heitor et al., (2014), ``Corporate Liquidity
Management: A Conceptual Framework and Survey'', Available at SSRN:
\url{http://ssrn.com/abstract_id=2336367}

\hypertarget{ref:brav05}{}\hyperlink{ref:brav05}{[Brav(05)]} Brav, Alon, Graham, John~R., Harvey, Campbell~R., \&
Michaely, Roni (2005). ``Payout Policy in the 21st Century,''
\emph{Journal of Financial Economics}, 77(3), 483--527.

\hypertarget{ref:durr20}{}\hyperlink{ref:durr20}{[Dürr(20)]} Dürr, O., Sick, B., \& Murina, E. (2020). Probabilistic
deep learning: With Python, Keras and TensorFlow probability. Manning.

\hypertarget{ref:iasb21}{}\hyperlink{ref:iasb21}{[IASB(21)]} International Accounting Standards Board (IASB). (2021).
IAS 2: Inventories.
\url{https://www.ifrs.org/content/dam/ifrs/publications/pdf-standards/english/2024/issued/part-a/ias-2-inventories.pdf?bypass=on}

\hypertarget{ref:koller20}{}\hyperlink{ref:koller20}{[Koller(20)]} Koller, T., Goedhart, M., \& Wessels, D. (2020).
\emph{Valuation: Measuring and Managing the Value of Companies} (7th ed.). McKinsey \& Company.

\hypertarget{ref:lintner56}{}\hyperlink{ref:lintner56}{[Lintner(56)]} Lintner, John (1956). ``Distribution of Incomes of
Corporations Among Dividends, Retained Earnings, and Taxes,''
\emph{American Economic Review}, 46(2), 97--113.

\hypertarget{ref:liu23}{}\hyperlink{ref:liu23}{[Liu(23)]} Liu, Nelson F. et al, (2023), ``Lost in the Middle: How
Language Models Use Long Contexts'',
\url{https://arxiv.org/abs/2307.03172}

\hypertarget{ref:merton74}{}\hyperlink{ref:merton74}{[Merton(74)]} Merton, Robert C. (1974). ``On the Pricing of
Corporate Debt: The Risk Structure of Interest Rates,''
\emph{Journal of Finance}, 29(2), 449--470.

\hypertarget{ref:myers84}{}\hyperlink{ref:myers84}{[Myers(84)]} Myers, Stewart~C., \& Majluf, Nicholas~S. (1984).
``Corporate Financing and Investment Decisions When Firms Have
Information That Investors Do Not Have,'' \emph{Journal of Financial
Economics}, 13(2), 187--221.

\hypertarget{ref:pareja07}{}\hyperlink{ref:pareja07}{[Pareja(07)]} Vélez Pareja, Ignacio, (2007), ``Forecasting Financial
Statement with No Plugs and No Circularity'', Available at SSRN:
\url{http://ssrn.com/abstract=1031735}

\hypertarget{ref:pareja09}{}\hyperlink{ref:pareja09}{[Pareja(09)]} Vélez Pareja, Ignacio, (2009), ``Constructing
Consistent Financial Planning Models for Valuation'' (August 15),
Available at SSRN: \url{http://ssrn.com/abstract=1455304}

\hypertarget{ref:wei22}{}\hyperlink{ref:wei22}{[Wei(22)]} Wei, Jason et al., (2022), ``Chain-of-Thought Prompting
Elicits Reasoning in Large Language Models'',
\url{https://arxiv.org/abs/2201.11903}

\hypertarget{ref:zhang25}{}\hyperlink{ref:zhang25}{[Zhang(25)]} Zhang, Haichao et al., (2025), ``Research on Financial
Statement Checking Relationship Recognition System Based on Large
Language Models'',
\url{https://dl.acm.org/doi/pdf/10.1145/3745238.3745291}

\appendix

\section{Reviewer's Feedback}

\subsection{Feedback on Part 1}

\emph{(1) You identified the principle of ``no plug'' and some
accounting identities, such as non-current assets (NCA) being a function
of depreciation and capex:}

\emph{NCA\_t = NCA\_t-1 -- Depreciation\_t + CapEx\_t}

\emph{However, later, while formulating the time-series model, you
consider the non-current asset at time t to be driven by a ratio:}

\emph{NCA\_t = NCA\_t-1 * (1 + \%AG)}

\emph{Is it possible to design a workflow that fully respects the
accounting identities (like the first equation) and can be driven by
drivers?}

\textbf{Response:} Thank you for catching this mistake. This was a typo in
the report; the actual script had the correct relationship of CapEx\_t =
\%AG * Sales\_t . We fixed the typo in the report

\emph{(2) What is the difference between the FIFO and LIFO principles in
accounting? Under what circumstances will these two accounting choices
make a difference? Is it possible to avoid using COGS as a plug?}

\textbf{Response:} FIFO (First-In, First-Out): Under FIFO, the oldest
inventory costs are matched against current sales revenues.

LIFO (Last-In, First-Out): Under LIFO, the most recent inventory costs
are matched against current sales revenues.

The most difference can arise from these two accounting standards when
there is a rapid inflationary (deflationary) environment. Under FIFO,
the cost ratio will be smoother and slightly lower (higher) than the
actual current cost of replacement. Under LIFO, the cost ratio will be
more volatile since the recent inventory cost is used, and it will
tightly track the inflation (deflation) and will squeeze (grow) the
margin immediately. Under inflation = 0\%, FIFO and LIFO should equal
each other because old cost = new cost. Under high inflation, FIFO
profit \textgreater{} LIFO profit.

In our model, we avoid using COGS as a plug by updating it so that we
learn the cost ratio (CR\textsubscript{t}) directly from the historical
COGS and sales data by modeling cost ratio as a logit and learning the
underlying parameters $\alpha$ and $\beta$ to ensure positivity :

$\text{logit}(\text{CR}_t) = \ln(\text{CR}_t / (1 - \text{CR}_t)) = \alpha + \beta \cdot t$ (from \eqref{eq:cost_ratio_logit})

$\text{CR}_t = \text{COGS}_t / \text{Sales}_t$
(from \eqref{eq:cost_ratio})

(3) Similar to (1), can you list out all the entities (rows) in the
income statement, cash budget, and balance sheet (simplified versions
are fine)? Then analyze how these three statements connect with each
other. You may start with revenue on the income statement.

Response: We now include the 10-year forecasted Balance Sheet, Income
Statement, and Cash Budget in the \emph{Projected Balance Sheet, Income
Statement, and Cash Budget} subsection.

\emph{(4) Advance payment is defined as: AdvPP\_t = \%AdvPP *
Purchase\_t+1}

\emph{Is it possible to remove the future dependency on this variable?
It may break causal timing and introduce leakage during training. Are
there alternative approaches to model this term?}

\textbf{Response:} We updated our model such that the advance payments for
purchases (AdvPP) are proportional to this year's purchases instead of
next year, with the justification:

\begin{itemize}
\item
  From \emph{Advance Payments for Purchases}: ``Advance payments to
  suppliers for purchases (\$AdvPP\$) can be modeled as strictly
  proportional to the total purchases made during the current year. This
  approach relies on the simplifying assumption that current purchasing
  volumes serve as a reliable proxy for near-term future demand.''
\end{itemize}

Same is applied for advance payments from customers for sales (AdvPS):

\begin{itemize}
\item
  From \emph{Advance Payments for Future Sales}: ``Advance payments
  received from customers for future sales (\$AdvPC\$) are modeled as a
  direct proportion of the current year's total sales.
  This assumes that the volume of upfront customer payments scales
  linearly with overall revenue generation.''
\end{itemize}

\emph{(5) For CapEx, it's great that you split this term into
maintenance capex and growth capex:}

\emph{CapEx\_t = MaintenanceCapex\_t + GrowthCapex\_t}

\emph{You further assume maintenance capex is tied to depreciation. This
is a common approach when forecasting future cash flows by assuming a
firm is in a steady state. What if the firm is not in a steady state?}

\textbf{Response:} We updated our model to handle the case where the asset
replacement rate is not unitary (= 1.0). We reflect that change in the
report as well:

\begin{itemize}
\item
  From \emph{Non-Current Assets (NCA):} ``To model CapEx, we diverge
  from a rigid steady-state assumption where maintenance perfectly
  offsets depreciation. Instead, we decompose CapEx into a maintenance
  component (scaled by a trainable parameter, \%AM, to reflect partial,
  full, or inflationary replacement) and a growth component driven by
  sales expansion:
\end{itemize}

$\text{CapEx}_t = \text{\%AM} \cdot \text{Depreciation}_t + \text{\%AG} \cdot \text{Sales}_t$ \eqref{eq:capex}''

\emph{(6) In your report, the financing dynamics on debt are as
follows:}

\emph{Debt\_t-1 =\textgreater{} Interest\_t, PrincipalDue\_t
=\textgreater{} deficit\_t =\textgreater{} Debt\_t}

\emph{Can you plot the amount of debt and see whether the outcome is
sensible?}

\begin{figure}[H]\centering
\noindent\begin{minipage}{0.48\textwidth}\centering\includegraphics[width=\linewidth]{report_latex_media/media/image12.png}\end{minipage}\hfill\begin{minipage}{0.48\textwidth}\centering\includegraphics[width=\linewidth]{report_latex_media/media/image13.png}\end{minipage}
\caption{Forecasted debt levels for Apple: short-term effective debt (left) and non-current liabilities (right), demonstrating that the model avoids explosive debt spirals.}
\label{fig:debt_levels}
\end{figure}

\textbf{Response:} We have plotted the forecasted debt levels (see above).
The outcome is sensible and mathematically stable. We verified that the
model does not fall into an explosive `debt
spiral' The financing dynamics converge realistically.

Specifically for Apple, the model accurately captures two key trends:

\begin{enumerate}
\def\labelenumi{\arabic{enumi}.}
\item
  Short-Term Debt: The model predicts short-term debt scaling
  proportionally with Sales. This correctly reflects
  Apple's real-world strategy of utilizing revolving
  credit and commercial paper for working capital needs, despite
  possessing a large cash balance.
\item
  Long-Term Debt: The forecast accurately continues the historical trend
  of Apple slowly paying down its non-current liabilities.Ultimately,
  the plot confirms that the endogenous deficit calculation handles debt
  rollover gracefully without compounding errors.
\end{enumerate}

\emph{(7) Some rules seem a bit ad hoc. Is it possible to tie this to
terms with a stronger financial relationship?}

\textbf{Response:} We can redefine our training variables from arbitrary
percentages to these time-based efficiency metrics.

Receivables: Instead of \%AR, we can train a target DSO (Days Sale
Outstanding) parameter.

AR\textsubscript{t} = DSO / 365 * Sales\textsubscript{t}

Inventory: Inventory is a function of the cost to produce goods, not the
final sales price. Tie it to COGS using DSI (Days Sales of Inventory)

Inv\textsubscript{t} = DSI / 365 * COGS\textsubscript{t}

Payables: We can use DPO (Days Purchases Outstanding) instead of \%AP:

AP\textsubscript{t} = DPO / 365 * Purchases\textsubscript{t}

(8) Can you specify the possible constraints for the variables and
equations? Please add them to the report and implement them in the code.

\textbf{Response:} Thank you. We have added a new subsection,
\emph{Parameter Constraints and Bounds}.

\emph{(9) After you have trained the model, can you also perform
inference and simulate a firm's financial statements for several years?}

\textbf{Response:} We now include the 10-year forecasted Balance Sheet,
Income Statement, and Cash Budget in the \emph{Projected Balance Sheet,
Income Statement, and Cash Budget} subsection.

(10) Remember to write test cases.

Response: Thank you for the reminder. We wrote 87 test cases across 8 test files
covering all modules in the codebase:

\begin{itemize}
\item \texttt{test\_simple\_model\_forecast.py} -- forward simulation model
\item \texttt{test\_trainable\_financial\_model\_simpleopex.py} -- trainable model with deterministic OpEx
\item \texttt{test\_trainable\_financial\_model\_bayesianopex.py} -- trainable model with Bayesian OpEx
\item \texttt{test\_llm\_forecast.py} -- LLM-only balance sheet forecasting
\item \texttt{test\_reasoning\_prompt.py} -- CEO/CFO recommendation prompt
\item \texttt{test\_extraction.py} -- financial statement PDF extraction
\item \texttt{test\_statement\_normalization.py} -- LLM field normalization
\item \texttt{test\_tax\_anomalies.py} -- tax anomaly extraction and model integration
\end{itemize}

\subsection{Feedback on Part 2}

\emph{(1) In the guideline, we have asked you to ``you must submit your
code in TensorFlow and Python.'' TensorFlow is not only for deep
learning and training, but also for vectorization management. That means
you need to rewrite your codes by using TensorFlow in place of NumPy
with the existing framework. NumPy are not strictly forbidden, as they
can still be applied when downloading data from Yahoo Finance, data
preprocessing and plotting.}

\textbf{Response:} We refactored all core modules (training, inference,
model) to use TensorFlow exclusively. All trainable parameters are
\texttt{tf.Variable} instances, the forecast loop and training steps are
compiled into TensorFlow graphs via \texttt{@tf.function}, and gradient
updates use \texttt{GradientTape} for automatic differentiation. NumPy
is retained only for data serialization (\texttt{np.savez}/\texttt{np.load})
and preprocessing. We observed a performance improvement of over 100x
from graph compilation alone.

\emph{(2) Specifically, what is the merit of TensorFlow compared to
NumPy?}

\textbf{Response:} In the context of this codebase, TensorFlow provides
six concrete advantages over NumPy:

\begin{enumerate}
\item \textbf{Automatic differentiation.}
  \texttt{GradientTape} computes gradients for over 50 coupled financial
  parameters without manual calculus. A NumPy implementation would
  require finite-difference approximation or hand-derived gradients for
  every parameter, which is impractical given the interdependencies of the
  financial statements.

\item \textbf{Graph compilation.}
  The \texttt{@tf.function} decorator compiles the forecast loop and
  training step into optimized computation graphs, enabling over 100x
  speedup compared to eager Python/NumPy loops across 25{,}000+ training
  epochs.

\item \textbf{Constrained trainable variables.}
  \texttt{tfp.util.TransformedVariable} with bijectors (Softplus,
  Sigmoid, Chain) automatically enforces economic constraints (such as
  positivity for growth rates, or $[0, 1]$ bounds for payout
  ratios) during gradient updates. NumPy has no equivalent; manual
  clamping after each optimization step would be required.

\item \textbf{Vectorized Monte Carlo simulation.}
  \texttt{tf.TensorArray} and \texttt{tf.while\_loop} inside
  \texttt{@tf.function} efficiently batch 1{,}000+ stochastic
  trajectories. NumPy would require explicit Python loops or complex
  pre-allocated array management that cannot be graph-compiled.

\item \textbf{Variational inference.}
  TensorFlow Probability enables Bayesian OpEx modeling with automatic
  KL divergence gradients through sampling operations. Computing the
  Evidence Lower Bound (ELBO) loss requires differentiating through
  stochastic nodes, which is not possible natively in NumPy.

\item \textbf{Module system.}
  \texttt{tf.Module} provides automatic variable discovery and tracking
  across a hierarchy of policy submodules (CapEx, working capital,
  liquidity, debt, dividends, buyback, tax). This eliminates manual
  parameter bookkeeping that a NumPy implementation would require.
\end{enumerate}

\emph{(3) Codebase structure: Ultimately, we are delivering a Python
package to our client. Please refactor it to an OOP style, with proper
codebase hierarchy.}

\textbf{Response:} We refactored the entire codebase following four
guiding principles:

\begin{enumerate}
\item OOP principles: proper encapsulation, inheritance, polymorphism,
  and abstraction
\item Clean Code principles (Robert Martin): meaningful naming,
  single-responsibility functions and classes, SOLID principles
\item PEP~8 and PEP~20 compliance
\item Pythonic style per Fluent Python: Python idioms over Java-style
  patterns
\end{enumerate}

We used Claude Code's Plan Mode to conduct a systematic code review
before refactoring. The following prompt was used to leverage the tool
effectively:

\begin{lstlisting}[basicstyle=\ttfamily\scriptsize, breaklines=true, captionpos=b, caption={Claude Code planning-mode prompt used for systematic code review before refactoring}, label=lst:code_review_prompt]
Enter planning mode. Do not modify any files yet.

I am a finance/quant intern candidate at JP Morgan and need a thorough,
honest code review of this financial forecasting codebase before I submit it.
The interviewer specifically referenced these three books as the standard
I should meet:

1. Clean Code (Robert C. Martin)
2. Clean Code in Python (Mariano Anaya)
3. Fluent Python (Luciano Ramalho)

Please review the entire codebase systematically across these dimensions:
1. OOP Design & SOLID Principles
2. Module Cohesion & Coupling
3. Pythonic Style (Fluent Python lens)
4. Naming & Readability (Clean Code lens)
5. Type Hints & Contracts
6. Test Coverage Gaps

After completing the review, produce a structured report with:
- A severity rating per issue (Critical / Major / Minor)
- The specific file and line range for each issue
- A one-line explanation of which principle it violates
- Do NOT suggest fixes yet
\end{lstlisting}

\emph{(4) Add docstring and comment: please ensure your code is
well-documented in English and follow established docstring conventions
(such as PEP257 for Python).}

\textbf{Response:} All public interfaces are now documented with
PEP~257-compliant docstrings as part of the refactoring effort.

\emph{(5) Attempting LaTeX: Can you try illustrating your ideas in a
LaTeX-generated PDF?}

\textbf{Response:} The report has been rewritten in LaTeX.

\emph{(6) Double Entry Principle: Can you start with a section to discuss
in detail what double entry principle means by providing with some
examples? An example from items on income statement and one from
statement of cashflow will help the reader understanding why it can keep
the balance sheet balanced.}

\textbf{Response:} We have added a dedicated subsection in the
Introduction (see Section~\ref{sec:double-entry}, \emph{The
Double-Entry Principle}), which explains the principle with two worked
examples: a credit sale (income statement to balance sheet) and a
receivable collection (cash budget to balance sheet).


\section{Cross-Company Forecast Results}\label{app:cross_company}

The following figures present the full model-fit and ten-year Bayesian forecast
results for Microsoft and Pfizer, as discussed in
Section~\ref{sec:cross_company}.

\subsection{Microsoft (MSFT)}

\begin{figure}[H]\centering
\includegraphics[width=\textwidth]{training_results/msft/opex_probabilistic_fit_20260406_070422_monte_carlo.png}
\caption{Bayesian variational inference fit of operating expenses for Microsoft,
showing the posterior mean and 95\% credible band.}
\label{fig:msft_opex_fit}
\end{figure}

\begin{figure}[H]\centering
\includegraphics[width=\textwidth]{training_results/msft/all_elements_forecast_20260406_070419.png}
\caption{Ten-year Bayesian forecast with 95\% Monte Carlo confidence intervals
for Microsoft: historical fit (red) and forward projection (blue) of all
balance-sheet, income-statement, and cash-flow elements.}
\label{fig:msft_forecast}
\end{figure}

\subsection{Pfizer (PFE)}

\begin{figure}[H]\centering
\includegraphics[width=\textwidth]{training_results/pfe/opex_probabilistic_fit_20260406_081331_monte_carlo.png}
\caption{Bayesian variational inference fit of operating expenses for Pfizer,
showing the posterior mean and 95\% credible band.}
\label{fig:pfe_opex_fit}
\end{figure}

\begin{figure}[H]\centering
\includegraphics[width=\textwidth]{training_results/pfe/all_elements_forecast_20260406_081328.png}
\caption{Ten-year Bayesian forecast with 95\% Monte Carlo confidence intervals
for Pfizer: historical fit (red) and forward projection (blue) of all
balance-sheet, income-statement, and cash-flow elements. The COVID-era revenue
spike and subsequent contraction is visible in the income-statement panels.}
\label{fig:pfe_forecast}
\end{figure}

\end{document}
